% -----------------------------------------------------------
\section{1831}
% -----------------------------------------------------------
	\label{EMS-1831}
	
	% ----------------------
	\subsection{Joseph Smith}
	% ----------------------
		\label{EMS-1831-JS}
		
		% -----
		\subsubsection{Introduction to 1831}
		% -----
		
			sdf
			
		% -----
		\subsubsection{Revelation, 2 January 1831 [D\&C 38]}
		% -----
			\label{EMS-1831-JS-2-JAN-1831}
			% \label{EMS-1830-JS-DAY-3LETTERMONTH-YEAR}
		
			\begin{description}
				\item[Print Sources]
			\end{description}

				\begin{tabular}{p{0.5in}p{1in}p{0.5in}p{1in}}
				JSP	 	& D1:229--233	& JSR 		& 98--100 \\
				DC	 	& Secion 38	 	& HC 		& 1:140--1433 \\
				HDDC 	& 490--503		& TCHC 		& 1:105--107 \\
				\end{tabular}
				% HDDC = woodford's DC diss; JSR = Marquardt's JS Revelations; TCHC = Vogel HC
			
			\begin{description}
				\item[Online Access] \href{https://www.josephsmithpapers.org/paper-summary/revelation-2-january-1831-dc-38/1}{Here}
				% \item[Analysis] \hyperref[DC38]{Here}
			\end{description}
			
			\begin{description}
				\item[Background]
			\end{description}
			
				% It might also be useful to give a two sentence context of the period: e.g., "This speech was given in Nauvoo to a small congregation one month prior to the destruction of the Nauvoo Expositor. These and other tensions may have been on the prophet's mind when he gave the speech."
			
			\begin{description}
				\item[Original Source]
			\end{description}
			
				% brief description of original source material, with reference to JSP or somewhere else for more info
				% Background material: it might be helpful to have a note that says whether a given document was presented as a speech, was written in a journal, etc.
				% Also a comment here about how reliable the source might be, or what shape it is in, if there are large omissions or parts that are hard to understand, and so on.
			
			\begin{description}
				\item[Text]
			\end{description}
			
				\begin{tightcenter}
				41\supers{st} Commandment Jan 2\supers{nd} AD 1831
				\end{tightcenter}
				
				Received at farette [Fayette] Seneca County State of New york A Comandment to the Churches in New York at a conference they being Commanded to flee to Ohio \&c
				
				Saying thus saith the Lord God even Jesus Christ the great I am Alph[a] \& Omega the begining \& the end the same which looked upon the wide expance of eternity \& all the Scerifick [seraphic] hosts of Heaven before the world was made the same which k[n]oweth all things for all things are present before mine eyes I am the same which spoke \& the world was made \& all things came by me I am the same which hath taken the Zion of Enoch into mine own bosom \& verily I say even as many as have believed on my name for I am Christ \& in mine own name by the Virtue of the blood which I have spilt have I pled before the Father for them but Behold the residue of the wicked have I kept in Chains of darkness untill the Judgement of the great day which shall come at the end of the Earth \& even so will I cause the wicked that will not hear my voice but harden their hearts \& wo, wo, is their doom But Behold Verily Verily I say unto you that mine eyes are upon you I am in your midst \& ye cannot see me but the day soon cometh that ye shall see me \& know that I am for the \sout{chains} \sout{<vails>} \sout{of} vail\sout{s} of darkness shall soon be rent \& he that is not purified shall not abide the day wherefore gird up your loins \& be prepared Behold the Kingdom is yours \& the enemy shall not overcome Verily I say unto you that ye are clean but not all \& there is none else with whoom I am well pleased for all flesh is corruptabl before me \& the powers of darkness prevail upon the Earth among the Children of men in the presence of all the host <of> Heaven which causeth silence to reign \& all eternity is pained \& the Angels are waiting the great command to Reap down the Earth to gether the tears [tares] that they may be burned \& Behold the enemy is combined \& now I shew unto you a Mystery a thing which is had in seecret Chambers to bring to pass even your distruction in process of time \& ye knew it not but now I tell it <unto> you \& ye are blessed not because of your iniquity neither your hearts of unbelief, for Verily <some of> \sout{ye} you are guilty before me Therefore be ye strong from henceforth fear not for the Kingdom is yours \& for your Salvation I gave it unto you a commandment for I have heard your prayers \& the poor have complained before me \& the rich have I made \& all flesh is mine \& I am no respector to \sout{\& I} persons \& I have made the earth rich \& Behold it is my footstool Wherefore again I will stand upon it \& I hold forth \& deign to give unto you greater Riches even a land of promise a land flowing with milk \& Honey upon which there shall be no curse \& I will give it unto you for the land of your enheritance if you seek it with all your hearts \& this shall be my covenant with you ye shall have it for the land of your inheritence \& for the inheritance of your Children forever while the Earth shall stand \& ye shall Possess it again in eternity no more to pass away But Verily I say unto you that in time ye shall have no King nor Ruler for I will be your King \& watch over you Wherefore hear my voice \& follow me \& ye shall be a free People \& ye shall have no laws but my laws for I am your Law giver \& who can stay my hand But Verily I say unto you teach one another according to the Office wherewith I have appointed you \& let evry man esteem his brother as himself \& practice Virtue \& Holyness before me \& again I say unto you let evry man esteem his Brother as himself for what man among you having twelve sons \& is no respector to them \& they Serve him obediently \& he saith unto the one be thou clothed in Robes \& sit thou here \& to the other be thou clothed in Raggs \& sit thou there \& looketh upon his sons \& saith I am Just Behold I have given unto you a Parable \& it is even as I am I say unto you be one \& if ye are not one ye are not mine \& again I say unto you that the Enemy in the Seecret Chambers seek<[e]th> your lives ye hear of wars in far Countries \& you say in your hearts there will soon be great wars in far Countries but ye know not the hearts of they in your own Land I tell you these things because of your prayers Wherefore treasure up Wisdom in your bosoms lest the wickedness of men \sout{reveal} reveal these things \sout{in your} unto you by their wicke[d]ness in a manner which shall speak in your ears with a voice louder than that which shall shake the Earth but if ye are prepared ye need not fear \& that ye might escape the power of the enemy \& be gethered unto me a Righteous people without spot \& blameless Wherefore for this cause I gave unto you The commandment that ye should go to the Ohio \& there I will give unto you my law \& there you shall be endowed with power from on high \& from thence whomsoever I will shall go forth among all Nations \& it shall be told them what they shall do for I have a great work laid up in store for Israel shall be saved \& I will lead them whithersoever I will \& no power can stay my hand And now I give unto the church in these parts a commandment that certain men among them shall be appointed \& they shall be appointed by the voice of the Church \& they shall look to the poor \& the needy \& administer to their relief that they shall not suffer \& send them forth to the place which I have commanded them \& this shall be their Work to govern the affairs of the \sout{Church} Property of \sout{this} the Church \& they that have farms that cannot be sold let them be left or rented \sout{as} as seemeth them good see that all things are preserved \& when men are endowed with power from on high \& are sent forth all these things shall be gethered unto the Bosom of the Church \& if ye seek the riches which is the will of the Father to give unto you ye shall be the richest of all People for ye shall have the riches of eternity \& it must needs be that the riches of the Earth is mine to give but beware of Pride lest ye become as the Nephites of old \& again I say unto you I give unto you a commandment that evry man both Elder Priest \sout{\&} Teacher \& also Member go to with his might with the Labour of his hands to prepare \& accomplish these things which I have commanded \& let your preaching be the warning voice evry man to his Neighbour in mildness \& in meekness \& go ye out from among the wicked save yourselves be ye clean that bear the vesels of the Lord even so amen
				
		% -----
		\subsubsection{Revelation, 5 January 1831 [D\&C 39]}
		% -----
			\label{EMS-1831-JS-5-JAN-1831}
			% \label{EMS-1830-JS-DAY-3LETTERMONTH-YEAR}
		
			\begin{description}
				\item[Print Sources]
			\end{description}

				\begin{tabular}{p{0.5in}p{1in}p{0.5in}p{1in}}
				JSP	 	& D1:233--236	& JSR 		& 101--102 \\
				DC	 	& Section 39 	& HC 		& 1:143--145 \\
				HDDC 	& 504--515		& TCHC 		& 1:107--108 \\
				\end{tabular}
				% HDDC = woodford's DC diss; JSR = Marquardt's JS Revelations; TCHC = Vogel HC
			
			\begin{description}
				\item[Online Access] \href{https://www.josephsmithpapers.org/paper-summary/revelation-5-january-1831-dc-39/1}{Here}
				% \item[Analysis] \hyperref[DC39]{Here}
			\end{description}
			
			\begin{description}
				\item[Background]
			\end{description}
			
				% It might also be useful to give a two sentence context of the period: e.g., "This speech was given in Nauvoo to a small congregation one month prior to the destruction of the Nauvoo Expositor. These and other tensions may have been on the prophet's mind when he gave the speech."
			
			\begin{description}
				\item[Original Source]
			\end{description}
			
				% brief description of original source material, with reference to JSP or somewhere else for more info
				% Background material: it might be helpful to have a note that says whether a given document was presented as a speech, was written in a journal, etc.
				% Also a comment here about how reliable the source might be, or what shape it is in, if there are large omissions or parts that are hard to understand, and so on.
			
			\begin{description}
				\item[Text]
			\end{description}
				
				\begin{tightcenter}
				42\supers{nd} Commandment Rec\supers{d} J\uline{an}. 5\supersu{th}\supers{.} 1831
				\end{tightcenter}
				
				there was a man by the name of James [Covel] who covenanted with the Lord that he would obey any commandment that the Lord would give through his servent Joseph \& <accordingly> he enquird of the Lord \& he received these words as follows
				
				given at Fayette Seneca County state New York
				
				Saying hearken ye \& listen to the voice of him who is from all eternity to all eternity the great I am even Jesus Christ th[e] light \& the life of the world a light which shineth in darkness \& the darkness comprehendeth it not the same which came in the maridian of time unto my own \& my own Received me not but to as many as received me gave I power to become my Sons \& even so will I give unto as many as Receive me power to become my Sons. \& Verily Verily I say unto you he that receiveth my Gospel Receiveth me \& this is my Gospel Repentance \& Baptism by water \& then Cometh the Baptism of fire \& the Holy ghost yea even the comforter which knoweth all things \& teacheth the peacibl things of the Kingdom \& Now Behold I say unto you my servent James I have looked upon thee \& thy works \& I know thee \& now verily I say unto the[e] thine heart is right before me at present Behold I have bestowed great blessings upon thy head Nevertheless thou hast seen great Sorrow for thou hast rejected me many times because of pride \& be cause of the world but behold the days of thy deliverance is come arise \& be baptized \& wash away your sins calling on my name \& ye shall receive my spirit \& a blessing so great as ye have never known \& I have prepared thee for a greater work thou shalt Preach the fulness of my Gospel which I have sent forth in these last days. yea even the covenant which I have sent forth to recover my People which are of the house of Israel \& it shall come to pass that power shall rest upon thee thou shalt have great faith \& I will be with thee \& go before thy face yea thou art called to Labour in my Vineyard \& to build up my Church \& to bring forth Zion that it may Rejoice upon the hills \& flourish Behold Verily Verily I say unto you thou art not called to go \sout{to} unto the Eastern countries but thou art called to go to Ohio \& in asmuch as my People shall assemble themselves at the Oohio I have kept in store a blessing such as is not known among the children of men \& it shall be poured forth upon \sout{your} their heads \& from thence ye shall go forth into all Nations Behold Verily Verily I say unto you that the people in Ohio call upon me in much faith believeing I would stay my hand in Judgement upon the Nations but I cannot deny my word Wherefore lay to with your might \& call forth Labourers into my Vinyard that it may be pruned for the last time \& inasmuch as they do Repent \& receive the fulness of my Gospel \& become sanctified \& I will stay my hand in Judgement wherefore go forth crying with a loud voice saying the Kingdom of Heaven is at hand crying Hosannah blessed is the name of the most high God go forth Baptizing with water preparing for the time of my coming is at hand the day nor the hour no man knoweth but it shurely shall come \& he that Receiveth these things receiveth me \& they shall be gethered unto me \sout{into} <unto> in time \& in eternity \& again it shall come to pass that on as many as ye shall baptize with water ye shall lay your hands in the name of Christ \& they shall receive the Holy ghost \& shall be \sout{a} looking forth for the time of my coming \& shall know me Behold I come quickly even so amen
				
		% -----
		\subsubsection{Revelation, 6 January 1831 [D\&C 40]}
		% -----
			\label{EMS-1831-JS-6-JAN-1831}
			% \label{EMS-1830-JS-DAY-3LETTERMONTH-YEAR}
		
			\begin{description}
				\item[Print Sources]
			\end{description}

				\begin{tabular}{p{0.5in}p{1in}p{0.5in}p{1in}}
				JSP	 	& D1:236--237	& JSR 		& 102--103 \\
				DC	 	& Section 40 	& HC 		& 1:145 \\
				HDDC 	& 504--515		& TCHC 		& 1:108 \\
				\end{tabular}
				% HDDC = woodford's DC diss; JSR = Marquardt's JS Revelations; TCHC = Vogel HC
			
			\begin{description}
				\item[Online Access] \href{https://www.josephsmithpapers.org/paper-summary/revelation-6-january-1831-dc-40/1}{Here}
				% \item[Analysis] \hyperref[DC40]{Here}
			\end{description}
			
			\begin{description}
				\item[Background]
			\end{description}
			
				% It might also be useful to give a two sentence context of the period: e.g., "This speech was given in Nauvoo to a small congregation one month prior to the destruction of the Nauvoo Expositor. These and other tensions may have been on the prophet's mind when he gave the speech."
			
			\begin{description}
				\item[Original Source]
			\end{description}
			
				% brief description of original source material, with reference to JSP or somewhere else for more info
				% Background material: it might be helpful to have a note that says whether a given document was presented as a speech, was written in a journal, etc.
				% Also a comment here about how reliable the source might be, or what shape it is in, if there are large omissions or parts that are hard to understand, and so on.
			
			\begin{description}
				\item[Text]
			\end{description}
			
				\begin{tightcenter}
				43\supersu{rd}\supers{.} Commandment January 6\supersu{th}\supers{.} 1831
				\end{tightcenter}
				
				A Revelation to Joseph \& Sidney [Rigdon] Rec\supers{d} at Fayette Seneca County state of NY telling them why James [Covel] \sout{Receivd} \sout{not the} obeyed not the Command which he Received \&c
				
				Behold verily I say unto you that his heart was right before me for he covenanted with me that he would obey my word \& he Received the word with Gladness but Straitway Satan came \& tempted him \& the fear of persecution \& the cares of the world caused him to reject the word wherefore he broke the covenant which he had made \& it Remaineth in me to do with him as seemeth me good
				
		% -----
		\subsubsection{Revelation, 4 February 1831 [D\&C 41]}
		% -----
			\label{EMS-1831-JS-4-FEB-1831}
			% \label{EMS-1830-JS-DAY-3LETTERMONTH-YEAR}
		
			\begin{description}
				\item[Print Sources]
			\end{description}

				\begin{tabular}{p{0.5in}p{1in}p{0.5in}p{1in}}
				JSP	 	& D1:241--245	& JSR 		& 105--106 \\
				DC	 	& Section 41 	& HC 		& 1:147 \\
				HDDC 	& 516--524		& TCHC 		& 1:109 \\
				\end{tabular}
				% HDDC = woodford's DC diss; JSR = Marquardt's JS Revelations; TCHC = Vogel HC
			
			\begin{description}
				\item[Online Access] \href{https://www.josephsmithpapers.org/paper-summary/revelation-4-february-1831-dc-41/1}{Here}
				% \item[Analysis] \hyperref[DC41]{Here}
			\end{description}
			
			\begin{description}
				\item[Background]
			\end{description}
			
				% It might also be useful to give a two sentence context of the period: e.g., "This speech was given in Nauvoo to a small congregation one month prior to the destruction of the Nauvoo Expositor. These and other tensions may have been on the prophet's mind when he gave the speech."
			
			\begin{description}
				\item[Original Source]
			\end{description}
			
				% brief description of original source material, with reference to JSP or somewhere else for more info
				% Background material: it might be helpful to have a note that says whether a given document was presented as a speech, was written in a journal, etc.
				% Also a comment here about how reliable the source might be, or what shape it is in, if there are large omissions or parts that are hard to understand, and so on.
			
			\begin{description}
				\item[Text]
			\end{description}
			
				\begin{tightcenter}
				44 Commandment given Fe\supersu{b}\supers{.} 4\supersu{th}\supers{.} 1831
				\end{tightcenter}
				
				at Kirtland Geauga County Ohio given to the Church in these parts it pointing at the office of Edward [Partridge] \&c \& there was a man by the name of Coply [Leman Copley] in the Township of Thompson who had requested <his> Brother [JS] \& Sidney [Rigdon] \sout{\&} to live with him \& he would furnish them houses \& provisions \&c then \sout{By} Joseph enquired of the lord \& Received as follows
				
				Hearken \& hear oh! my People saith your lord \& your God ye whom I delight to bless with the greatest of blessings ye that hear me \& ye that hear me not will I curse with that have professed my name with the heaviest of all cursings hearken oh ye Elders of my Church whom I have called Behold I give unto you a commandment that ye shall assemble yourselves to gether to agree upon my my word \& by the prayer of your faith ye shall receive my law that ye may know how to govern my \sout{Church} Church \& have all things right before me \& I will be your ruler \& ye shall see that my law is kept he that Receiveth my law \& doeth it the same is my Deciple \& he that saith he Receiveth it \& Doeth it not the same is not my Deciple \& shall be cast out from among you for it is not meet that the things which belong to the Children of the Kingdom should be cast before Swine \& again it is meet that my servent Joseph should have a house built in which to live \& translate \& again it is meet that my Servent Sidney should have a comfortable Room to live in \& again I have called my Servent Edward \& give him a commandment that he should be appointed by the voice of the Church \& be ordained a bishop unto the Church \& leave his merchandise \& spend all his time in the labours of the Church \& see to all things as it shall be appointed in my Laws in the day that I shall give them \& this because his heart is pure before me for he is like unto Nathaniel of old in whome there is no guile these words are given unto you \& they are pure before me wherefore be ye aware how you hold them for they are to be answered upon your souls in the day of Judgement even so amen
				
		% -----
		\subsubsection{Revelation, 9 February 1831 [D\&C 42:1--72]}
		% -----
			\label{EMS-1831-JS-9-FEB-1831}
			% \label{EMS-1830-JS-DAY-3LETTERMONTH-YEAR}
		
			\begin{description}
				\item[Print Sources]
			\end{description}

				\begin{tabular}{p{0.5in}p{1in}p{0.5in}p{1in}}
				JSP	 	& D1:245--256	& JSR 		& 107--115 \\
				DC	 	& Section 42 	& HC 		& 1:148--154 \\
				HDDC 	& 525--569		& TCHC 		& 1:110--113 \\
				\end{tabular}
				% HDDC = woodford's DC diss; JSR = Marquardt's JS Revelations; TCHC = Vogel HC
			
			\begin{description}
				\item[Online Access] \href{https://www.josephsmithpapers.org/paper-summary/revelation-9-february-1831-dc-421-72/1}{Here}
				% \item[Analysis] \hyperref[DC42]{Here}
			\end{description}
			
			\begin{description}
				\item[Background]
			\end{description}
			
				% It might also be useful to give a two sentence context of the period: e.g., "This speech was given in Nauvoo to a small congregation one month prior to the destruction of the Nauvoo Expositor. These and other tensions may have been on the prophet's mind when he gave the speech."
			
			\begin{description}
				\item[Original Source]
			\end{description}
			
				% brief description of original source material, with reference to JSP or somewhere else for more info
				% Background material: it might be helpful to have a note that says whether a given document was presented as a speech, was written in a journal, etc.
				% Also a comment here about how reliable the source might be, or what shape it is in, if there are large omissions or parts that are hard to understand, and so on.
			
			\begin{description}
				\item[Text]
			\end{description}
			
				\begin{tightcenter}
				The Laws of the Church of Christ
				\end{tightcenter}
				
				Kirtland Geauga Ohio May 23\supers{d} 1831 A Commandment to <the> Elder[s]
				
				[First Shall the Church come to gether into one place or continue in seperate establishments?]
				
				Hearken oh! ye Elders of my Church who have assembled yourselves together in my name even Jesus Christ the Son of \sout{God} the living God the Saveiour of the world in as much as they believe on my name \& keep my commandments again I say unto you hearken \& hear \& obey the laws which I Shall give unto you for verily I say as ye have assembled yourselves together according to the commandment wherewith I commanded you and are agreed as touching this one thing \& have asked me in my name even so ye shall receive Behold verily I say unto you I give unto you this first Commandment that ye shall go forth in my name every one of you except my servant Joseph \& Sidny [Sidney Rigdon] \& I give unto them a commandment that they shall go forth for a little Season \& it shall be given by the power of my spirit when they Shall return \& ye Shall go forth in the power of my Spirit preaching my Gospel two by two in my name lifting up your voices as with the voice of a trump declaring my word like unto Angels of God \& ye shall go forth baptizing with water <saying> repent ye repent ye for the kingdom of Heaven is at hand \& from this place ye shall go forth in to the regions westward \& in as much as ye shall find my deciples ye shall build up my church in every region untill the time shall come when it Shall be revealed unto you from on high \& the City of the New Jerusalam Shall be prepared that ye may be gathered in one that ye may be my people \& I will be your God \& again I say unto you that my Servant Edward [Partridge] shall stand in the office wherewith I have appointed him \& it shall come to pass that if he transgress another shall be appointed in his Stead even so Amen------

				The Law? Again I say unto you that it shall not be given unto any one to go forth to preach my gospel or to build up my church except they be ordained by some one that hath authority \& it is known to the church that he hath authority \& have been regularly ordained by the leaders of the church \& again the Elders priests \& teachers of this Church shall teach the scriptures which are in the Bible \& the Book of Mormon in the which is the fullness of the Gospel \& thou Shalt observe the covenants \& church articles to do them \& <this> Shall be thy teaching<s> \& thou shalt be directed by the spirit it shall be given the[e] by the prayer of faith \& if ye receive not the spirit ye shall not teach \& all this ye shall observe to do as I have commanded concerning your teaching untill the fulness of my Scripture be given \& as ye Shall lift your voices by the comforter ye shall speak and prophecy as seemeth me good for behold the Comforter knoweth all things \& beareth record of the father \& the Son \& now behold I Speak unto the church------
				
				Thou shalt not kill \& he that killeth shall not have forgiveness neither in this world neither in the world to come \& again thou shallt not kill he that killeth shall die, Thou shalt not steel \& he that steeleth \& will not repent shall be cast out Thou shalt not lie he that lieth \& will not repent shall be cast out Thoue shalt love thy wife with all thy heart \& shalt cleave unto her \& none else \& he that looketh on a woman to lust after her shall deny the faith \& shall not have the spirit \& if he repent not he shall be cast out------
				
				Thou shalt not commit adultery \& he that committeth adultery \& repenteth not shall be cast out \& he that committeth adultery \& repentteth with all his heart \& forsaketh \& doeth it no more thou shalt forgive him but if he doeth it again he shall not be forgiven but shall be cast out Thou shalt not speak evil of thy neighbour or do him any harm Thou knowest my laws they <are> given in my Scriptures he that sinneth \& repenteth not shall be cast out if thou lovest me thou shall serve \& keep all my commandments \& Behold thou shalt conscrate all thy \sout{property} properties that which thou hast unto me with a covena[n]t and Deed which cannot be broken \& they Shall be laid before the Bishop of my church \& two of the Elders such as he shall appoint \& set apart for that purpose \& it shall come to pass that the Bishop of my church after that he has received the properties of my church that it cannot be taken from \sout{him} you he shall appoint every man a Steward over his own property or that which he hath received in as much as shall be sufficient for him self and family \& the residue shall be kept to administer to him that hath not that every man may receive according as he stands in need \& the residue shall be kept in my store house to administer to the poor and needy as shall be appointed by the Elders of the church \& the Bishop \& for the purpose of purchaseing Land \& building up of the New Jerusalem which is here after to be revealed that my covenant people may be gathered in one in the day that <I> shall come to my temple \& this I do for the salvation of my people \& it shall come to pass that he that sinneth \& rep[e]nteth not shall be cast out \& shall not receive again that which he hath consecrated unto me, for it shall come to pass that which I spake by the mouth of my prophets shall be fulfilld for I will consecrate the riches of the Gentiles unto my people which are of the house of Israel \& Again thou Shalt not be proud in thy heart let all thy garments be plain \& their beauty the beauty of the work of thine own hands \& let all things be done in cleanliness before me------
				
				Thou Shalt not be Idle for he that is Idle shall <not> eat the bread nor wear the garment of the labourer \& whosoever among you that is sick \& hath not faith to be healed but believeth Shall be nourished in all tenderness with herbs and mild food \& that not of the world \& the Elders of the church two or more Shall be called \& shall pray for and lay their hands upon them in my name \& if they die they shall die unto me \& if they shall live <they shall live> unto me Thou \sout{shall} <shalt> live together in love insomuch that thou shalt weep for the loss of them that die \& more especially for those that have not hope of a glorious resurrection \& it shall come to pass that \sout{that} they that die in me shall not taste of death for it shall be sweet unto them \& they that die not in me wo is them for their death is bitter \& again it shall come to pass that he that hath faith in me to be healed \& is not appointed unto death shall be healed he that hath faith to see shall see he that hath faith to hear shall hear the Lame that have faith <to> leep shall <leep> \& they that have not faith to do these things but believe in me hath power to become my sons \& in asmuch as they break not my Laws thou Shalt bear their infirmities thou shalt stand in the place of thy Stewardship thou shalt not take thy brothers garment thou shalt pay for that which thou \sout{<shalt receive>} shall receive of thy Brother \& if thou obtain more than that which would be for thy support thou shalt give it unto my store house that it may be done according to that which I have Spoken, Thou shalt ask \& my scriptures shall be given as I have appointed \& for thy Salvation thou shalt hold thy peace concerning them untill ye have received them \& then I give unto you a commandment that ye should teach them unto all men \& they also shall be taught unto all nations kindreds Toungs \& People thou Shalt take the things which thou hast received which thou knowest to have been my Law to be my Law to govern my church \& he that doeth according to these things shall be saved \& he that doeth them not shall be damned if he continue[.] if thou shalt ask thou shalt receive Revelation upon Revelation knowledge upon knowledge that thou mayest know the mysteries \& the peacible things of the kingdom that which bringeth joy that which life Eternal thou shalt ask \& it shall be revealed unto you in mine own due time when the New Jerusalem shall be built thou shalt ask \& it shall be revealed in mine own due time \& behold it shall come to pass that my Servants Shall be sent both to the East \& to the west the north \& to the South \& even now let him that goes to the East teach them that are converted to flee to the west \& this because [of] that which is to come \& secret combinations Behold thou shalt observe all these things \& great Shall be thy reward thou shalt observe to keep the mysteries of the Kingdom unto thy Self for it is not given unto the world to know the mysteries \& these Laws which ye have received are sufficient for \sout{<you>} both here \& in the New Jerusalem but he that lacketh knowledge let him ask of me \& I will give him liberally \& \sout{\&} upbraid him not Lift up your hearts \& rejoice for unto you the Kingdom is given
				
				\begin{tightright}
				Even so Amen------
				\end{tightright}
				
				3\supers{d} How the Elders are to dispose of their families while they are proclaiming repentance or are otherwise engaged in the Service of the Church?
				
				The Priests and Teachers Shall have their Stewardship given them as the members. And the Elders are to assist the Bishop in all things \& he is to see that their families are supported out of the property which is consecrated to the Lord either a stewardship or otherwise as may be thought best by the Elders \sout{\&} \& Bishop
				
				4\supers{th} How far it is the will of the Lord that we Should have dealings with the wo[r]ld \& how we Should conduct our dealings with them?
				
				Thou Shalt contract no debts with them \& again the Elders \& Bishop shall Council together \& they shall do by the direction\sout{s} of the spirit as it must be necessary------
				
				5\supers{th} What preperations we shall make for our Brethren from the East \& when \& how?
				
				There shall be as many appointed as must needs be necessary to assist the Bishop in obtaining places that they may be together as much as can be \& is directed by the holy Spirit \& every family Shall have places that they may live by themselves \& every Church Shall be organized in as close bodies as they can be in consequence of the enemy!------
				
		% -----
		\subsubsection{Revelation, February 1831--A [D\&C 43]}
		% -----
			\label{EMS-1831-JS-FEB-1831-A}
			% \label{EMS-1830-JS-DAY-3LETTERMONTH-YEAR}
		
			\begin{description}
				\item[Print Sources]
			\end{description}

				\begin{tabular}{p{0.5in}p{1in}p{0.5in}p{1in}}
				JSP	 	& D1:256--260	& JSR 		& 115--118 \\
				DC	 	& Section 43 	& HC 		& 1:154--156 \\
				HDDC 	& 570--583		& TCHC 		& 1:114--115 \\
				\end{tabular}
				% HDDC = woodford's DC diss; JSR = Marquardt's JS Revelations; TCHC = Vogel HC
			
			\begin{description}
				\item[Online Access] \href{https://www.josephsmithpapers.org/paper-summary/revelation-february-1831-a-dc-43/1}{Here}
				% \item[Analysis] \hyperref[DC43]{Here}
			\end{description}
			
			\begin{description}
				\item[Background]
			\end{description}
			
				% It might also be useful to give a two sentence context of the period: e.g., "This speech was given in Nauvoo to a small congregation one month prior to the destruction of the Nauvoo Expositor. These and other tensions may have been on the prophet's mind when he gave the speech."
			
			\begin{description}
				\item[Original Source]
			\end{description}
			
				% brief description of original source material, with reference to JSP or somewhere else for more info
				% Background material: it might be helpful to have a note that says whether a given document was presented as a speech, was written in a journal, etc.
				% Also a comment here about how reliable the source might be, or what shape it is in, if there are large omissions or parts that are hard to understand, and so on.
			
			\begin{description}
				\item[Text]
			\end{description}
			
				\begin{tightcenter}
				45\supersu{th}\supers{.} Commandment AD 1831
				\end{tightcenter}
				
				given to the Elders of this Church at Kirtland Geauga Ohio
				
				Oh hearken ye Elders of my Church \& give ere [ear] to the words which I shall speak unto you for Behold verily Verily I say unto you that ye have received a commandment for a law unto my Church through him whom I have appointed unto you to receive commandments \& Revelations from my hand \& this ye shall know asshuredly that there is none other appointed unto you to receive commandments \& Revelations untill he be taken if he abide in me. but Verily Verily I say unto you that none else shall be appointed unto this gift except it be through him for if it be taken from him he shall not have power except to appoint another in his stead \& this shall be a law unto you that ye receive not the teachings of any that shall Come before you \sout{with} as Revelations or commandments \& this I give unto you that you may not be deceived, that you may know they are not of me for Verily I say unto you that he that is ordained of me shall come in at the gate \& be ordained as I have told you before to teach those Revelations which you have received \& shall receive through him whom I have appointed--- \& now Behold I give unto you a commandment that when ye are assembled yourselves together ye shall note with a Pen how to act, \& for my Church to act upon the points of my law \& commandments which I have given \& thus it shall become a law unto you being Sanctified by that which ye have received that ye shall bind yourselves to act in all holiness before me that in as much as ye do this glory shall be ad[d]ed to the Kingdom which ye have received, inasmuch as ye do it not it shall be taken even that which ye have received, purge ye out the iniquity which is among you Sanctify yourselves before me \& if ye desire the glories of the Kingdom appoint ye my Servent \& uphold him before me by the prayer of faith \& again I say unto you that if you desire the mysteries of the Kingdom provide for him food \& raiment \& whatsoever \sout{is} thing he needeth to accomplish the work which I have commanded him \& if ye do it not he shall remain unto them that have received him that I may reserve unto myself a pure People before me---
				
				Again I say hearken ye Elders of my Church whom I have appointed ye are not sent forth to be taught but to teach the Children of men the things which I have put in your hands by the power of my Spirit \& ye are to be taught from on high Sanctify yourselves \& ye shall be endowed with power \sout{from on high} that ye may give even as I have spoken hearken ye for Behold the great day of the Lord is nigh at hand for the day cometh that the Lord shall utter his voice out of Heaven the Hevens shall shake \& the Earth shall tremble \& the Trump of God shall sound both long \& loud \& shall say to the sleeping Nations ye saints arise \& live ye sinners stay \& sleep untill I shall call again Wherefore gird up your loins lest ye are found among the wicked lift up your voices \& spare not call upon the Nations to repent both old \& young both bond \& free saying prepare yourselves for the great day of the Lord for if I who am a man do lift up my voice \& call upon you to repent \& ye hate me\sout{n} what will you say when the day cometh when the Thunders shall utter their voices from the ends of the Earth speaking in the ears of all that live saying repent \& prepare for the great day of the Lord yea \& again when the lightning shall streak forth from the East unto the west \& shall utter forth their voices unto all that live \& make the ears of all tingle that hear saying these words repent ye for the great day of the Lord is come \& again the lord shall utter his voice out of Heaven saying hearken O ye Nations of the Earth \& hear the words of that God who made you O ye Nations of the Earth how often would I have gethered you as a hen gethereth her chickens under her wings but ye would not how oft have I called upon you by the mouth of my Servents \& by the ministering of Angels \& by the voice of lightnings \& by the voice of tempests \& by the voice of Earthquakes \& great hailstorms \& by the voice of famines \& pestilences of evry kind \& by the great sound of a trump \& by the voice of Judgements \& by the voice of mercy all the day long \& by the voice of Glory \& honour \& the riches of eternal life \& would have saved you with an everlasting salvation but ye would not Behold the day has come when the cup of the wrath of mine indignation is full Behold verily I say unto you that these are the words of the Lord your God Wherefore Labour ye Labour ye in my vineyard for the last time for the last time call ye upon the inhabitants of the Earth for in mine Own due time will I come upon the Earth in Judgement \& my People shall be redeemed \& shall reign with me on Earth for the great Millenial which I have spoken by the mouth\sout{s} of my Servents shall come for satan shall be bound \& when he is loosed again he shall only reign for a little Season \& then cometh the end of the world \& he that liveth in righteousness shall be changed in the twinkling of an eye \& the Earth shall pass away so as by fire \& the wicked shall go away into unquinchable fire \& their end no man knoweth on Earth nor ever shall know untill they come before me in Judgement hearken ye to these words Behold I am Jesus Christ the saveiour of the World treasure these things up in your hearts \& let the Solemn[i]ties of Eternity rest upon your minds be sober keep all the commandments even So amen
				
		% -----
		\subsubsection{Revelation, February 1831--B [D\&C 44]}
		% -----
			\label{EMS-1831-JS-FEB-1831-B}
			% \label{EMS-1830-JS-DAY-3LETTERMONTH-YEAR}
		
			\begin{description}
				\item[Print Sources]
			\end{description}

				\begin{tabular}{p{0.5in}p{1in}p{0.5in}p{1in}}
				JSP	 	& D1:260--261	& JSR 		& 118 \\
				DC	 	& Section 44 	& HC 		& 1:157 \\
				HDDC 	& 584--589		& TCHC 		& 1:115--116 \\
				\end{tabular}
				% HDDC = woodford's DC diss; JSR = Marquardt's JS Revelations; TCHC = Vogel HC
			
			\begin{description}
				\item[Online Access] \href{https://www.josephsmithpapers.org/paper-summary/revelation-february-1831-b-dc-44/1}{Here}
				% \item[Analysis] \hyperref[DC44]{Here}
			\end{description}
			
			\begin{description}
				\item[Background]
			\end{description}
			
				% It might also be useful to give a two sentence context of the period: e.g., "This speech was given in Nauvoo to a small congregation one month prior to the destruction of the Nauvoo Expositor. These and other tensions may have been on the prophet's mind when he gave the speech."
			
			\begin{description}
				\item[Original Source]
			\end{description}
			
				% brief description of original source material, with reference to JSP or somewhere else for more info
				% Background material: it might be helpful to have a note that says whether a given document was presented as a speech, was written in a journal, etc.
				% Also a comment here about how reliable the source might be, or what shape it is in, if there are large omissions or parts that are hard to understand, and so on.
			
			\begin{description}
				\item[Text]
			\end{description}
				
				\begin{tightcenter}
				46\supers{th} Commandment Feb\uline{u}. 1831
				\end{tightcenter}
				
				A Revelation to Joseph \& Sidney [Rigdon] Receivd at Kirtland Geauga Ohio a call to the Eldrs of this \sout{Church} \sout{\&c} Church \&c
				
				Blehold thus saith the Lord unto you my Servents it is expedient in me that the Elders of my Church should be called to gether from the East \& from the West \& from the North \& from the South by letter or some other way \& it shall come to pass that I will pour out my Spirit upon them in the day that they assemble themselves together \& it shall come to pass that they shall go forth unto the regions round about \& preach repentance unto this People \& many shall be converted insomuch that ye shall obtain power to organize yourselves according to the laws of man that your enemies may be under your feet in all things that ye may be enabled to keep my laws that evry band may be broken wherewith the enemy seeketh to destroy my People Behold I say unto you that ye must visit the poor \& the needy \& administer to their releaf that they may be kept untill all things may be done according to my law which ye have receivd amen
				
		% -----
		\subsubsection{Revelation, 23 February 1831 [D\&C 42:74--93]}
		% -----
			\label{EMS-1831-JS-23-FEB-1831}
			% \label{EMS-1830-JS-DAY-3LETTERMONTH-YEAR}
		
			\begin{description}
				\item[Print Sources]
			\end{description}

				\begin{tabular}{p{0.5in}p{1in}p{0.5in}p{1in}}
				JSP	 	& D1:264--267	& JSR 		& 118--120\footnote{Note how Marquardt breaks these up into two sections.} \\
				DC	 	& Section 42 	& HC 		& 1:148--154 \\
				HDDC 	& 525--569		& TCHC 		& 1:110--113 \\
				\end{tabular}
				% HDDC = woodford's DC diss; JSR = Marquardt's JS Revelations; TCHC = Vogel HC
				
			\begin{description}
				\item[Online Access] \href{https://www.josephsmithpapers.org/paper-summary/revelation-23-february-1831-dc-4274-93/1}{Here}
				% \item[Analysis] \hyperref[DC42]{Here}
			\end{description}
			
			\begin{description}
				\item[Background]
			\end{description}
			
				% It might also be useful to give a two sentence context of the period: e.g., "This speech was given in Nauvoo to a small congregation one month prior to the destruction of the Nauvoo Expositor. These and other tensions may have been on the prophet's mind when he gave the speech."
			
			\begin{description}
				\item[Original Source]
			\end{description}
			
				% brief description of original source material, with reference to JSP or somewhere else for more info
				% Background material: it might be helpful to have a note that says whether a given document was presented as a speech, was written in a journal, etc.
				% Also a comment here about how reliable the source might be, or what shape it is in, if there are large omissions or parts that are hard to understand, and so on.
			
			\begin{description}
				\item[Text]
			\end{description}
				
				\begin{tightcenter}
				February 23\supers{d} 1831 the rules and regulations of the <Law>
				\end{tightcenter}
				
				How the Elders of the church of Christ are to act upon the points of the Law given by Jesus Christ to the Church in the presents of twelve Elders February 9\supers{th} 1831 as agreed upon by Seven \sout{Elders} Elders Feby 23\supers{d} 1831 according to to the commandment of God------
				
				1\supers{th} The first commandment in the law teaches that all the Elders shall go unto the regions westward and labour to build up Churches unto Christ wheresoever they shall <find> any to receive them and obey the Gospel of Jesus Christ except Joseph \& Sidney [Rigdon] and Edward [Partridge] and Such as the Bishop Shall appoint to assist him in his duties according to the Law which we have received this commandment as far as it respects these Elders to be sent to the west is a special one for the time being incumbent on the present Elders who shall return when directed by the Holy Spirit------
				
				2\supers{d} Every person who belongeth to this church of Christ Shall observe all the Commandments and covenants of the Church and it shall come to pass that if any person among you shall kill they shall be delivered up and dealt with according to the Laws of the land for remember that he hath no forgiveness and it shall be proven according to the Laws of the land but if any man shall commit Adultery he Shall be tried before two Elders of the Church or more and every word shall be established against him by two witnesses of the Church and not of the world but if there are more than two witnesses it is better but he shall be condemned by the mouth of two witnesses and the Elders shall lay the case before the Church and the Church shall lift up their hands against them that they may be \sout{dealt} dealt with according to <the> Law and if it can be it is necessary that the Bishop is present also and thus ye shall do in all cases which shall come before you and if a\sout{n} man shall rob he shall be delivered up unto the Law and if he shall steal he shall be delivered up unto the Law and if he lie he shall be delivered up unto the Law if he do any manner of iniquity he shall be delivered up unto the Law even that of God and if thy Brother offend thee, thou shalt take him between him and thee alone and if he confess thou shalt be reconciled and if he confess not thou shalt take another with thee and then if he confess not thou Shalt deliver him up unto the Church not to the members but to the Elders and it shall be done in a meeting and that not before the world and if thy Brother offend many he shall be chastened before many and if any one offend openly he shall be rebuked openly that he may be ashamed and if he confess not he Shall be delivered up unto the law, if any shall offend in secret he shall be rebuked in Secret \sout{the} <that he> may have oportunity to confess in Secret to him whome \sout{he} he has offended and to God that the Brethren may not speak reproachfully of him and thus shall ye conduct in all things------
				
				\begin{tightcenter}
				How to act in cases of adultery
				\end{tightcenter}
				
				Behold verily I say unto you whatsoever person among you having put away their companion for the cause of fornication or in otherwords if he shall testify before you in all Lowliness of heart that this is the case ye shall not cast them out from among you but if ye shall find that any person hath left their companion for the sake of adultery and they themselves are the offender and their companions are living they shall be cast out \sout{out} from among you and again I say unto you that ye be watchful and careful with all inquiry that ye receive none such among you if they are married and if they are not married they shall repent of all their sins or ye shall not receive them------
				
		% -----
		\subsubsection{Letter to Hyrum Smith, 3--4 March 1831}
		% -----
			\label{EMS-1831-JS-3-4-MAR-1831}
			% \label{EMS-1830-JS-DAY-3LETTERMONTH-YEAR}
		
			\begin{description}
				\item[Print Sources]
			\end{description}

				\begin{tabular}{p{0.5in}p{1in}p{0.5in}p{1in}}
				JSP	 	& D1:267--273	& JSR 		& --- \\
				DC	 	& ---		 	& HC 		& --- \\
				HDDC 	& ---			& TCHC 		& --- \\
				\end{tabular}
				% HDDC = woodford's DC diss; JSR = Marquardt's JS Revelations; TCHC = Vogel HC
			
			\begin{description}
				\item[Online Access] \href{https://www.josephsmithpapers.org/paper-summary/letter-to-hyrum-smith-3-4-march-1831/1}{Here}
				% \item[Analysis] \hyperref[XX]{Here}
			\end{description}
			
			\begin{description}
				\item[Background]
			\end{description}
			
				% It might also be useful to give a two sentence context of the period: e.g., "This speech was given in Nauvoo to a small congregation one month prior to the destruction of the Nauvoo Expositor. These and other tensions may have been on the prophet's mind when he gave the speech."
			
			\begin{description}
				\item[Original Source]
			\end{description}
			
				% brief description of original source material, with reference to JSP or somewhere else for more info
				% Background material: it might be helpful to have a note that says whether a given document was presented as a speech, was written in a journal, etc.
				% Also a comment here about how reliable the source might be, or what shape it is in, if there are large omissions or parts that are hard to understand, and so on.
			
			\begin{description}
				\item[Text]
			\end{description}
			
					\hypertarget{EMS-1831-JS-3-4-MAR-1831-a}%
				Brother
					\marginnote{a}
				Hyram [Hyrum Smith]
				
				we arived here safe and are all well I hav[e] been ingageed in regulating the Churches here as the deciples are numerous and the devil had made many attempts to over throw them it has been a Serious job but the Lord is with us and we have overcome and have all things regular the work is brakeing forth on the <right> hand and on the left and there is a great Call for Elders in this place...
				
					\hypertarget{EMS-1831-JS-3-4-MAR-1831-b}%
				I
					\marginnote{b}%
				<have> had much Concirn about you but I always remember you in \sout{your} <my> prayers Calling upon god to keep <you> Safe in spite <of> men or devils I think <you> had better Come into this Country immediately for the Lord has Commanded us that we Should Call the Elders of \sout{the} this Chursh to gether unto this plase as soon as possable
				
					\hypertarget{EMS-1831-JS-3-4-MAR-1831-c}%
				March
					\marginnote{c}%
				forth this morning after being Colled out of my bed in the night to go a small distance I went and had \sout{and} an awful strugle with satan <but> being armed with the power of God he was cast out and the woman is Clothed in hir right mind the Lord worketh wonders in this land
				
		% -----
		\subsubsection{Revelation, circa 7 March 1831 [D\&C 45]}
		% -----
			\label{EMS-1831-JS-CIR-7-MAR-1831}
			% \label{EMS-1830-JS-DAY-3LETTERMONTH-YEAR}
		
			\begin{description}
				\item[Print Sources]
			\end{description}

				\begin{tabular}{p{0.5in}p{1in}p{0.5in}p{1in}}
				JSP	 	& D1:274--280	& JSR 		& 121--124 \\
				DC	 	& Section 45 	& HC 		& 1:159--163 \\
				HDDC 	& 590--609		& TCHC 		& 1:117--120 \\
				\end{tabular}
				% HDDC = woodford's DC diss; JSR = Marquardt's JS Revelations; TCHC = Vogel HC
			
			\begin{description}
				\item[Online Access] \href{https://www.josephsmithpapers.org/paper-summary/revelation-circa-7-march-1831-dc-45/1}{Here}
				% \item[Analysis] \hyperref[DC45]{Here}
			\end{description}
			
			\begin{description}
				\item[Background]
			\end{description}
			
				% It might also be useful to give a two sentence context of the period: e.g., "This speech was given in Nauvoo to a small congregation one month prior to the destruction of the Nauvoo Expositor. These and other tensions may have been on the prophet's mind when he gave the speech."
			
			\begin{description}
				\item[Original Source]
			\end{description}
			
				% brief description of original source material, with reference to JSP or somewhere else for more info
				% Background material: it might be helpful to have a note that says whether a given document was presented as a speech, was written in a journal, etc.
				% Also a comment here about how reliable the source might be, or what shape it is in, if there are large omissions or parts that are hard to understand, and so on.
			
			\begin{description}
				\item[Text]
			\end{description}
			
				\begin{tightcenter}
				47 A prophecy March 7\supersu{th}\supers{.} 1831
				\end{tightcenter}
				
				given to Joseph the seer at Kirtland geauga County Ohio
				
				Saying hearken O ye people of my Church to whom the Kingdom \sout{was} has been given hearken ye \& give ere [ear] to him who laid the foundation of the Earth who made the Heavens \& all the hosts thereof \& by whom all things were made which live \& move \& have a being \& again I say hearken unto my voice lest death shall overtake you in an hour when ye think not the Summer shall be past \& the harvest ended \& your souls not saved listen to him who is the advocate with the Father who is pleading your case before him saying Father behold the sufferings \& death of him who did no sin in whom thou wast well pleased Behold the Blood of thy son which was shed the blood of him whom thou gavest that thyself might be glorified wherefore Father spare these my Brethren that Believe on my name that they may come unto me And have everlasting life Hearken O ye people of my Church \& ye Elders listen together \& hear my voice whilest it is called to day \& harden not your hearts For verily I say unto you that I am Alpha \& Omega the Begining \& the end the light \& the life of the world a light that shineth in darkness \& the darkness comprehendeth it not I came unto my own \& my own Received me not but unto as many as received me gave I power to do many Miricles \& to become the Sons of God \& even unto them that believed on my name gave I power to obtain eternal life \& even so I have sent mine everlasting covenant unto the World to be a light to the world \& to be a standerd for my people \& for the gentiles to seek to it And to be a mesinger before my face to prepare the way before me
				
				Wherefore come ye unto it With him that cometh I will reason as with men \sout{of old} in days of old And I will shew unto you my strong reasoning Wherefore hearken ye together \& let me shew it unto you even my wisdom the wisdom of him whom ye say is the God of Enoch \& his Brethren who were seperated from the Earth \& were reserved unto myself a City reserved untill a day of righteousness shall come a day which was sought for by all Holy men \& they found it not Because of wickedness \& abominations \& confessed that they were strangers \& pilgrims on the Earth but obtained a promise that they should find it \& see it in their flesh wherefore hearken \& I will reason with you \& I will speak unto you \& prophecy as unto men in days of old \& I will shew it plainly as I shewed it unto my Deciples as I stood before them in the flesh \& spake unto them saying <as> ye have asked of me concerning these signs of my coming in the day when I shall come in my glory in the clouds of Heaven to fulfill the promises that I have made unto your fathers for as you have looked upon the long absence of your spirits from your bodies to be a bondage I will shew unto you how the day of redemption shall come \& also the restoration of the scattered Israel \& now ye behold this temple which is in Jerusalem which ye call the House of God \& your enemies say that this House shall never fall but verily I say unto you that desolation shall come upon this generation as a thief in the night And this people shall be destroyed \& scattered among all Nations \& this Temple which ye now see shall be thrown down that there shall not be left an stone upon another \& it shall come to pass that this generation of Jews shall not pass away [until every] desolation which I have told you concerning them shall come to pass ye say that ye know that the end of the World cometh ye say also that ye know that the Heavens \& the Earth shall pass away And in this ye say truly for so it is But these things which I have told you shall not pass away but all shall be fulfilled \& this I have told you concerning Jerusalum \& when that day shall come shall a remnant \sout{shall a} be scattered among all Nations but they shall be gethered again but they shall remain untill the times of the gentiles be fulfelled \& in that day shall be heard of wars \& rumours of wars \& the whole Earth shall be in commotion \& mens hearts shall fail them \& \sout{\&} shall say that Christ delayeth his coming until the end of the world \& the love of men shall wax cold \& inequity shall abound \& when the times of the gentiles \sout{shall be} is come in \sout{And} a light shall break forth among them that sit in darkness \& it shall be the fulness of my Gospel but they receive it not for they perceive not the light \& they turn their hearts from me because of the precepts of men \& in that generation shall the times of the gentiles be fulfilled \& there shall be men standing in that generation that shall not pass untill they shall see an overflowing scourge for a desolating sicknes shall cover the land \sout{\& shall not be} \sout{moved} but my Deciples shall stand in Holy places \& shall not be moved but among the wicked men shall lift up their voices \& curse God \& die \& there shall be earthqakes also in diverse places \& many desolations yet men will harden their hearts against me \& they will take up the sword one against another \& they will kill one another And now when I the Lord had spoken these words unto my Deciples they were troubled \sout{for when all these} \sout{things shall come} \& I said unto them be not troubled for when all these things shall come to pass ye may know that the promises which have been made unto you shall be fulfilled \& when the light shall begin to break forth it shall be with them like unto a Parable which I will shew you ye look \& behold the figgtrees \& ye see them with your eyes \& ye say when they begin to shoot forth \& th$\diamond\diamond\diamond$ [their] leaves are yet tender ye say that summer is now nigh at hand even so it shall be in that day when they shall see all these things then shall they know that the hour is nigh \& it shall come to pass that he that feareth me shall be looking for the great day of the lord to come even for the signs of the coming of the son of man \& they shall see signs \& wonders for they shall be shewn forth in the heavens above \& in the Earth beneath \& they shall behold blood \& fires \& vapors of smoke \& before the day of the lord come the sun shall be darkened \& the moon be turned into blood \& some stars shall fall from Heaven \& the remnant shall be gethered unto this place \& then they shall look for me \& Behold I will come \& they shall see me in the clouds of heaven clothed with power \& great glory with all the holy Angels \& he that watches not for me shall be cut off But before the arm of the Lord shall fall \& the Angel shall sound his Trump \& the saints that have slept shall come forth to meet me in the cloud wherefore if ye have slept in peace blessed are you for as you now Behold me \& know that I am even so shall ye come unto me \& your souls shall live And your redemption shall be perfected And the saints shall come forth from the four quarters of the Earth then shall the arm of the Lord fall upon the Nations \& then shall the lord set his foot upon this mount \& it shall cleave in twain \& the Earth shall tremble \& reel to \& fro \& the Heavens also shall shake \& the Lord shall utter his voice \& all the ends of the Earth shall hear it \& the Nations of the Earth shall mourn \& they that have laughed shall see their folly \& calamity shall cease the mocker \& the scorner shall be consumed \& they that have watched for iniquity shall be \sout{cut off} hewn down \& cast into the fire \& then shall the Jews look upon me \& say what are these wounds in thine hands \& in thy feet then shall they know that I am the Lord for I will say unto them these wounds are the wounds With which I was wounded in the house of my friends I am he that was lifted up I am Jesus which was crusified I am the Son of God \& then shall they weep because of their iniquities then shall they lament because they persecuted their King \& then shall the heathen Nations be redeemed \& they which knew no law shall have part in the first resurrection \& it shall be tolerable for them \& Satan shall be bound that he shall have no place in the hearts of the children of men \& at that day when I shall come in my glory, Shall the palable [parable] be fulfilled of which I spoke concerning the ten virgins for he that is wise \& hath received the truth \& has taken the Holy Spirit for their guide \& have not been deceived Verily I say unto you they shall not be hewn down \& cast into the fire but shall abide the day \& the Earth shall be given unto them for an inheritance \& they shall multiply \& wax strong \& their children shall g$\diamond\diamond\diamond$ [grow] up without Sin unto salvation for the Lord shall be in their midst \& his glory shall be upon them \& he shall be their King \& their law giver \& now behold I say unto you it shall not be given unto you to know any farther then this until the New Testament be translated \& \sout{it} in it all things shall be made known Wherefore I give unto you that ye may now Translate it that ye may be prepared for the things to come for Verily I say unto you that great things await you ye hear of wars in foreign lands but behold I say unto you they are nigh even unto your doors \& not many years hence ye Shall hear of wars in your own lands wherefore I the Lord have said gether ye out from the Eastern lands assemble ye yourselves together ye Elders of my Church ge [go] ye forth into the western countries call upon the inhabitants to repent \& in as much as they do repent build up Churches unto me \& with one heart \& with one mind gether up your riches that you may purchase an inheritance which shall hereafter be appointed you \& it shall be called the New Jerusalem a land of peace a City of refuge a place of safety for the saints of The most high God \& the glory of the Lord shall be there \& the terer of the Lord also shall be there insomuch that the wicked will not come unto it \& it shall be called Zion \& it shall come to pass among the wicked that evry man that will not take his sword against his Neighbour must needs flee unto Zion for safety \& there shall be getherd unto it out of evry Nation under Heaven \& it shall be the only people that shall not be at war one with another \& it shall be said among the wicked let us not go up to battle against Zion for the inhabitants of Zion are terible wherefore we cannot stand \& it shall come to pass that the righteous shall be gethered out from among all Nations \& shall come to Zion singing with songs of everlasting Joy \& now I say unto you keep these things from going abroad unto the world that ye may accomplish this work in the eyes of the people \& in the eyes of your enemies that they may not know your works untill ye have accomplished the thing which I have commanded you that when they shall know it it may be terible unto them that fear may sieze upon them \& they shall stand afar off \& tremble \& all nations shall be afraid because of the teror of the lord \& the power of his might even so amen
				
		% -----
		\subsubsection{Revelation, circa 8 March 1831--A [D\&C 46]}
		% -----
			\label{EMS-1831-JS-CIR-8-MAR-1831-A}
			% \label{EMS-1830-JS-DAY-3LETTERMONTH-YEAR}
		
			\begin{description}
				\item[Print Sources]
			\end{description}

				\begin{tabular}{p{0.5in}p{1in}p{0.5in}p{1in}}
				JSP	 	& D1:280--283	& JSR 		& 124--126 \\
				DC	 	& Section 46 	& HC 		& 1:163--165 \\
				HDDC 	& 610--619		& TCHC 		& 1:120--121 \\
				\end{tabular}
				% HDDC = woodford's DC diss; JSR = Marquardt's JS Revelations; TCHC = Vogel HC
			
			\begin{description}
				\item[Online Access] \href{https://www.josephsmithpapers.org/paper-summary/revelation-circa-8-march-1831-a-dc-46/1}{Here}
				% \item[Analysis] \hyperref[DC46]{Here}
			\end{description}
			
			\begin{description}
				\item[Background]
			\end{description}
			
				% It might also be useful to give a two sentence context of the period: e.g., "This speech was given in Nauvoo to a small congregation one month prior to the destruction of the Nauvoo Expositor. These and other tensions may have been on the prophet's mind when he gave the speech."
			
			\begin{description}
				\item[Original Source]
			\end{description}
			
				% brief description of original source material, with reference to JSP or somewhere else for more info
				% Background material: it might be helpful to have a note that says whether a given document was presented as a speech, was written in a journal, etc.
				% Also a comment here about how reliable the source might be, or what shape it is in, if there are large omissions or parts that are hard to understand, and so on.
			
			\begin{description}
				\item[Text]
			\end{description}
			
				\begin{tightcenter}
				48\supers{th} Commandment March $\diamond$\supersu{th}\supers{.} 1831
				\end{tightcenter}
				
				given at Kirtland geauga County Ohio to the Church concerning conformation \& sacrament meetings \&c
				
				Hearken oh ye my people of my Church for Verily I say unto you that these things are spoken unto you for your profit \& learning but notwithstanding those things which are written it always has been given to the Elders of my Church from the begining \& ever shall be to conduct all meetings as they are conducted \& guided by the Holy spirit nevertheless ye are commanded never to cast any one out from your publick meetings which are held before the world ye are also commanded never to cast any one out who belongeth to the Church out of your sacrament meetings nevertheless if any have trespassed let him not partake untill he makes reconciliation And again I say unto you ye shall not cast any out of your sacrement meetings who is earnestly seeking the Kingdom I speak this concerning \sout{this} those who are not of the Church And again I say unto you concerning your confirmation meetings that if there be any that is not of the Church that is earnestly seeking after the Kingdom ye Shall not cast them out but ye are commanded in all things to ask of God who giveth liberally \& that which the spirit testifies unto you even so I would that ye should do in all Holyness of heart walking uprightly before me considering the end of your salvation doing all things with prayer \& thanksgiving that ye may not be seduced by evil spirits or doctrines of Devils or the commandments of men for some are of men \& others of Devils Wherefore beware lest ye are deceived \& that ye may not be deceived seek ye earnestly the best gifts always remembering for what they \sout{were} are given for verily I say unto you they are given for the benefit of those who love me \& keep all my commandments \& him that seeketh so to do that all may be benefitted that seeketh or that asketh of me that asketh \& not for a sign that he may consume it upon his lusts And again Verily I say unto you I would that ye should always remember \& always retain in your minds what these gifts are that are given unto the Church for all have not every gift given unto them for there are many gifts \& to evry man is given a gift by the spirit of God to some is given one \& to some is given another that all may be profited thereby to some is given by the Holy Ghost to know that Jesus Christ is the son of God \& that he was crusified for the sins of the World to others it is given to believe on their words that they also \sout{may} might have eternal life if they continue faithful And again to some it is given by the Holy Ghost to know the Defferences of administeration as it will be pleasing unto the same Lord according as the Lord will suiting his mercies according to the conditions of the children of men And again it is given by the Holy Ghost to some to know the diversities of opperations whether it be of God or not so that the manifestations of the spirit may be given to evry man to prophet [profit] withall And again Verily I say unto you to some it is given by the spirit of God the word of wisdom to another it is given the word of Knowledge that all may be taught to be wise \& to have knowledge \& again to some it is given to have faith to be healed \& to others it is given to have faith to heal And again to some it is given the working of miracles \& to others it is given to prophecy \& to others the decerning of spirits \& again it is given to some to speak with tongues \& to another it is given the interpretation of tongues \& all these gifts cometh from [God] for the benefit of the children of God \& unto the Bishop of the Church \& unto such as God sahall appoint \& ordain to watch over the Church \& to <be> \sout{the} Elders unto the Church are to have it given unto [them] to decern all those gifts lest there shall be any among you prophecying \& yet not be of God \& it shall come to pass that he that asketh in spirit shall receive in spirit that unto some it may be given to have all those gifts that there may be a head in order that evry member may be \sout{propheted} profited thereby he that asketh in spirit asketh according to the will of God wherefore it is done even as he asketh \& again I say unto you all things must be done in the name of Christ whatsoever you do in the spirit \& ye must give thanks unto God in the spirit for whatsoever blessing ye are blessed with \& ye must practice virtue \& holyness before me continually even so amen
				
		% -----
		\subsubsection{Revelation, circa 8 March 1831--B [D\&C 47]}
		% -----
			\label{EMS-1831-JS-CIR-8-MAR-1831-B}
			% \label{EMS-1830-JS-DAY-3LETTERMONTH-YEAR}
		
			\begin{description}
				\item[Print Sources]
			\end{description}

				\begin{tabular}{p{0.5in}p{1in}p{0.5in}p{1in}}
				JSP	 	& D1:284--286	& JSR 		& 126 \\
				DC	 	& Section 47 	& HC 		& 1:166 \\
				HDDC 	& 620--627		& TCHC 		& 1:121 \\
				\end{tabular}
				% HDDC = woodford's DC diss; JSR = Marquardt's JS Revelations; TCHC = Vogel HC
			
			\begin{description}
				\item[Online Access] \href{https://www.josephsmithpapers.org/paper-summary/revelation-circa-8-march-1831-b-dc-47/1}{Here}
				% \item[Analysis] \hyperref[DC47]{Here}
			\end{description}
			
			\begin{description}
				\item[Background]
			\end{description}
			
				% It might also be useful to give a two sentence context of the period: e.g., "This speech was given in Nauvoo to a small congregation one month prior to the destruction of the Nauvoo Expositor. These and other tensions may have been on the prophet's mind when he gave the speech."
			
			\begin{description}
				\item[Original Source]
			\end{description}
			
				% brief description of original source material, with reference to JSP or somewhere else for more info
				% Background material: it might be helpful to have a note that says whether a given document was presented as a speech, was written in a journal, etc.
				% Also a comment here about how reliable the source might be, or what shape it is in, if there are large omissions or parts that are hard to understand, and so on.
			
			\begin{description}
				\item[Text]
			\end{description}
			
				\begin{tightcenter}
				50\supers{th} Commandment March 8\supersu{th}\supers{.} 1831
				\end{tightcenter}
				
				Given at Kirtland Geauga Ohio--- given to John [Whitmer] in consequenc[e] of not feeling reconsiled to write at the request of Joseph with[o]ut a commandment \&c
				
				Behold it is expedient that my servent John should write \& keep a regulal [regular] history \& assist my servent Joseph in Transcribing all things which shall be given him And again verily I say unto you that ye \sout{can} can also lift Up your voice in Meetings when ever it shall be expedient \& again I say unto you that it shall be appointed unto you to Keep the Church Record \& History continually for Oliver [Cowdery] I have appointed to an other office wherefore it shall be given thee by th[e] comforter to write these things even so amen
				
		% -----
		\subsubsection{Revelation, 10 March 1831 [D\&C 48]}
		% -----
			\label{EMS-1831-JS-10-MAR-1831}
			% \label{EMS-1830-JS-DAY-3LETTERMONTH-YEAR}
		
			\begin{description}
				\item[Print Sources]
			\end{description}

				\begin{tabular}{p{0.5in}p{1in}p{0.5in}p{1in}}
				JSP	 	& D1:286--288	& JSR 		& 127 \\
				DC	 	& Section 48 	& HC 		& 1:166--167 \\
				HDDC 	& 628--633		& TCHC 		& 1:122 \\
				\end{tabular}
				% HDDC = woodford's DC diss; JSR = Marquardt's JS Revelations; TCHC = Vogel HC
			
			\begin{description}
				\item[Online Access] \href{https://www.josephsmithpapers.org/paper-summary/revelation-10-march-1831-dc-48/1}{Here}
				% \item[Analysis] \hyperref[DC48]{Here}
			\end{description}
			
			\begin{description}
				\item[Background]
			\end{description}
			
				% It might also be useful to give a two sentence context of the period: e.g., "This speech was given in Nauvoo to a small congregation one month prior to the destruction of the Nauvoo Expositor. These and other tensions may have been on the prophet's mind when he gave the speech."
			
			\begin{description}
				\item[Original Source]
			\end{description}
			
				% brief description of original source material, with reference to JSP or somewhere else for more info
				% Background material: it might be helpful to have a note that says whether a given document was presented as a speech, was written in a journal, etc.
				% Also a comment here about how reliable the source might be, or what shape it is in, if there are large omissions or parts that are hard to understand, and so on.
			
			\begin{description}
				\item[Text]
			\end{description}
			
				\begin{tightcenter}
				49 Commandment March 10\supersu{th}\supers{.} 1831
				\end{tightcenter}
				
				A Revelation Received at Kirtland Geauga County Ohio concerning the Bretheren in New York how to Manage with their property \&c
				
				It is nessessary that ye should remain for the present time in your places of abode as it shall be suitable to your circumstances \& inasmuch as ye [have] lands ye shall impart to the Eastern Brethren \& in as much as ye have not lands let them buy for the present time in those regions round about as seemeth them good for it must needs be nessessary that they have places to live for the present time it must needs be nessessary that ye save all the money that ye can (\& that ye obtain all that ye can) that in time ye may be enabled to purchase lands for an inheritance (even the City) the place is not yet to be revealed but after your Brethren come from the East there are to be certain men to be appointed \& to them it shall be given to know the place as to them it shall be revealed \& they shall be appointed to purchase the lands \& to \sout{lay the foundation} make a commencement to lay the foundation of the City \& then ye shall begin to be gethered with your families evry man according to his family according to his circumstances \& as is appointed to them by the Bishop \& Elders of the Church according to the laws \& commandments which ye have received \& which ye shall hereafter receive
				
		% -----
		\subsubsection{Revelation, 7 May 1831 [D\&C 49]}
		% -----
			\label{EMS-1831-JS-7-MAY-1831}
			% \label{EMS-1830-JS-DAY-3LETTERMONTH-YEAR}
		
			\begin{description}
				\item[Print Sources]
			\end{description}

				\begin{tabular}{p{0.5in}p{1in}p{0.5in}p{1in}}
				JSP	 	& D1:297--303	& JSR 		& 128--129 \\
				DC	 	& Section 49 	& HC 		& 1:167--169 \\
				HDDC 	& 634--642		& TCHC 		& 1:122--123 \\
				\end{tabular}
				% HDDC = woodford's DC diss; JSR = Marquardt's JS Revelations; TCHC = Vogel HC
			
			\begin{description}
				\item[Online Access] \href{https://www.josephsmithpapers.org/paper-summary/revelation-7-may-1831-dc-49/1}{Here}
				% \item[Analysis] \hyperref[DC49]{Here}
			\end{description}
			
			\begin{description}
				\item[Background]
			\end{description}
			
				% It might also be useful to give a two sentence context of the period: e.g., "This speech was given in Nauvoo to a small congregation one month prior to the destruction of the Nauvoo Expositor. These and other tensions may have been on the prophet's mind when he gave the speech."
			
			\begin{description}
				\item[Original Source]
			\end{description}
			
				% brief description of original source material, with reference to JSP or somewhere else for more info
				% Background material: it might be helpful to have a note that says whether a given document was presented as a speech, was written in a journal, etc.
				% Also a comment here about how reliable the source might be, or what shape it is in, if there are large omissions or parts that are hard to understand, and so on.
			
			\begin{description}
				\item[Text]
			\end{description}
			
				\begin{tightcenter}
				51\supersu{st}\supers{.} Commandment May 7\supersu{th}\supers{.} 1831
				\end{tightcenter}
				
				A Revelation given to Sidney [Rigdon] \& Parley [P. Pratt] \& Leman [Copley] Received at Kirtland Geauga Ohio th[e]ir mission to the Shakers \& thus saith the Lord unto them as follows
				
				Hearken unto my word my Servent Sidney \& Parley \& Leman for Behold verily I say unto you that I give unto you a commandment that you shall go \& preach my Gospel which ye have received even as ye have received it unto the Shakers. Behold I say unto you that they desire to know the truth in Part but not all for they are not right before me \& must needs repent wherefore I send you my Servents sidney \& Parley to preach the Gospel unto them \& my servent Leman shall be ordained unto this work that he may reason with them not according to that which he hath received of them but according to that which shall be taught \sout{them} him by you my Servents \& by so doing I will bless him otherwise he shall not prosper thus saith the Lord for I am God \& have sent mine only begotten Son into the world for the redemption of the world \& have decreed that he that receiveth him shall be saved \& he that receiveth him not shall be damned \& they have done unto the Son of man even as they listed \& he hath taken his power on the right hand of his glory \& now reigneth in the Heavens \sout{<\& will>} till he decends on the Earth to put all enemies under his feet which time is nigh at hand I the Lord hath spoken it but the hour \& the day no man knoweth neither the angels in Heaven nor shall they know untill he come wherefore I will that all men <shall> repent for all are under sin except them which I have reserved unto myself Holy men that ye know not of wherefore I say unto you that I have sent unto you mine everlasting Covenant even that which was from the begining \& that which I have promised I have so fulfilled \& the Nations of the earth shall bow to it \& if not of them selves they shall come down for that which is now exalted of itself shall be laid low of power. wherefore I give unto you a commandment that ye go among this People \& say unto them like unto mine Apostle of old whose name was Peter believe on the name of the Lord Jesus who was on the Earth \& is to come the begining \& the end. Repent \& be baptized in the name of Jesus christ according to the holy commandment for the remission of sins \& whoso doeth this shall receive the gift of the Holy Ghost by the laying on of the hands of the Elders of this Church. \& again I say unto you that whoso forbideth to marry is not ordained of God for it is ordained of God unto man wherefore it is lawful that he should have one wife \& they twain shall be one flesh \& all this that the Earth might answer the end of its Creation \& that it might be filled with the measure of man according to his creation before the world was made \& whoso forbideth to abstain from meats that man should not eat the same is not ordained of God for behold the beasts of the field \& the fowls of the air \& that which cometh of the Earth is ordained for the use of man for food \& for raiment \& that he might have in abundance but it is not given that one man should possess that which is above an other wherefore the world lieth in sin \& wo be unto man that shedeth blood or that wasteth flesh \& hath no need \& again verily I say unto you that the son of man cometh not in the form of a woman neither of <a> man traveling on the Earth wherefore be not deceived but continue in steadfastness looking forth for the Heavens to <be> shaken \& the Earth to tremble \& to reel to \& fro as a drunken man \& for the vallies to be exhalted \& for the Mountains to be made low \& for the rough places to become smooth \& all this when the Angel shall sound his trumpet but before this great day of the Lord shall come Jacob shall flourish in the wilderness \& the Lamanit[e]s shall blossom as the rose \& Zion shall flourish upon the Hills \& rejoice upon the Mountains \& shall be assembled together Unto the place which I have appointed behold I say unto you go forth as I have commanded you repent of all your sins ask \& ye shall receive knock \& it shall be opened unto you behold I will go before you \& be your rearward \& I will be in your midst \& you shall not be confounded behold I am Jesus Christ \& I come quickly even so Amen------
				
		% -----
		\subsubsection{Revelation, 9 May 1831 [D\&C 50]}
		% -----
			\label{EMS-1831-JS-9-MAY-1831}
			% \label{EMS-1830-JS-DAY-3LETTERMONTH-YEAR}
		
			\begin{description}
				\item[Print Sources]
			\end{description}

				\begin{tabular}{p{0.5in}p{1in}p{0.5in}p{1in}}
				JSP	 	& D1:303--308	& JSR 		& 130--132 \\
				DC	 	& Section 50 	& HC 		& 1:170--173 \\
				HDDC 	& 643--663		& TCHC 		& 1:124--125 \\
				\end{tabular}
				% HDDC = woodford's DC diss; JSR = Marquardt's JS Revelations; TCHC = Vogel HC
			
			\begin{description}
				\item[Online Access] \href{https://www.josephsmithpapers.org/paper-summary/revelation-9-may-1831-dc-50/1}{Here}
				% \item[Analysis] \hyperref[DC50]{Here}
			\end{description}
			
			\begin{description}
				\item[Background]
			\end{description}
			
				% It might also be useful to give a two sentence context of the period: e.g., "This speech was given in Nauvoo to a small congregation one month prior to the destruction of the Nauvoo Expositor. These and other tensions may have been on the prophet's mind when he gave the speech."
			
			\begin{description}
				\item[Original Source]
			\end{description}
			
				% brief description of original source material, with reference to JSP or somewhere else for more info
				% Background material: it might be helpful to have a note that says whether a given document was presented as a speech, was written in a journal, etc.
				% Also a comment here about how reliable the source might be, or what shape it is in, if there are large omissions or parts that are hard to understand, and so on.
			
			\begin{description}
				\item[Text]
			\end{description}
			
				\begin{tightcenter}
				52\supers{nd} Commandment May 9\supersu{th}\supers{.} 1831
				\end{tightcenter}
				
				A Revelation to the Elders of this Church given at Kirtland geauga Ohio in consequence of their not being perfectly acquainted with the different opperations of the Spirits which are abroad in the Earth \& thus saith the Lord unto them as follows
				
				Hearken o ye Elders of my Church \& give ear to the voice of the living God \& attend to the words of wisdom which shall be given unto you according as ye have asked \& are agreed as touching the Church \& the spirits which have gone abroad in the Earth Behold verily I say unto you that there are many spirits which are false spirits which have gone forth in the Earth deceiving the world \& also Satan hath sought to deceive you that he might overthrow you Behold I the Lord have looked upon you \& have seen abominations in the Church which profess my name but blessed are they who are faithfull \& endure whether in life or in death for they shall inherit eternal life but wo be unto them that are deceivers \& hypocrites for thus saith the Lord I will bring them to Judgement behold verily I say unto you there are hypocrites among you \& have deceived some which have given the adve[r]sary power but behold such shall be reclaimed but the hypocrites shall be detected \& shall be cut off either in life or in death even as I will \& wo is them that is cut off from my Church for the same is overcome of the world wherefore let every man be aware lest he do that which is not in truth \& righteousness before me \& now come saith the Lord by the spirit by the Elders of his Church \& let us reason together that ye may understand let us reason even as a man reasoneth one with another face to face now when a man reasoneth he under[stand]eth of man because he reasoneth as a man even so will I the Lord reason with you that you may understand wherefore I the Lord asketh you this question unto what were ye ordained to Preach my Gospel by the spirit even the comforter which was sent forth to teach the truth \& then received you spirits which ye could not understand \& received them to be of God \& in this are ye Justified? Behold ye shall answer this yourselves nevertheles I will be mercyfull unto you he that is weak among you hereafter shall be made strong verily I say unto you he that is ordained of me \& sent forth to preach the word of truth by the comforter the spirit of truth doth he preach it by the spirit of truth or some other way \& if by some other way it be not of God \& again he that receiveth the word of truth doth he receive it by the spirit of truth or some other way if it be some other way it be not of God therefore why is it that ye cannot understand \& know that he <that> receiveth the word by the spirit of truth receiveth it as it is preached by the spirit of truth wherefore he that <preacheth \& he that> receiveth understandeth one another \& both are edified \& rejoice together \& that which doth not edify is not of God \& is darkness that which is of God is light \& he that receiveth light \& continueth in god receiveth more light \& that light groweth \sout{lighter} brighter and brighter untill the perfect day \& again verily I say unto you \& I say it that you may know the truth that you may \sout{choose} chase darkness from among you for he that is ordaned of God \& sent forth the same is apponted to be the greatest notwithstanding he is least \& the servent of all wherefore he is possessor of all things <for all things> are subject unto him both in Heaven \& on the Earth the life \& the light the spirit \& the power sent forth by the will of the father through Jesus Christ his son but no man is possessor of all things except he be purified \& cleansed from all sin \& if ye are purified \& cleansed from all sin ye shall ask whatsoever you will in the name of Jesus \& it shall be done but know this it shall be given you what ye shall ask \& as ye are appointed to the head The spirits shall be subject unto you wherefore it shall come to pass that if ye behold a spirit manifested that ye cannot understand \& you receive not that spirit ye shall ask of the father in the name of Jesus \& if he \sout{gave} give not unto you that spirit\sout{s} then ye may know that it is not of God \& it shall be given unto you power \sout{that spi} over that spirit \& you shall proclaim against that spirit with a loud voice that it is not of God not with railing accusation that ye be not over come neither with boasting nor rejoicing lest you be seased [seized] therewith he that receiveth of God let it [him] account it of God \& let him rejoice that he is accounted of God worthy to receive \& by giving heed \& doing these things which ye have received \& which ye shall hereafter receive \textit{[illegible]} the Kingdom is given unto you of the father \& power to over come all things which is not ordained of him \& Behold verily I say unto you blessed are you that hear these words of mine from the mouth of my servent for your sins are forgiven you Let my Servent Joseph [Wakefield] in whom I am well pleased \& my servent Parley [P. Pratt] go forth among the Churches \& strengthen them by the word of exhortation \& also my servent John [Corrill] or as many of my servents as are ordained unto this office \& let them labour in the vinyard \& let no man hinder them of doing that which I have appointed unto them wherefore in this thing my Servent Edward [Partridge] is not Justified nevertheless let him repent \& he shall be forgiven------ Behold ye are little Children \& ye cannot bear all things now ye must grow in grace \& in the knowledge of the truth fear not little children for you are mine \& I have overcome the world \& you are of them which my father hath given me \& none of them which my father hath given me shall be lost \& the father \& I are one I am in the father \& the father in me \& I in \sout{you} as much as ye have received me ye are in me \& \sout{am} I in you wherefore I am in your midst \& I am the good shepherd \& the day cometh that you shall hear my voice \& see me \& know that I am watch therefore that ye may be ready even So Amen
				
		% -----
		\subsubsection{Revelation, 15 May 1831}
		% -----
			\label{EMS-1831-JS-15-MAY-1831}
			% \label{EMS-1830-JS-DAY-3LETTERMONTH-YEAR}
		
			\begin{description}
				\item[Print Sources]
			\end{description}

				\begin{tabular}{p{0.5in}p{1in}p{0.5in}p{1in}}
				JSP	 	& D1:309--314	& JSR 		& 135--136 \\
				DC	 	& ---		 	& HC 		& --- \\
				HDDC 	& ---			& TCHC 		& --- \\
				\end{tabular}
				% HDDC = woodford's DC diss; JSR = Marquardt's JS Revelations; TCHC = Vogel HC
			
			\begin{description}
				\item[Online Access] \href{https://www.josephsmithpapers.org/paper-summary/revelation-15-may-1831/1}{Here}
				% \item[Analysis] \hyperref[XX]{Here}
			\end{description}
			
			\begin{description}
				\item[Background]
			\end{description}
			
				% It might also be useful to give a two sentence context of the period: e.g., "This speech was given in Nauvoo to a small congregation one month prior to the destruction of the Nauvoo Expositor. These and other tensions may have been on the prophet's mind when he gave the speech."
			
			\begin{description}
				\item[Original Source]
			\end{description}
			
				% brief description of original source material, with reference to JSP or somewhere else for more info
				% Background material: it might be helpful to have a note that says whether a given document was presented as a speech, was written in a journal, etc.
				% Also a comment here about how reliable the source might be, or what shape it is in, if there are large omissions or parts that are hard to understand, and so on.
			
			\begin{description}
				\item[Text]
			\end{description}
			
				\begin{tightcenter}
						\hypertarget{EMS-1831-JS-15-MAY-1831-a}%
					53
						\marginnote{a}
					Commandment May 15\supersu{th}\supers{.} 1831
				\end{tightcenter}
				
				givn to Ezra Thayer \& Joseph Smith [Sr.] concerning a farm \&c
				
				Hearken unto my words \& behold I will make known unto you what ye shall do as it shall be pleasing unto me for verily I say unto you it must needs be that ye let the bargain stand that ye have made concerning these farms untill it be so fulfilled Behold ye are holden for the one even so likewise thine advisary is holden for the other wherefore it must needs be that ye pay no more money for the present time untill the contract be fulfilled
					\hypertarget{EMS-1831-JS-15-MAY-1831-b}%
				\phantom{ }
					\marginnote{b}%
				\phantom{ }\& let my Servent Joseph [Smith Sr.] \& his family go into the House after thine advisary is gone \& let my Servent Ezra board with him \& let all the Brethren immediately assemble together \& put up an house for my Servent Ezra \& let my Servent\sout{s} Frederick [G. Williams]'s family remain \& let the house be repaired \& their wants be supplied \& when my Servent Frederick returns from the west Behold he taketh his family to the west Let that which belongeth to my Servent Frederick be secured unto him by deed or bond \& thus he willeth that the Brethren reap the good thereof let my Servent Joseph [Smith Sr.] govern the things of the farm \& provide for the families \& let him have help in as much as he standeth in need
					\hypertarget{EMS-1831-JS-15-MAY-1831-c}%
				let
					\marginnote{c}%
				my servent Ezra humble himself \& at the conference meeting he shall be ordained unto power from on high \& he shall go from thence (if he be obedient unto my commandments) \& proclaim my Gospel unto the western regions with my Servents that must go forth even unto the borders of the Lamanit[e]s for Behold I have a great work for them to do \& it shall be given unto you to know what ye shall do at the conferenc[e] meeting even so Amen------
				
					\hypertarget{EMS-1831-JS-15-MAY-1831-d}%
				What
					\marginnote{d}%
				shall the Brethren do with their money------
				
				Ye shall go forth \& seek dilligently among the Brethren \& obtain lands \& save the money that it may be consecrated to purcchase lands in the west for an everlasting enheritance
				
				even So Amen
				
		% -----
		\subsubsection{Revelation, 20 May 1831 [D\&C 51]}
		% -----
			\label{EMS-1831-JS-20-MAY-1831}
			% \label{EMS-1830-JS-DAY-3LETTERMONTH-YEAR}
		
			\begin{description}
				\item[Print Sources]
			\end{description}

				\begin{tabular}{p{0.5in}p{1in}p{0.5in}p{1in}}
				JSP	 	& D1:314--317	& JSR 		& 133--135 \\
				DC	 	& Section 51 	& HC 		& 1:173--174 \\
				HDDC 	& 664--674		& TCHC 		& 1:125--126 \\
				\end{tabular}
				% HDDC = woodford's DC diss; JSR = Marquardt's JS Revelations; TCHC = Vogel HC
			
			\begin{description}
				\item[Online Access] \href{https://www.josephsmithpapers.org/paper-summary/revelation-20-may-1831-dc-51/1}{Here}
				% \item[Analysis] \hyperref[DC51]{Here}
			\end{description}
			
			\begin{description}
				\item[Background]
			\end{description}
			
				% It might also be useful to give a two sentence context of the period: e.g., "This speech was given in Nauvoo to a small congregation one month prior to the destruction of the Nauvoo Expositor. These and other tensions may have been on the prophet's mind when he gave the speech."
			
			\begin{description}
				\item[Original Source]
			\end{description}
			
				% brief description of original source material, with reference to JSP or somewhere else for more info
				% Background material: it might be helpful to have a note that says whether a given document was presented as a speech, was written in a journal, etc.
				% Also a comment here about how reliable the source might be, or what shape it is in, if there are large omissions or parts that are hard to understand, and so on.
			
			\begin{description}
				\item[Text]
			\end{description}
			
				\begin{tightcenter}
				54 Commandment------
				\end{tightcenter}
				
				54 <Com> A Revelation given to the Bishop at Thompson Ohio May 20\supersu{th}\supers{.} 1831 concerning the property of the Church \&c
				
				Hearken unto me saith the lord your God \& I will speak unto my Servent Edward [Partridge] \& give unto him directions for it must needs be that he receive directions how to organize this people for it must needs be that they are organized according to my laws if otherwise they will be cut off wherefore let my Servent Edward receive the properties of this People which have covenanted with me to obey the Laws which I have given \& let my Servent Edward receive the money as it shall be laid before him according to the covenant \& go \& obtain a deed or Article of this land unto himself for I have appointed him to receive these things \& thus through him the Properties of this Church shall be covenanted unto me wherefore let my Servent Edward \& those whom he has chosen in whom I am well pleased appoint unto this People their portion every man alike according to their families according to their wants \& their needs \& let my servent Edward when he shall appoint a man his portion give unto him a writing that shall secure unto him his portion that he shall hold it of the Church untill he transgress \& is not counted worthy by the \sout{Church} voice of the Church according to the laws to belong to the Church \& thus all things shall be made sure according to the laws of the land \& let that which belongeth to this people be appointed unto this people \& the money which is left unto this people let there be an agent appointed unto this people to take the money to provide food \& raiment according to the wants of this people \& let every man deal honestly \& be alike among \sout{you} this People \& receive alike that ye may be one even as I have commanded you \& let that which belongeth to this people not be taken \& given unto that of another church wherefore if another Church would receive money of this Church let them pay unto this church again according as they shall agree \& this shall be done through the Bishop or the agent which shall be appointed by the voice of the church \& again let the Bishop appoint a storehouse unto this Church \& let all things both in in money \& in meat which is more then is needful for the want of this People be kept in the hands of the Bishop \& let him also reserve unto himself for his own wants \& for the wants of his family as he shall be employed in doing this Business \& thus I grant unto this People a privelige of organizeing themselves according to my laws \& I consecrate unto them this land for a little season untill I the Lord shall provide for them otherwise \& command them to go hence \& the hour \& the day is not given unto them wherefore let them act upon this land as for years \& this shall turn unto them for their good Behold this shall be an example unto my Servent Edward in other places in all Churches \& whoso is found a faithful \& Just \& a wise stewart shall enter into the Joy of his lord \& shall inherit eternal life verily I say unto you I am Jesus Christ who cometh quickly in an hour you think not even so Amen
				
		% -----
		\subsubsection{Minutes, circa 3--4 June 1831}
		% -----
			\label{EMS-1831-JS-CIR-3-4-JUN-1831}
			% \label{EMS-1830-JS-DAY-3LETTERMONTH-YEAR}
		
			\begin{description}
				\item[Print Sources]
			\end{description}

				\begin{tabular}{p{0.5in}p{1in}p{0.5in}p{1in}}
				JSP	 	& ---		& JSR 		& --- \\
				DC	 	& --- 		& HC 		& --- \\
				HDDC 	& ---		& TCHC 		& --- \\
				\end{tabular}
				% HDDC = woodford's DC diss; JSR = Marquardt's JS Revelations; TCHC = Vogel HC
			
			\begin{description}
				\item[Online Access] \href{https://www.josephsmithpapers.org/paper-summary/minutes-circa-3-4-june-1831/1}{Here}
				% \item[Analysis] \hyperref[DC52]{Here}
			\end{description}
			
			\begin{description}
				\item[Background]
			\end{description}
			
				% It might also be useful to give a two sentence context of the period: e.g., "This speech was given in Nauvoo to a small congregation one month prior to the destruction of the Nauvoo Expositor. These and other tensions may have been on the prophet's mind when he gave the speech."
			
			\begin{description}
				\item[Original Source]
			\end{description}
			
				% brief description of original source material, with reference to JSP or somewhere else for more info
				% Background material: it might be helpful to have a note that says whether a given document was presented as a speech, was written in a journal, etc.
				% Also a comment here about how reliable the source might be, or what shape it is in, if there are large omissions or parts that are hard to understand, and so on.
			
			\begin{description}
				\item[Text]
			\end{description}
			
					\hypertarget{EMS-1831-JS-CIR-3-4-JUN-1831-a}%
				...Conference
					\marginnote{a}%
				opened by br. Joseph Smith jr. in exhortation \& prayer, Prayer by br. Sidney Rigdon \& exhortation by the same. Exhortation by most of the Elders present.

				Brs. Lyman Wight John Murdock Reynolds Cahoon Harvey Whitlock \& Hyrum Smith were ordained to the high Priesthood under the hand <of> br. Joseph Smith jr.

				Exhortation by brs. Lyman Wight \& Harvey Whitlock

					\hypertarget{EMS-1831-JS-CIR-3-4-JUN-1831-b}%
				Brs.
					\marginnote{b}%
				Parley P. Pratt Thomas B. Marsh Isaac Morley Edward Partridge Joseph Wakefield Martin Harris Ezra Thayer Ezra Booth, (denied the faith) John Corrill Samuel H. Smith Solomon Hancock Simeon Carter Wheeler Baldwin Jacob Scott, (denied the faith) Joseph Smith [Sr.] John Whitmer Joseph Smith jr. \& Sidney Rigdon were ordained to the High Priesthood under the hand of br. Lyman Wight; The Bishop, (Edward Partridge) then blessed those who were ordained in the name of Christ according to commandment br. John Corrill \& Isaac Morley were ordained assistants to the Bishop under the hand of Lyman Wight.

				Exhortation by br. Sidney Rigdon \& Joseph Smith jr. Closed by prayer by br. Sidney Rigdon.

				\begin{tightright}
					John Whitmer Clerk
				\end{tightright}
				
		% -----
		\subsubsection{Revelation, 6 June 1831 [D\&C 52]}
		% -----
			\label{EMS-1831-JS-6-JUN-1831}
			% \label{EMS-1830-JS-DAY-3LETTERMONTH-YEAR}
		
			\begin{description}
				\item[Print Sources]
			\end{description}

				\begin{tabular}{p{0.5in}p{1in}p{0.5in}p{1in}}
				JSP	 	& D1:327--332	& JSR 		& 136--138 \\
				DC	 	& Section 52 	& HC 		& 1:177--179 \\
				HDDC 	& 675--692		& TCHC 		& 1:127--129 \\
				\end{tabular}
				% HDDC = woodford's DC diss; JSR = Marquardt's JS Revelations; TCHC = Vogel HC
			
			\begin{description}
				\item[Online Access] \href{https://www.josephsmithpapers.org/paper-summary/revelation-6-june-1831-dc-52/1}{Here}
				% \item[Analysis] \hyperref[DC52]{Here}
			\end{description}
			
			\begin{description}
				\item[Background]
			\end{description}
			
				% It might also be useful to give a two sentence context of the period: e.g., "This speech was given in Nauvoo to a small congregation one month prior to the destruction of the Nauvoo Expositor. These and other tensions may have been on the prophet's mind when he gave the speech."
			
			\begin{description}
				\item[Original Source]
			\end{description}
			
				% brief description of original source material, with reference to JSP or somewhere else for more info
				% Background material: it might be helpful to have a note that says whether a given document was presented as a speech, was written in a journal, etc.
				% Also a comment here about how reliable the source might be, or what shape it is in, if there are large omissions or parts that are hard to understand, and so on.
			
			\begin{description}
				\item[Text]
			\end{description}
			
				\begin{tightcenter}
				55\supersu{th}\supers{.} Commandment given at Kirtland June 6\supersu{th}\supers{.} 1831
				\end{tightcenter}
				
				Directions to the Elders of the Church of Christ \&c
				
				Behold thus saith the Lord unto the Elders whom he hath called \& chosen in these last days by the voice of his Spirit saying I the Lord will make known unto you what I will make known that ye sh$\diamond\diamond$ld do from this time untill the next conference which shall be held in Missorie upon the land which I will consecrate unto my People which are a remnant of Jacob \& those who are heirs according to the covenant wherefore verily I say unto you let my Servents Joseph \& Sidney [Rigdon] take their Journey as Soon as preperations can be made to leave their homes \& Journey to the land of Missorie \& in as much as they are faithfull unto me it shall be made known unto them what they shall do \& it shall also in as much as they are faithfull be made known unto them the land of your inheritance \& in as much as they are not faithfull they shall be cut off even as I will as Seemeth me good \& again verily I say unto you let my Servent Lyman (Wight) \& my Servent John (Corrill) take their Journey speedily \& also my Servent John (Murdock) \& my Servent Hyram [Hyrum] (Smith) take their Journey unto the Same place by the way of Detroit \& let them Journey from thence preaching the word by the way saying none other things than the Prophets \& Apostles have written \& that which is Taught them by the Comforter through the prayer of faith[.] let them go two by two \& thus let them Preach by [the] way in every congregation Baptizing \& t by waterhe laying on the hands by the water side for thus sayeth the lord I will cut my work short in righteousness for the days cometh that I will send forth Judgement unto victory \& let my Servent Lyman be aware for Satan desireth to sift him as Chaff \& behold he that is faithfull shall be made ruler over many things \& again I will give unto you a Pattern in all things that ye may not be deceived for Satan is abroad in the land \& he goeth forth deceiveing the Nations wherefore he that prayeth whose spirit is contrite the same is accepted of me if he obey mine ordinances[.] he that speaketh whose spirit is contrite whose language \sout{in} is meek \& edifieth the Same is of God if he obey mine ordinances \& again he that trembeleth under my power shall be made strong \& shall bring forth fruits of Praise \& wisdom according to the Revelations \& truths which I have given you \& again he that is overcome \& bringeth not forth fruits even according to this Pattern is not of me wherefore by this Pattern ye shall know all the spirits in all cases under the whole Heavens \& the days have come according to mens faith it shall be done unto them behold this commandment is given unto all the Elders whom I have chosen \& again verily I say unto you let my Servent\sout{s} Thomas [B.] (Marsh) \& my Servent Ezra (Thayer) take their Journey also preaching the word by the way unto this same land \& again let my servent Isaac (Morl[e]y) \& my Servent Ezra (Booth) take their Journey also preaching the word by the way to the same land let my Servent Edward (Pa[r]tridge) \& Martin (Harris) take their Journey with my servents Sidney \& Joseph let my Servent David (Whitmer) \& Harvey Whitlock also take their Journey \& preach by the way let my Servent Parly (Pratt) [Parley P. Pratt] \& Or[s]on (Pratt) also take their Journey \& Preach by the way unto this same land \& my Servent Solomon (Hancock) \& simeon (Carter) also take their Journey to the same land \& preach by the way let my Servent Edson (Fuller) \& Jacob (Schott [Scott]) also take their Journey let my servent Levi Hancock \& Zebedee (Coltrin) take their Journey let my Servent Reynolds (Cahoon) \& Samuel (Smith) also take their Journey let my Servent Wheeler (Baldwin) \& william (Carter) also take their Journey let my Servent Newel Knight \& Sealy [Selah] (Griffin) both be ordained \& also take their Journey yea verily I say \sout{unto you} let all these take their Journey unto one Place in their Several courses \& one man shall not build upon anothers foundation neither Journey in an others tracks he that is faithfull the same shall be kept \& blest with much fruit \& again I say unto you let my servent Joseph (Wakefield) \& Solomon (Hmphrey) [Solomon Humphrey Jr.] take their Journey into the easteren lands \& let them labour with their families declaring none other things than the Prophets \& Apostles that which they have seen \& heard \& most shuredly believe that the Prophecies may be fulfilled[.] in consequence of transgression let that which was bestowed upon Heman (Bassett) be taken from him \& placed upon the head of \sout{Simonds} Simonds (Rider) \& again verily I say unto you let Jared Carter be ordained a Priest \& also George (James) be ordained a Priest let the residue of the Elders watch over the Churches \& declare the word in the regiones among them \& let them labour with their own hands that there be no Idolitry nor wickedness practiced \& remember in all things the poor \& the Needy the Sick \& the afflicted for he that doeth not these things is not my deciple these things the same is not my Deciple \& again let my Servent Joseph \& Sidney \& Edward take with them a recomend from the Church \& let there be one obtained for my Servent Oliver [Cowdery] also \& thus even as I have said if ye are faithfull ye shall assemble yourselves together to rejoice upon the land of your inheritance which is now the land of your enemies but behold I the lord will hasten the City in its time \& will crown the faithfull with Joy \& rejoicing Behold I am Jesus Christ the Son of God \& I will lift them up at the last day even so amen
				
		% -----
		\subsubsection{Revelation, 8 June 1831 [D\&C 53]}
		% -----
			\label{EMS-1831-JS-8-JUN-1831}
			% \label{EMS-1830-JS-DAY-3LETTERMONTH-YEAR}
		
			\begin{description}
				\item[Print Sources]
			\end{description}

				\begin{tabular}{p{0.5in}p{1in}p{0.5in}p{1in}}
				JSP	 	& D1:332--334	& JSR 		& 138--139 \\
				DC	 	& Section 53 	& HC 		& 1:179--180 \\
				HDDC 	& 693--699		& TCHC 		& 1:129 \\
				\end{tabular}
				% HDDC = woodford's DC diss; JSR = Marquardt's JS Revelations; TCHC = Vogel HC
			
			\begin{description}
				\item[Online Access] \href{https://www.josephsmithpapers.org/paper-summary/revelation-8-june-1831-dc-53/1}{Here}
				% \item[Analysis] \hyperref[DC53]{Here}
			\end{description}
			
			\begin{description}
				\item[Background]
			\end{description}
			
				% It might also be useful to give a two sentence context of the period: e.g., "This speech was given in Nauvoo to a small congregation one month prior to the destruction of the Nauvoo Expositor. These and other tensions may have been on the prophet's mind when he gave the speech."
			
			\begin{description}
				\item[Original Source]
			\end{description}
			
				% brief description of original source material, with reference to JSP or somewhere else for more info
				% Background material: it might be helpful to have a note that says whether a given document was presented as a speech, was written in a journal, etc.
				% Also a comment here about how reliable the source might be, or what shape it is in, if there are large omissions or parts that are hard to understand, and so on.
			
			\begin{description}
				\item[Text]
			\end{description}
			
				\begin{tightcenter}
				\sout{55} 56\supers{th} Commandment June 8\supersu{th}\supers{.} 1831
				\end{tightcenter}
				
				A Revelation to Sidney Gilbert his Call \&c
				
				Behold I say unto you my servent Sidney that I have heard your prayers \& ye have called upon me that it should be made known unto you of the lord your God concerning your calling \& election in this Church which I the Lord have raised up in these last days Behold I \sout{my} the Lord who was crusified for the sins of the world giveth unto you a commandment that you shall forsake the world take upon you mine ordinances even that of an Elder to Preach faith \& repentance \& remission of sins according to my word \& the reception of the Holy spirit by the laying on of hands \& also to be an agent unto this Church in the Place which shall be appointed by the Bishop according to commandments which shall be given hereafter \& again verily I say unto you you shall take your Journey with my servent Joseph \& sidney [Rigdon] Behold these are the <firist> ordinances which you shall receive \& the residue shall be made known unto you in a time to come according to your labour in my vinyard \& again I would that ye should learn that it is hi$\diamond$ only who is saved that endureth unto the end even so amen
				
		% -----
		\subsubsection{Revelation, 10 June 1831 [D\&C 54]}
		% -----
			\label{EMS-1831-JS-10-JUN-1831}
			% \label{EMS-1830-JS-DAY-3LETTERMONTH-YEAR}
		
			\begin{description}
				\item[Print Sources]
			\end{description}

				\begin{tabular}{p{0.5in}p{1in}p{0.5in}p{1in}}
				JSP	 	& D1:334--336	& JSR 		& 139 \\
				DC	 	& Section 54 	& HC 		& 1:181 \\
				HDDC 	& 700--708		& TCHC 		& 1:130 \\
				\end{tabular}
				% HDDC = woodford's DC diss; JSR = Marquardt's JS Revelations; TCHC = Vogel HC
			
			\begin{description}
				\item[Online Access] \href{https://www.josephsmithpapers.org/paper-summary/revelation-10-june-1831-dc-54/1}{Here}
				% \item[Analysis] \hyperref[DC54]{Here}
			\end{description}
			
			\begin{description}
				\item[Background]
			\end{description}
			
				% It might also be useful to give a two sentence context of the period: e.g., "This speech was given in Nauvoo to a small congregation one month prior to the destruction of the Nauvoo Expositor. These and other tensions may have been on the prophet's mind when he gave the speech."
			
			\begin{description}
				\item[Original Source]
			\end{description}
			
				% brief description of original source material, with reference to JSP or somewhere else for more info
				% Background material: it might be helpful to have a note that says whether a given document was presented as a speech, was written in a journal, etc.
				% Also a comment here about how reliable the source might be, or what shape it is in, if there are large omissions or parts that are hard to understand, and so on.
			
			\begin{description}
				\item[Text]
			\end{description}
			
				\begin{tightcenter}
				57\supersu{th}\supers{.} Commandment June 10\supersu{th}\supers{.} 1831
				\end{tightcenter}
				
				A Revelation to the Church at Thompson giving them Directions what to do \&c
				
				Behold thus saith the Lord even Alpha \& Omega the begining \& the end even him that was crusified for the sins of the World Behold verily verily I say unto you my servent Newel [Knight] you shall stand fast in the office wherewith I have appointed you \& if your Brethren desire to escape their enemies let them repent of all their sins \& become truly humble before me \& contrite \& as the covenant which they make unto me has been broken even so it hath become void \& of none affect \& wo to him by whom this offence cometh for it had been better for him that he had been drownded in the depth of the sea but blessed are they who have kept the covenant \& observed the commandment for they shall obtain mercy wherefore go to now \& flee the land lest your enemies come upon you And take your Journey \& appoint whom you will to be your leader \& to pay moneyes for you \& thus you shall take your Journeys into the regions westward unto Missorie unto the borders of the Lamanites \& after you have done Journeying Behold I say unto you seek ye a living like unto men untill I prepare a place for you \& again be patient in tribulation untill I come \& Behold I come quickly \& my reward is with me \& he that hath sought me early shall find rest to their Souls even So amen ------
				
		% -----
		\subsubsection{Revelation, 14 June 1831 [D\&C 55]}
		% -----
			\label{EMS-1831-JS-14-JUN-1831}
			% \label{EMS-1830-JS-DAY-3LETTERMONTH-YEAR}
		
			\begin{description}
				\item[Print Sources]
			\end{description}

				\begin{tabular}{p{0.5in}p{1in}p{0.5in}p{1in}}
				JSP	 	& D1:336--339	& JSR 		& 140 \\
				DC	 	& Section 55 	& HC 		& 1:185--186 \\
				HDDC 	& 709--715		& TCHC 		& 1:132	 \\
				\end{tabular}
				% HDDC = woodford's DC diss; JSR = Marquardt's JS Revelations; TCHC = Vogel HC
			
			\begin{description}
				\item[Online Access] \href{https://www.josephsmithpapers.org/paper-summary/revelation-14-june-1831-dc-55/1}{Here}
				% \item[Analysis] \hyperref[DC55]{Here}
			\end{description}
			
			\begin{description}
				\item[Background]
			\end{description}
			
				% It might also be useful to give a two sentence context of the period: e.g., "This speech was given in Nauvoo to a small congregation one month prior to the destruction of the Nauvoo Expositor. These and other tensions may have been on the prophet's mind when he gave the speech."
			
			\begin{description}
				\item[Original Source]
			\end{description}
			
				% brief description of original source material, with reference to JSP or somewhere else for more info
				% Background material: it might be helpful to have a note that says whether a given document was presented as a speech, was written in a journal, etc.
				% Also a comment here about how reliable the source might be, or what shape it is in, if there are large omissions or parts that are hard to understand, and so on.
			
			\begin{description}
				\item[Text]
			\end{description}
			
				\begin{tightcenter}
				58 Comandment June 14\supersu{th}\supers{.} 1831
				\end{tightcenter}
				
				A Revelation to William [W.] Phelps \& Joseph Coe their Calling \&c------
				
				Behold thus saith the lord unto you my servent William yea even the lord of the whole Earth thou art called \& chosen \& \sout{hast} after thou hast been baptized by water which if you do with an eye single to my glory you shall have a remission of your sins \& a reception of the Holy spirit by the laying on of hands \& then thou shalt be ordained by the hand of my servent Joseph to be an Elder unto this Church to Preach repentance \& remission of sins by way of baptism in the name of Jesus Christ the son of the living God \& on whomsoever you shall lay your hands \sout{\&} if they are contrite before me you shall have power t$\diamond$ give the holy spirit \& again thou shalt be ordained to assist my servent Oliver [Cowdery] to do the work of Printing \& of Selecting \& writing Books for Schools in this Church that little Children also may receive instruction before me as is pleasing unto me \& again verily I say unto you for this cause thou shalt take thy Journey with my servents Joseph \& Sidney [Rigdon] that thou mayest be planted in the land of thine inheritance to do this work \& again let my servent Joseph (Coe) also take his Journey with them the residue shall be made known hereafter even as I will Amen
				
		% -----
		\subsubsection{Revelation, 15 June 1831 [D\&C 56]}
		% -----
			\label{EMS-1831-JS-15-JUN-1831}
			% \label{EMS-1830-JS-DAY-3LETTERMONTH-YEAR}
		
			\begin{description}
				\item[Print Sources]
			\end{description}

				\begin{tabular}{p{0.5in}p{1in}p{0.5in}p{1in}}
				JSP	 	& D1:339--342	& JSR 		& 140--142 \\
				DC	 	& Section 56 	& HC 		& 1:186--188 \\
				HDDC 	& 716--725		& TCHC 		& 1:133--134 \\
				\end{tabular}
				% HDDC = woodford's DC diss; JSR = Marquardt's JS Revelations; TCHC = Vogel HC
			
			\begin{description}
				\item[Online Access] \href{https://www.josephsmithpapers.org/paper-summary/revelation-15-june-1831-dc-56/1}{Here}
				% \item[Analysis] \hyperref[DC56]{Here}
			\end{description}
			
			\begin{description}
				\item[Background]
			\end{description}
			
				% It might also be useful to give a two sentence context of the period: e.g., "This speech was given in Nauvoo to a small congregation one month prior to the destruction of the Nauvoo Expositor. These and other tensions may have been on the prophet's mind when he gave the speech."
			
			\begin{description}
				\item[Original Source]
			\end{description}
			
				% brief description of original source material, with reference to JSP or somewhere else for more info
				% Background material: it might be helpful to have a note that says whether a given document was presented as a speech, was written in a journal, etc.
				% Also a comment here about how reliable the source might be, or what shape it is in, if there are large omissions or parts that are hard to understand, and so on.
			
			\begin{description}
				\item[Text]
			\end{description}
			
				\begin{tightcenter}
				59\supersu{th}\supers{.} Commandment June 15\supersu{th}\supers{.} 1831
				\end{tightcenter}
				
				Thomas [B.] Marsh was desirous to know what he should do as the Lord had commanded him \& Ezra Thayer to take their Journey to the land of Missorie but Thayer could not get ready as \sout{he said} soon as Thomas wanted that he should \& these are the words of the Lord as follows
				
				Hearken O ye people which Profess my name saith the lord your God for Behold mine anger is kindeled against the rebelious \& they shall know mine arm \& mine indignation in the day of visitation \& of wrath upon the Nations \& he that will not take up his cross \& follow me \& keep my commandments the same shall not be saved behold I the Lord commandeth \& he that will not obey shall be cut off in mine own due time And after that I have commanded \& the commandment is broken Wherefore I the lord command \& revoke as it seemeth me good \& all this to be answered upon the heads of the rebelious saith the Lord Wherefore I revoke the commandment which was given unto my servent Thomas \& Ezra \& give a new commandment unto my Servent Thomas that he shall take up his Journey speedily to the land of Missorie \& my servent Sealy [Selah] Griffin shall also go with him for behold I revoke the commandment which was given to my servents Sealy \& Newel [Knight] in consequence of the stiffneckedness of my people which are in Thompson \& their rebelions Wherefore let my servent Newel remain with them \& as many as will go may go that are contrite before me \& be led by him to the land which I have appointed \& again verily I say unto you that my servent Ezra Thayer must repent of his pride \& his selfishness \& obey the former commandment which I have given him concerning the place upon which he lives \& if he will do this as there shall no divisions \sout{be} <be> made upon the land he shall be appointed still to go to the land of Missorie otherwise he shall receive the money which he has paid \& shall leave the place \& shall be cut off out of my Church saith the Lord god of host \& though the Heavens \& the Earth pass away these words shall not pass away but shall be fulfilled \& if my servent Joseph must needs pay the money Behold I the Lord will pay <it> unto him again in the land of Missorie that those of whom he shall receive may be rewarded again according to that which they do they shall receive even in lands for their in heritance Behold thus saith the lord unto my people you have many things to do \& to repent of for behold your sins have come up unto me \& are not pardoned because you seek to council in your own ways \& your hearts are not satisfied \& y$\diamond$ obey not the truth but have pleasure in unrighteousness wo! unto you rich men that will not give your substance to the poor for your riches will kanker your souls And this shall be your lamentation in the day of visitation \& of Judgement \& of indignation the Harvest is past the summer is ended \& my soul is not saved wo! unto you poor men whose hearts are not broken whose spirits are not contrite \& whose bellys are not satisfied \& whose hands are not stayed from laying hold upon other mens goods whose eyes are full of greediness who\sout{s} will not labour with their own hands But blessed are the poor <who are pure in heart whose hearts are broken \&> whose spirits are contrite for they shall see the Kingdom of God coming with power \& great glory unto their deliverance for the fatness of the Earth shall be theirs for Behold the lord shall come \& his recompence shall be with him \& he shall reward every man \& the poor shall rejoice \& their generations shall inherit the Earth from generation to generation for ever \& ever \& now I make an end of speaking unto you even So amen
				
		% -----
		\subsubsection{Note on Ordinations, circa 16 June 1831}
		% -----
			\label{EMS-1831-JS-CIR-16-JUN-1831}
			% \label{EMS-1830-JS-DAY-3LETTERMONTH-YEAR}
		
			\begin{description}
				\item[Print Sources]
			\end{description}

				\begin{tabular}{p{0.5in}p{1in}p{0.5in}p{1in}}
				JSP	 	& ---		& JSR 		& --- \\
				DC	 	& --- 		& HC 		& --- \\
				HDDC 	& ---		& TCHC 		& --- \\
				\end{tabular}
				% HDDC = woodford's DC diss; JSR = Marquardt's JS Revelations; TCHC = Vogel HC
			
			\begin{description}
				\item[Online Access] \href{https://www.josephsmithpapers.org/paper-summary/note-on-ordinations-circa-16-june-1831/1}{Here}
				% \item[Analysis] \hyperref[DC53]{Here}
			\end{description}
			
			\begin{description}
				\item[Background]
			\end{description}
			
				% It might also be useful to give a two sentence context of the period: e.g., "This speech was given in Nauvoo to a small congregation one month prior to the destruction of the Nauvoo Expositor. These and other tensions may have been on the prophet's mind when he gave the speech."
			
			\begin{description}
				\item[Original Source]
			\end{description}
			
				% brief description of original source material, with reference to JSP or somewhere else for more info
				% Background material: it might be helpful to have a note that says whether a given document was presented as a speech, was written in a journal, etc.
				% Also a comment here about how reliable the source might be, or what shape it is in, if there are large omissions or parts that are hard to understand, and so on.
			
			\begin{description}
				\item[Text]
			\end{description}
			
				\begin{tightcenter}
					June 6\supers{th.} 1831.
				\end{tightcenter}
			
				Brothers Simonds Rider Selah J. Griffin Daniel Stanton Peter Dustin Sidney Gilbert \& William W. Phelps were ordained Elders of this Church according to the church covenants under the hand of Joseph Smith jr.
				
		% -----
		\subsubsection{Revelation, 20 July 1831 [D\&C 57]}
		% -----
			\label{EMS-1831-JS-20-JUL-1831}
			% \label{EMS-1830-JS-DAY-3LETTERMONTH-YEAR}
		
			\begin{description}
				\item[Print Sources]
			\end{description}

				\begin{tabular}{p{0.5in}p{1in}p{0.5in}p{1in}}
				JSP	 	& D2:5--12		& JSR 		& 142--144 \\
				DC	 	& Section 57 	& HC 		& 1:189--190 \\
				HDDC 	& 726--737		& TCHC 		& 1:135 \\
				\end{tabular}
				% HDDC = woodford's DC diss; JSR = Marquardt's JS Revelations; TCHC = Vogel HC
			
			\begin{description}
				\item[Online Access] \href{https://www.josephsmithpapers.org/paper-summary/revelation-20-july-1831-dc-57/1}{Here}
				% \item[Analysis] \hyperref[DC57]{Here}
			\end{description}
			
			\begin{description}
				\item[Background]
			\end{description}
			
				% It might also be useful to give a two sentence context of the period: e.g., "This speech was given in Nauvoo to a small congregation one month prior to the destruction of the Nauvoo Expositor. These and other tensions may have been on the prophet's mind when he gave the speech."
			
			\begin{description}
				\item[Original Source]
			\end{description}
			
				% brief description of original source material, with reference to JSP or somewhere else for more info
				% Background material: it might be helpful to have a note that says whether a given document was presented as a speech, was written in a journal, etc.
				% Also a comment here about how reliable the source might be, or what shape it is in, if there are large omissions or parts that are hard to understand, and so on.
			
			\begin{description}
				\item[Text]
			\end{description}
			
				\begin{tightcenter}
				60 Commandment
				\end{tightcenter}
				
				Given in Missorie Independence Jackson Co July 20\supersu{th}\supers{.} 1831 giving directions <to the Bishop \& Agent> how to preceed concerning purchuseing Lands \&c. \&c.%
					\footnote{%
						% \label{}%
						JSP note: ``This heading likely did not appear in the original manuscript; John Whitmer likely added it when he copied the revelation into Revelation Book 1. It is not included in other manuscript copies of the revelation. At some point, Whitmer added `Not to be printed at present' to the copy in Revelation Book 1, and the revelation was not printed until the 1835 edition of the Doctrine and Covenants.'' (If you look at the image of this page of RB1, the ``[Not to be printd at present]'' is very prominent---it's large and in darker ink.)
						}
				
				Hearken Oh ye Elders of my Church, saith the Lord your God, Who have assembelled yourselves together, according to my commandment in this land which is the land of Missorie which is the Land which I, have appointed \& consecrated for the gethering of the Saints. Wherefore, this is the land of promise \& the place for the City of Zion. yea thus saith the Lord your God, If ye will receive wisdom here is wisdom. Behold the place which is now called Independence is the centre place, \& the spot for the Temple is lying westward upon a lot which is not far from the court-house. Wherefore it is wisdom that the land should be purchased by the saints \& also every tract lying westward even unto the line runing directly betwen Jew \& gentile And also every tract bordering by the Prairies in as much as my Deciples are enabled to buy lands. Behold this is wisdom that they may obtain it for an everlasting inheritance \& let my Servent Sidney Gilbert stand in the office which I have appointed to receive moneys to be an agent unto the church to buy lands in all the regions round about in as much as can be in righteousness, \& as wisdom shall direct. And let my servent Edward [Partridge] stand in the office which I have appointed him to [divide] unto the saints their inheritance even as I have commanded \& also them whom he has appointed to assist him And again verily I say unto you let my servent Sidney Gilbert plant himself in this place, \& establish a store that \sout{his} he may sell goods without fraud\sout{s} that he may obtain money to buy lands for the good\sout{s} of the Saints \& that he may Obtain provisions \& whatsoever things the Deciples may need to plant them in their inheritance \& also let my servent\sout{s} Sidney obtain license (behold here is wisdom \& whoso readeth let him understand) that he may send goods also unto the lamanites even by whom he will as clerks employed in his service \& thus the gospel may be preached unto them And again verily I say unto you let my servent William [W. Phelps] also be planted in this place \& be established as a Printer unto the Church \& lo. if the world receiveth his writings (behold this is wisdom) let him obtain whatsoever he can <obtain> in righteousness for the good of the saints. And let my servent\sout{s} Oliver [Cowdery] assist him even as I have commanded in Whatsoever place I shall appoint unto him to copy \& to correct \& select \&c that all things may be right before me as it shall be proved by the Spirit through him \& thus let those of whom I have spoken be planted in the Land of Zion as speedily as can be with their families to do these things even as I have spoken And concerning the gethering let the bishop \& the agent make preperations for th$\diamond$se families which have been commanded to come to this land as soon as posible \& plant them in their inheritance \& unto the residue of both Elders \& members further directions shall be given hereafter even So Amen------
				
		% -----
		\subsubsection{Revelation, 1 August 1831 [D\&C 58]}
		% -----
			\label{EMS-1831-JS-1-AUG-1831}
			% \label{EMS-1830-JS-DAY-3LETTERMONTH-YEAR}
		
			\begin{description}
				\item[Print Sources]
			\end{description}

				\begin{tabular}{p{0.5in}p{1in}p{0.5in}p{1in}}
				JSP	 	& D2:12--21		& JSR 		& 145--149 \\
				DC	 	& Section 58 	& HC 		& 1:191--195 \\
				HDDC 	& 738--751		& TCHC 		& 1:136--138 \\
				\end{tabular}
				% HDDC = woodford's DC diss; JSR = Marquardt's JS Revelations; TCHC = Vogel HC
			
			\begin{description}
				\item[Online Access] \href{https://www.josephsmithpapers.org/paper-summary/revelation-1-august-1831-dc-58/1}{Here}
				% \item[Analysis] \hyperref[DC58]{Here}
			\end{description}
			
			\begin{description}
				\item[Background]
			\end{description}
			
				% It might also be useful to give a two sentence context of the period: e.g., "This speech was given in Nauvoo to a small congregation one month prior to the destruction of the Nauvoo Expositor. These and other tensions may have been on the prophet's mind when he gave the speech."
			
			\begin{description}
				\item[Original Source]
			\end{description}
			
				% brief description of original source material, with reference to JSP or somewhere else for more info
				% Background material: it might be helpful to have a note that says whether a given document was presented as a speech, was written in a journal, etc.
				% Also a comment here about how reliable the source might be, or what shape it is in, if there are large omissions or parts that are hard to understand, and so on.
			
			\begin{description}
				\item[Text]
			\end{description}
			
				\begin{tightcenter}
				61 Commandment August 1\supersu{st}\supers{.} 1831
				\end{tightcenter}
				
				A Revelation given to the Elders who were assembeled \sout{in} on the land of Zion Directions what to do \&c \&c \&c
				
				Hearken oh ye Elders of my Church \& give ear to my word \& learn of me what I will concerning you \& also concerning this land unto which I have sent you for verily I say unto you blessed is he that keepeth my commandments whether in life or in death \& he that is faithfull in tribulation the reward of the same is greater in the kingdom of heaven ye cannot behold with your natural eyes for the present time the design of your God concerning those things which shall come hereafter \& the glory which shall follow after much tribulation for after much tribulation cometh the blessings Wherefore the day cometh that ye shall be crowned with much glory the hour is not yet but is nigh at hand remember this which I tell you before that you may lay it to heart \& receive that which shall follow behold verily I say unto you for this cause I have sent you that you might be obedient \& that your hearts might be prepared to bear testimony of the things which are to come \& also that you might be honoured of laying the foundation \& of bearing record of the land upon which the Zion of God shall stand \& also that a feast of fat things might be prepared for the poor yea a feast of fat things of wine on the lees well refined that the earth may know that the mouths of the Prophets shall not fail yea a supper of the house of the Lord well prepared unto which all nations shall be invited firstly the rich \& the learned the wise \& the Noble \& after that cometh the day of my Power then shall the poor the lame and the blind \& the deaf come in unto the marriage of the lamb \& partake of the supper of the Lord prepared for the great day to come Behold I the Lord have spoken it \& that the testimony might go forth from Zion yea from the mouth of the City of the heritage of God yea for this cause I have Sent you hither \& have Selected my Servent Edward [Partridge] \& appointed him his mission in this land but if he repent not of his sins which is unbelief \& blindness of heart let him take heed lest he fall behold his mission is given unto him \& it shall not be given again \& whoso standeth in that mission is appointed to be a Judge in Israel like as it was in ancient days to divide the lands of the heritage of God unto his children \& to Judge his people by the testimony of the Just \& by the assistance of his councillors according to the laws of the kingdom which are given by the Prophets of God for verily I say unto you my laws shall be kept on this land let no man think that he is ruler but let god rule him that Judgeth according to the council of his own will (or in other words) him that councileth or seteth upon the Judgement Seat let no man break the laws of the land for he that keepeth the laws of God hath no need to break the laws of the land Wherefore be subject to the powers that be untill he reigns whose right it is to reign \& subdues all enemies under his feet behold the laws which ye have received from my hand are the laws of the Church \& in this light ye shall hold them forth behold here is wisdom \& now as I spoke concerning my Servent Edward this land is the land of his residence \& those whom he has appointed for his councillors \& also the land of the residence of him whom I have appointed to keep my storehouse Wherefore let them bring their families to this land as they shall council between them \& me for behold it is not meet that I should command in all things for he that is compelled in all things the same is a slothfull \& not a wise Servent Wherefore he receiveth no reward verily I say men should be anxiously engaged in a good cause and do many things of their own free will, \& bring to pass much righteousness for the power is in them wherein they are agents unto themselves \& in-as-much as men are good they shall in no wise loose their reward but he that doeth not any thing untill he is commanded \& receiveth a commandment with a doubtfull heart \& keepeth it with slothfullness the same is damned Who am I that made man saith the Lord that will hold him guilty that obey not my commandments who am I saith the Lord that have promised \& have not fulfilled I command \& a man obeys not I revoke \& they receive not the blessing then they say in their hearts this is not the work of the Lord for his promises are not fulfilled but wo unto such for their reward lurketh beneath \& not from above \& now I give unto you further directions concerning this Land it is wisdom in me that my servent Martin [Harris] should be an example unto the church in laying his money before the bishop of the Church \& also this is a law unto every man that cometh unto the Land to receive an inheritance and he shall do with his moneys according as the law directs \& it is wisdom also that \sout{it} there should be lands purchased in Independence for the place of the storehouse \& also for the house of the Printing \& other directions concerning my servent Martin shall be given him of the spirit that he may receive his inheritance as seemeth him good \& let him repent of his sins for he seeketh the praise of the world \& also let my servent William [W. Phelps] stand in the office which I have appointed him \& receive his inheritance in the Land \& also he hath need to repent for I the lord am not pleased with him for he seeketh to excell \& he is not sufficiently \sout{humble} meek in his heart behold he that hath repented of his sins the same is forgiven \& I the Lord remembereth them no more by this ye may know if a man repenteth of his sins behold he will confess them \& forsake them \& now verily I say concerning the residue of the Elders of my Church the time has not yet come for many years for them to receive their inheritance in this land except they desire it through prayer only as \sout{they} it shall be appointed unto them for Behold they shall push the people together from the ends of the Earth wherefore assemble yourselves together \& he that is not appointed to stay in this land let them preach the gospel in the regions round about \& after that let them return to their homes let them preach by the way \& bear testimony of the truth in all places \& call upon the rich the high \& the low \& the poor to repent \& let them build up churches in-as-much as the inhabitants of the Earth will repent \& let there be an agent appointed by the voice of the Church \& I give unto my servent Sidney [Rigdon] a commandment that he shall write a discription of the Land of Zion \& a statement of the will of God as it shall be made known by the spirit unto him and an Epistle \& subscription to be presented unto all the Churches to obtain money to be put into the hands of the Bishop to purchase lands for an inheritance for the children of God of himself or the agent as seemeth him good or as he shall direct for behold verily I say unto you the Lord willeth that the deciples \& the children of men should open their hearts even to purchase this whole region of country as soon as time will permit behold here is wisdom let them do this lest they reserve none inheritance save it be by the sheding of blood \& again in as much as there is lands obtained let there be workmen sent forth of all kinds unto this land to labour for the saints of God let all these things be done in order \& let the priveliges of the lands be made known from time to time by the Bishop or the agent of the Church \& let the work of the gethering be not in haste nor by flight but let it be done as it shall be councelied by the Elders of the Church at the conferences according to the knoweledge which they receive from time to time \& let my servent Sidney consecrate \& dedicate this land \& the spot of the temple unto the Lord \& let a conference meeting be called \& after that let my servent Sidney \& Joseph return \& also Oliver [Cowdery] with them to accomplish the residue of the work which I have appointed unto them in their own land \& the residue as shall be ruled by the conference \& let no man return from this land except he bear record by the way of that which he knows \& most assuredly believes let that which has been bestowed upon Ziba [Peterson] be taken from him \& let him stand as a member in the Church \& labour with his own hands with the brethren untill he is sufficiently chastened for all his sins for he confeseth them not \& he thinketh to hide them let the residue of the Elders of this church which are coming to this land some of whom are exceedingly blessed even above measure also hold a conference which shall be held by them \& let them also return preaching the gospel by the way bearing record of the things which are revealed unto them for verily the sound must go forth from this place into all the world \& unto the uttermost parts of the Earth the gospel must be preached unto every creature with signs following them that believe \& behold the son of man cometh Amen
				
		% -----
		\subsubsection{Minutes, 4 August 1831}
		% -----
			\label{EMS-1831-JS-4-AUG-1831}
			% \label{EMS-1830-JS-DAY-3LETTERMONTH-YEAR}
		
			\begin{description}
				\item[Print Sources]
			\end{description}

				\begin{tabular}{p{0.5in}p{1in}p{0.5in}p{1in}}
				JSP	 	& D2:22--24		& JSR 		& --- \\
				DC	 	& --- 			& HC 		& 1:199 \\
				HDDC 	& ---			& TCHC 		& 1:142 \\
				\end{tabular}
				% HDDC = woodford's DC diss; JSR = Marquardt's JS Revelations; TCHC = Vogel HC
			
			\begin{description}
				\item[Online Access] \href{https://www.josephsmithpapers.org/paper-summary/minutes-4-august-1831/1}{Here}
				% \item[Analysis] \hyperref[XX]{Here}
			\end{description}
			
			\begin{description}
				\item[Background]
			\end{description}
			
				% It might also be useful to give a two sentence context of the period: e.g., "This speech was given in Nauvoo to a small congregation one month prior to the destruction of the Nauvoo Expositor. These and other tensions may have been on the prophet's mind when he gave the speech."
			
			\begin{description}
				\item[Original Source]
			\end{description}
			
				% brief description of original source material, with reference to JSP or somewhere else for more info
				% Background material: it might be helpful to have a note that says whether a given document was presented as a speech, was written in a journal, etc.
				% Also a comment here about how reliable the source might be, or what shape it is in, if there are large omissions or parts that are hard to understand, and so on.
			
			\begin{description}
				\item[Text]
			\end{description}
			
				...E[x]hortation by br. Joseph Smith jr. to acts of righteousness and keeping the commandments of the Lord with promise of blessing<s>...
				
		% -----
		\subsubsection{Revelation, 7 August 1831 [D\&C 59]}
		% -----
			\label{EMS-1831-JS-7-AUG-1831}
			% \label{EMS-1830-JS-DAY-3LETTERMONTH-YEAR}
		
			\begin{description}
				\item[Print Sources]
			\end{description}

				\begin{tabular}{p{0.5in}p{1in}p{0.5in}p{1in}}
				JSP	 	& D2:30--35		& JSR 		& 149--150 \\
				DC	 	& Section 59	& HC 		& 1:200--201 \\
				HDDC 	& 752--763		& TCHC 		& 1:142--143 \\
				\end{tabular}
				% HDDC = woodford's DC diss; JSR = Marquardt's JS Revelations; TCHC = Vogel HC
			
			\begin{description}
				\item[Online Access] \href{https://www.josephsmithpapers.org/paper-summary/revelation-7-august-1831-dc-59/1}{Here}
				% \item[Analysis] \hyperref[DC59]{Here}
			\end{description}
			
			\begin{description}
				\item[Background]
			\end{description}
			
				% It might also be useful to give a two sentence context of the period: e.g., "This speech was given in Nauvoo to a small congregation one month prior to the destruction of the Nauvoo Expositor. These and other tensions may have been on the prophet's mind when he gave the speech."
			
			\begin{description}
				\item[Original Source]
			\end{description}
			
				% brief description of original source material, with reference to JSP or somewhere else for more info
				% Background material: it might be helpful to have a note that says whether a given document was presented as a speech, was written in a journal, etc.
				% Also a comment here about how reliable the source might be, or what shape it is in, if there are large omissions or parts that are hard to understand, and so on.
			
			\begin{description}
				\item[Text]
			\end{description}
			
				Behold blessed saith the Lord are they who have come up unto this land with an eye single to my glory according to my Commandments for them that live shall inherit the earth and them that die shall rest from all their labours \& their works shall follow them they shall receive a crown in the mansions of my Father which I have prepared for them. Yea blessed are they whose feet stand upon the land of Zion who have obeyed my Gospel for they shall receive for their reward the good things of the earth \& it shall bring forth in her strength \& they also shall be crowned with blessings from above yea \& with commandments not a few \& with revelations in their time they that are faithful \& diligent before me. Wherefore I give unto them a commandment saying thus Thou shalt love the Lord thy God with all thy heart with all thy might mind \& strength \& in the name of Jesus Christ thou shalt serve him thou shalt love thy neighbour as thyself thou shalt not steal neither commit adultry nor kill or do any thing like unto it thou shalt thank the Lord thy God in all things thou shalt offer a sacrafice unto the Lord thy God in righteousness even that of a broken heart \& a contrite spirit. And that thou mayest more fully keep thyself unspotted from the world thou shalt go to the house of prayer \& offer up thy sacraments upon my holy day for verily this is a day appointed unto you to rest from your labours \& to pay thy devotions unto the most high Nevertheless thy vows should be offered up in righteousness <in> all days \& at all times but remember that on this the Lords day thou shalt offer thine oblations \& thy sacraments unto the most High Confessing thy sins unto thy brethren \& before the Lord \& on this day thou shalt do none other <things> only let thy food be prepared with singleness of heart that thy fastings may be perfect or in other words that thy joy may be full verily this is fasting and prayer or in other words rejoicing \& prayer. And inasmuch as ye do these things with thanksgiving with cheerful hearts \& coutenances not with much laughter (for this is sin) but with a glad heart \& a cheerful countenance verily I say that inasmuch as ye do this the fulness of the earth is yours the beasts of the fields \& the fowls of the air \& that which climbeth upon trees \& walketh upon the earth yea \& the herb \& the good things which cometh of the earth whether for food or for raiment or for houses or for barns or for orchards or for gardens or for vineyards yea all things which cometh of the earth in the season therof is made for the benefit \& the <use> of man both to please the eye \& to gladen the heart yea for food \& for raiment for taste \& for smell to strengthen the body \& to enliven the soul \& it pleaseth God that he hath given all these things unto man for unto this end were they made to be used with judgement not to excess neither by extortion \& in nothing doth man offend God or against none is his wrath kindled save those who Confess not his hand in all things \& obey not his commandments behold this is according to the law \& the prophets. Wherefore trouble me no more concerning this matter but learn that he who doeth the works of righteousness shall receive his reward even peace in this world \& eternal life in the world to come I the Lord hath spoken it \& the spirit beareth record Amen
				
				Given by Joseph the translatior \& written by Oliver [Cowdery] August 7. 1831 in the land of Zion%
					\footnote{%
						% \label{}%
						HDDC 753 says that there is a version of this revelation in the handwriting of Samuem Harrison Smith, and at the end it says: ``Given by Joseph the translator \& written by Oliver August 7, 1831 in the land of Zion \& copied by Samuel H. Smith Brother to the Seer.'' But I can't find this mentioned on the JSP website.
						}
				
		% -----
		\subsubsection{Revelation, 8 August 1831 [D\&C 60]}
		% -----
			\label{EMS-1831-JS-8-AUG-1831}
			% \label{EMS-1830-JS-DAY-3LETTERMONTH-YEAR}
		
			\begin{description}
				\item[Print Sources]
			\end{description}

				\begin{tabular}{p{0.5in}p{1in}p{0.5in}p{1in}}
				JSP	 	& D2:35--37		& JSR 		& 150--151 \\
				DC	 	& Section 60 	& HC 		& 1:201--202 \\
				HDDC 	& 764--770		& TCHC 		& 1:143 \\
				\end{tabular}
				% HDDC = woodford's DC diss; JSR = Marquardt's JS Revelations; TCHC = Vogel HC
			
			\begin{description}
				\item[Online Access] \href{https://www.josephsmithpapers.org/paper-summary/revelation-8-august-1831-dc-60/1}{Here}
				% \item[Analysis] \hyperref[DC60]{Here}
			\end{description}
			
			\begin{description}
				\item[Background]
			\end{description}
			
				% It might also be useful to give a two sentence context of the period: e.g., "This speech was given in Nauvoo to a small congregation one month prior to the destruction of the Nauvoo Expositor. These and other tensions may have been on the prophet's mind when he gave the speech."
			
			\begin{description}
				\item[Original Source]
			\end{description}
			
				% brief description of original source material, with reference to JSP or somewhere else for more info
				% Background material: it might be helpful to have a note that says whether a given document was presented as a speech, was written in a journal, etc.
				% Also a comment here about how reliable the source might be, or what shape it is in, if there are large omissions or parts that are hard to understand, and so on.
			
			\begin{description}
				\item[Text]
			\end{description}
			
				\begin{tightcenter}
				63 Commandment
				\end{tightcenter}
				
				given in Missorie Jackson County Independence August 8\supersu{th}{\supers{.} 1831 directions to some of the Elders to return to their \sout{homes} \sout{\& \&} own land \&c \&c
				
				Behold thus saith the Lord unto the Elders of this Church who are to return speedily to the land from whence they came behold it pleaseth me that you have come up hither but with some I am not well pleased for they will not open their mouths but hide the tallent which I have given <unto> them because of the fear of man wo unto such for mine anger is kindelled against them \& it shall come to pass if they are not more faithfull unto me it shall be taken away even that which they have for I the Lord ruleth in the heavens above \& among the armies of the Earth And in the day when I shall make up my Jewels all men shall know what it is that bespeaketh the power of God but verily I will speak unto you concerning your Journey unto the Land from whence you came let there be a craft made or bought as seemeth you good it mattereth not unto me \& take your Journey speedily for the place which is called St. Lewis. \& from thence let my Servent Sidney [Rigdon] \& Joseph \& Oliver [Cowdery] take their Journey for Cincinnati \& in this place let them lift up their voice \& declare my word with loud voices without wrath or doubting lifting up holy hands upon them for I am able to make you holy \& your sins are forgiven you \& let the residue take their Journey from St. Lowis two by two \& preach the word not in haste among the congregations of the wicked untill they return to the churches from whence they came \& all this for the good of the churches for this intent have I sent them \& let my servent Edward [Partridge] impart of the money which I have given him a portion unto mine Elders which are commanded to return \& he that is able let him return it by the way of the agent \& he that is not of him it is not required And now I speak of the residue which is to come to this Land Behold they have been sent to preach my gospel among the congregations of the wicked wherefore I give unto them a commandment Thus thou shalt not Idle away thy time neither shalt thou bury thy tallent that it may not be known \& after thou hast come up unto the land of Zion \& have proclaimed my word thou shalt speedily return preaching the word among the congregations not in haste neither in wrath nor with strife \& shake off the dust of thy feet against those who receive thee not not in their presence lest thou provoke them but in secret \& wash thy feet as a testimony against them in the day of Judgement this is sufficient for you \& the will of him who hath sent you \& by the mouth of my servent Joseph it shall be made known concerning sidney \& Oliver the residue hereafter even so amen------
				
		% -----
		\subsubsection{Revelation, 12 August 1831 [D\&C 61]}
		% -----
			\label{EMS-1831-JS-12-AUG-1831}
			% \label{EMS-1830-JS-DAY-3LETTERMONTH-YEAR}
		
			\begin{description}
				\item[Print Sources]
			\end{description}

				\begin{tabular}{p{0.5in}p{1in}p{0.5in}p{1in}}
				JSP	 	& D2:37--44		& JSR 		& 152--154 \\
				DC	 	& Section 61 	& HC 		& 1:203--205 \\
				HDDC 	& 771--782		& TCHC 		& 1:144--145 \\
				\end{tabular}
				% HDDC = woodford's DC diss; JSR = Marquardt's JS Revelations; TCHC = Vogel HC
			
			\begin{description}
				\item[Online Access] \href{https://www.josephsmithpapers.org/paper-summary/revelation-12-august-1831-dc-61/1}{Here}
				% \item[Analysis] \hyperref[DC61]{Here}
			\end{description}
			
			\begin{description}
				\item[Background]
			\end{description}
			
				% It might also be useful to give a two sentence context of the period: e.g., "This speech was given in Nauvoo to a small congregation one month prior to the destruction of the Nauvoo Expositor. These and other tensions may have been on the prophet's mind when he gave the speech."
			
			\begin{description}
				\item[Original Source]
			\end{description}
			
				% brief description of original source material, with reference to JSP or somewhere else for more info
				% Background material: it might be helpful to have a note that says whether a given document was presented as a speech, was written in a journal, etc.
				% Also a comment here about how reliable the source might be, or what shape it is in, if there are large omissions or parts that are hard to understand, and so on.
			
			\begin{description}
				\item[Text]
			\end{description}
			
				\begin{tightcenter}
				64 Commandment
				\end{tightcenter}
				
				given Aug 12\supersu{th}\supers{.} 1831 on the Bank of the River Distruction (or Missorie) unfolding some mysteries \&c \&c------
				
				Behold \& hearken unto \sout{him} the voice of him who hath all power who is from everlasting to everlasting even alpha \& omega the begining \& the end Behold verily thuss saith the lord unto you O ye Elders of my Church who are assembled upon this spot whose sins are now forgiven you for I the Lord forgiveth sins \& am mercyfull unto those who confess their sins with humble hearts but verily I say unto you that it is not needfull for this whole company of mine Elders to be moveing swiftly upon the waters whilst the Inhabitants on either sides are perishing in unbelief nevertheless I suffered it that ye might bear record Behold there are many dangers upon the waters \& more especially hereafter for I the Lord have decreed in mine anger many distructions upon the waters yea \& especially upon these waters nevertheless all flesh is in mine hand he that is faithfull among you shall not perish by the waters wherefore it is expedient that my servent Sidney Gilbert \& <my servent> William [W.] Phelps be in haste upon their errand \& mission nevertheless I would not suffer that ye should part untill \sout{ye} you are chastened for all your sins that you might be one that you might not perish in wickedness but now verily I say it Behooveth me that ye should part wherefore let them my servent Sidney [Gilbert] \& William take their former company \& let them take their Journey in haste that they may fill their mission \& through faith they shall overcome \& in as much as they are faithfull they shall be preserved \& I the Lord will be with them \& let the residue take that which is needfull for clothing let my servent sidney [Gilbert] take that which is not needfull with \sout{them} him as you shall agree \& now behold for your good I give unto you a commandment concerning these things \& I the Lord will reason with you as with men in days of old Behold I the Lord in the begining belessed the waters but in the last days by the mouth of my servent John I cursed the waters wherefore the days will come that no flesh shall be safe upon the waters \& it shall be said in days to come that none is able to go up to the land of Zion upon the waters but he that is upright in heart \& as I the Lord in the begining cursed the land even so in the last days have I blessed it in its time for the use of my saints that they may partake the fatness thereof \& now I give unto you a commandment \& what I say unto one I say unto all that you shall forewarn your brethren concerning these waters that they come not in Journeying \sout{on} upon them lest their faith fail \& they are caught in her snares I the Lord hath decreed \& the destroyer rideth upon the face thereof \& I revoke not the decree I the Lord was angery with you yesterday but to day mine anger is turned away wherefore let those whom I have spoken that should take their Journey in haste again I say unto you let them take their Journey in haste \& it mattereth not unto me after a little if it so be that they fill their mission whether they go by water or by land let this be as it is made known unto them according to their Judgement \& now concerning my servents Sidney [Rigdon] Joseph \& Oliver [Cowdery] let them come not again upon the waters save it be upon the canal while Journeying unto their homes or in other words they shall not come upon upon the waters <to Journey> save upon the canal Behold I the Lord have appointed a way for the Journeying of my saints \& behold this is the way that after they leave the canal they shall Journey by land in as much as they are commanded to Journey \sout{by} \& go up unto the land of Zion \& they shall do like unto the children of Israel pitching their tents by the way \& behold this commandment you shall give unto all your brethren nevertheless unto whom it is given power to command the waters unto him it is given by the spirit to know all his ways wherefore let him do as the spirit of the living God commandeth him whether upon the land or upon the waters as it remaineth with me to do hereafter \& unto you it is given the course of the saints or the way for the saints of the camp of the Lord to Journey \& again verily I say unto you my Servents Sidney [Rigdon] Joseph \& Oliver shall not open their mouths in the congregations of the wicked untill they arrive at cincinnati \& in that place they shall lift up their voices unto god against that People yea unto him whose anger is kindelled against their wickedness a people which is well ripened for distruction \& from thence let them Journy for the congregations of their brethren for their labours even now are wanted more abundantly among them then among the congregations of the wicked \& now concerning the residue let them Journey \& declare the word among the congregations of the wicked inasmuch as it is given \& in as much as they do this they shall rid their garments \& they shall be spotless before me \& let them Journey together or two by two as seemeth them good only let my servent reynolds [Cahoon] \& my Servent Samuel [Smith] with whom I am well pleased be not seperated untill they return to their homes \& this for a wise purpose in me \& now verily I say unto you \& what I say unto one I say unto all be of good cheer little children for I am in your midst \& I have not forsaken you \& in as much as ye have humbelled yourselves before me the blessings of the kingdom is yours gird up your loins \& be watchfull \& be sober looking forth for the coming of the Son of man in an hour you think not pray always that you enter not into temptation that you may abide the day of his coming whether in life or in death even So Amen
				
		% -----
		\subsubsection{Revelation, 13 August 1831 [D\&C 62]}
		% -----
			\label{EMS-1831-JS-13-AUG-1831}
			% \label{EMS-1830-JS-DAY-3LETTERMONTH-YEAR}
		
			\begin{description}
				\item[Print Sources]
			\end{description}

				\begin{tabular}{p{0.5in}p{1in}p{0.5in}p{1in}}
				JSP	 	& D2:44--48		& JSR 		& 154--155 \\
				DC	 	& Section 62 	& HC 		& 1:206 \\
				HDDC 	& 783--790		& TCHC 		& 1:145--146 \\
				\end{tabular}
				% HDDC = woodford's DC diss; JSR = Marquardt's JS Revelations; TCHC = Vogel HC
			
			\begin{description}
				\item[Online Access] \href{https://www.josephsmithpapers.org/paper-summary/revelation-13-august-1831-dc-62/1}{Here}
				% \item[Analysis] \hyperref[DC62]{Here}
			\end{description}
			
			\begin{description}
				\item[Background]
			\end{description}
			
				% It might also be useful to give a two sentence context of the period: e.g., "This speech was given in Nauvoo to a small congregation one month prior to the destruction of the Nauvoo Expositor. These and other tensions may have been on the prophet's mind when he gave the speech."
			
			\begin{description}
				\item[Original Source]
			\end{description}
			
				% brief description of original source material, with reference to JSP or somewhere else for more info
				% Background material: it might be helpful to have a note that says whether a given document was presented as a speech, was written in a journal, etc.
				% Also a comment here about how reliable the source might be, or what shape it is in, if there are large omissions or parts that are hard to understand, and so on.
			
			\begin{description}
				\item[Text]
			\end{description}
			
				\begin{tightcenter}
				65 Commandment
				\end{tightcenter}
				
				given Aug 13\supersu{th}\supers{.} 1831 on the Bank of the river Missorie at a meeting of some of the Elders which had not yet arived at their Journeys end \&c
				
				Behold \& hearken oh ye Elders of my Church saith the Lord your God even Jesus Christ your advocate who knoweth the weakness of man \& how to sucour they that are tempted \& verily mine eyes are upon \sout{you} those who have not as yet gone up unto the Land of Zion wherefore your mission is not yet full nevertheless ye are blessed for the testimony which ye have borne is recorded in heaven for the Angels to look upon \& they rejoice over you \& your sins are forgiven you \& now continue your Journey assemble yourselves upon the land of Zion \& hold a meeting \& rejoice together \& offer a sacrament unto the most high \& then you may return to bear record yea even all together or two by two as seemeth you good it mattereth not unto me only be faithfull \& declare glad tidings unto the inhabitants of the Earth or among the Congregations of the wicked Behold I the Lord have brought you together that the promise might be fulfilled that the faithfull among you should be preserved \& rejoice together in the Land of Missorie I the Lord promised the faithfull, \& cannot lie I the Lord am willing if any among you desireth to ride upon horses or upon mules or in chariots shall receive this blessing if he receive it from the hand of the Lord with <a> thankfull\sout{ness} hearts in all things these things remain with you to do according to Judgement \& the directions of the spirit Behold the kingdom is yours And Behold \& lo I am with the faithfull always even so Amen------
				
		% -----
		\subsubsection{Revelation, 30 August 1831 [D\&C 63]}
		% -----
			\label{EMS-1831-JS-30-AUG-1831}
			% \label{EMS-1830-JS-DAY-3LETTERMONTH-YEAR}
		
			\begin{description}
				\item[Print Sources]
			\end{description}

				\begin{tabular}{p{0.5in}p{1in}p{0.5in}p{1in}}
				JSP	 	& D2:48--55		& JSR 		& 155--158 \\
				DC	 	& Section 63 	& HC 		& 1:207--211 \\
				HDDC 	& 791--810		& TCHC 		& 1:146--149 \\
				\end{tabular}
				% HDDC = woodford's DC diss; JSR = Marquardt's JS Revelations; TCHC = Vogel HC
			
			\begin{description}
				\item[Online Access] \href{https://www.josephsmithpapers.org/paper-summary/revelation-30-august-1831-dc-63/1}{Here}
				% \item[Analysis] \hyperref[DC63]{Here}
			\end{description}
			
			\begin{description}
				\item[Background]
			\end{description}
			
				% It might also be useful to give a two sentence context of the period: e.g., "This speech was given in Nauvoo to a small congregation one month prior to the destruction of the Nauvoo Expositor. These and other tensions may have been on the prophet's mind when he gave the speech."
			
			\begin{description}
				\item[Original Source]
			\end{description}
			
				% brief description of original source material, with reference to JSP or somewhere else for more info
				% Background material: it might be helpful to have a note that says whether a given document was presented as a speech, was written in a journal, etc.
				% Also a comment here about how reliable the source might be, or what shape it is in, if there are large omissions or parts that are hard to understand, and so on.
			
			\begin{description}
				\item[Text]
			\end{description}
			
				Hearken O ye people and open your hearts and give ear from afar and listen you that call yourselves the people of the Lord \& hear the word of the Lord \& his will concerning you yea verily I say hear the word of him whose anger is kindled against \sout{you} <the> wicked \& rebellious who willeth to take even them whom he will take \& preserveth in life them whom he will preserve who buildeth up at his own will \& pleasure \& destroyeth when he please \& is able to cast the soul down to hell. Behold I the Lord uttereth my voice \& <\&> it shall be obeyed wherefore verily I say let the wicked take heed \& let the rebellious fear \& tremble \& let the unbelieving hold their lips for the day of wrath shall come upon them as a whirlwind \& all flesh shall know that I am God. And he that seeketh signs shall see signs but not unto salvation Verily I say unto you There are those among you who seeketh signs \& there has been such even from the beginning But behold faith cometh not by signs but signs follow those that believe yea signs cometh by faith not by the will of men nor as they please but by the will of God yea signs cometh by faith unto mighty works for without faith no man pleaseth God \& with whom God is angry he is not well pleased wherefore unto such he sheweth no signs only in wrath unto their condemnation Wherefore I the Lord am not \sout{well} pleased with those among you who have sought after signs \& wonders for faith \& not for the good of men unto my glory Nevertheless I gave commandments \& many have turned away from my commandments \& have not kept them There were among you adulterers \& adulteresses some of whom have turned away from you \& others remain with you that hereafter shall be revealed let such be aware \& repent speedily lest judgements shall come upon them as a snare \& their folly shall be made manifest \& their works shall follow them in the eyes of the people \& verily I say unto you as I have said before he that looketh on a woman to lust after her or if any shall commit adultry in his heart they shall not have the spirit but shall deny the faith \& shall fear Wherefore I the Lord have said that the fearful \& the unbelieving \& all liars \& whoso<ever> loveth \& maketh a lie \& the whoarmunger <\& the> sorcerer should have their part in that lake which burneth with fire \& brimstone which is the second death Verily I say <that> they shall not have part in the first resurrection And now behold I the Lord saith unto you that ye are not justified because \sout{of} these things are among you nevertheless he that endureth in faith \& doeth my will the same shall overcome \& shall receive an inheritance upon the Earth when the day of transfiguration shall come when the earth shall be transfigured even according to the pattern which was shown unto mine apostles upon the mount of which account the fulness ye have not yet received. And now verily I say unto you that as I said that I would make known my will unto you behold I will make it known unto you not by the way of commandment for their are many who observe not to keep my commandments but unto him that keepeth my commandments I will give the mysteries of my Kingdom \& the same shall be in him a well of <living> water springing up unto everlasting life And now behold this is the will of the Lord your God concerning his saints that they should assemble themselves together unto the land of Zion not in haste lest there should be confusion which bringeth pestilence Behold the land of Zion I the Lord holdeth it in mine own hands nevertheless I the Lord rendereth unto Cezar the things which are Cezars Wherefore I the Lord willeth that you should purchase the lands that you may have advantage of the world that you may have claim on the world that they may not be stired up unto anger for satan putteth it into their hearts to anger \& to the shedding of blood Wherefore the land of Zion shall not be obtained but by purchase or by blood otherwise there is none inheritance for you \& if by purchase\sout{d} behold \sout{ye} you are blessed \& if by blood as ye are forbidden to shed blood lo your enemies are upon you \& ye shall be scourged from city to city \& from Synagogue to synagogue \& but few shall stand to receive an inheritance I the Lord am angry with the wicked I am holding my spirit from the inhabitants of the earth I have sworn in my wrath \& decreed wars upon the face of the earth \& the wicked shall slay the wicked \& fear shall come upon every man \& the Saints also shall hardly escape nevertheless I the Lord am with them \& will come down in Heaven from the presence of God \& consume the wicked with unquenchable fire \& behold this is not yet but by \& by Wherefore seeing that I the Lord have decreed all these things upon the face of the earth I willeth that my saints should be assembled upon the land of Zion \& that every man should take righteousness in his hands \& faithfulness upon his loins \& lift a warning voice unto the inhabitants of the earth \& declare both by word \& by flight that desolation shall come upon the wicked Wherefore let my Desiples in Kirtland arrange their temporal concerns which dwell upon this farm let my servant Titus [Billings] who has the care thereof dispose of the land that he may be prepared in the coming spring to take his journey up unto the land of Zion with those that dwell upon the face thereof excepting those whom I shall reserve unto myself that shall not go until I shall command them \& let all the moneys which can be spared (it mattereth not unto me whether it be little or much) sent up unto the land of Zion unto them whom I have appointed to receive. Behold I the Lord will give unto my servant Joseph power that he shall be enabled to descern by the spirit those who shall go up unto the land of Zion \& those of my Desiples that shall tarry Let my servant Newel [K.] Whitney retain \sout{the} his store or in other words the store yet for a little season nevertheless let him impart all the money which he can impart to be sent up unto the land of Zion behold these things are in his own hands let him do according to wisdom verily I say let him be ordained as an agent unto the Desiples that shall tarry \& let him be ordained unto this power \& now speedily \sout{go with} <visit> the churches expounding these things unto them with my servant Oliver [Cowdery] behold this is my <will> \sout{with} obtaining moneys even as I have directed he that is faithful \& endureth shall overcome the world he that sendeth up treasures unto the land of Zion shall receive an inheritance in this world \& his works shall follow him \& also a reward in the world to come yea \& blessed are the dead that die in the Lord from henceforth when the Lord shall come \& old things shall pass away \& all things become new they shall rise from the dead \& shall not die \& shall receive an inheritance before the Lord in the Holy City \& he that liveth when the Lord shall come \& have kept the faith blessed is he nevertheless it is appointed unto him to die at the age of man Wherefore children shall grow up until they become old. Old men shall die but they shall not sleep in the dust but they shall be changed in the twinkling of an eye Wherefore for this cause preached the apostles unto the world the resurrection of the dead these things are the things that ye must look for \& speeking after the manner of the Lord they are now nigh at hand \& in a time to come even in the day of the coming of the Son of man \& until that hour there will be foolish virgins among the wise \sout{\&} \& at that hour cometh an entire separation of the righteous \& the wicked \& in that day will I send mine angels \& pluck out the wicked \& cast into unquenchable fire And now <behold> verily I say unto you I the Lord am not pleased with my servant Sidney [Rigdon] he exaulted himself in his heart \& received not counsel but grieved the spirit Wherefore his writing is not acceptable unto the Lord \& he shall make another \& if the Lord receive it not behold he standeth no longer in the office which he hath appointed him And again verily I say unto you let those who desire in their hearts in meekness to warn sinners to repentance let them be ordained unto this power for this is a day of warning \& not a day of many \sout{worlds} words for I the Lord am not to be mocked in the last days Behold I am from above \& my power lieth beneath I am over all \& in all \& through all \& searcheth all things \& the days cometh that all things shall be subject unto me. Behold I am Alpha \& Omega even Jesus Christ Wherefore let all men be ware how they take my name in their lips for behold verily I say that many there be who are under this condemnation who useth the name of the Lord \& useth it in vain having not authority Wherefore let the church repent of their sins \& I the Lord \sout{with} <will> own them otherwise they shall be cut off. Remember that that \sout{cometh} which cometh from above is sacred \& must be spoken with care \& by constraint of the spirit \& in this there is no condemnation and ye receive the spirit through prayer wherefore without this there remaineth condemnation. Let my servants Joseph \& Sidney seek them a home as they are \sout{led} taught through \sout{the spirit} prayer by the spirit these things remain to overcome through patience that such may receive a more exceeding \& eternal weight of glory otherwise a greater condemnation Amen.
				
				Given by Joseph the Seer in Kirtland August 31, 1831 and written by Oliver------
				
		% -----
		\subsubsection{Revelation, 31 August 1831}
		% -----
			\label{EMS-1831-JS-31-AUG-1831}
			% \label{EMS-1830-JS-DAY-3LETTERMONTH-YEAR}
		
			\begin{description}
				\item[Print Sources]
			\end{description}

				\begin{tabular}{p{0.5in}p{1in}p{0.5in}p{1in}}
				JSP	 	& D2:55--56		& JSR 		& --- \\
				DC	 	& ---		 	& HC 		& --- \\
				HDDC 	& ---			& TCHC 		& --- \\
				\end{tabular}
				% HDDC = woodford's DC diss; JSR = Marquardt's JS Revelations; TCHC = Vogel HC
			
			\begin{description}
				\item[Online Access] \href{https://www.josephsmithpapers.org/paper-summary/revelation-31-august-1831/1}{Here}
				% \item[Analysis] \hyperref[XX]{Here}
			\end{description}
			
			\begin{description}
				\item[Background]
			\end{description}
			
				% It might also be useful to give a two sentence context of the period: e.g., "This speech was given in Nauvoo to a small congregation one month prior to the destruction of the Nauvoo Expositor. These and other tensions may have been on the prophet's mind when he gave the speech."
			
			\begin{description}
				\item[Original Source]
			\end{description}
			
				% brief description of original source material, with reference to JSP or somewhere else for more info
				% Background material: it might be helpful to have a note that says whether a given document was presented as a speech, was written in a journal, etc.
				% Also a comment here about how reliable the source might be, or what shape it is in, if there are large omissions or parts that are hard to understand, and so on.
			
			\begin{description}
				\item[Text]
			\end{description}
			
				Behold thus saith the Lord by the voice of the Spirit it is wisdom in me that my Servent John Burk David Eliot [Elliott] Erastus Babit [Babbitt] should take their Journey this fall to the land of Zion
				
				\begin{tightright}
				Given August 31. 1831 by Joseph the Seer
				\end{tightright}
				
		% -----
		\subsubsection{Revelation, 11 September 1831 [D\&C 64]}
		% -----
			\label{EMS-1831-JS-11-SEP-1831}
			% \label{EMS-1830-JS-DAY-3LETTERMONTH-YEAR}
		
			\begin{description}
				\item[Print Sources]
			\end{description}

				\begin{tabular}{p{0.5in}p{1in}p{0.5in}p{1in}}
				JSP	 	& D2:61--67		& JSR 		& 158--161 \\
				DC	 	& Section 64 	& HC 		& 1:211--214 \\
				HDDC 	& 811--829		& TCHC 		& 1:149--150 \\
				\end{tabular}
				% HDDC = woodford's DC diss; JSR = Marquardt's JS Revelations; TCHC = Vogel HC
			
			\begin{description}
				\item[Online Access] \href{https://www.josephsmithpapers.org/paper-summary/revelation-11-september-1831-dc-64/1}{Here}
				% \item[Analysis] \hyperref[DC64]{Here}
			\end{description}
			
			\begin{description}
				\item[Background]
			\end{description}
			
				% It might also be useful to give a two sentence context of the period: e.g., "This speech was given in Nauvoo to a small congregation one month prior to the destruction of the Nauvoo Expositor. These and other tensions may have been on the prophet's mind when he gave the speech."
			
			\begin{description}
				\item[Original Source]
			\end{description}
			
				% brief description of original source material, with reference to JSP or somewhere else for more info
				% Background material: it might be helpful to have a note that says whether a given document was presented as a speech, was written in a journal, etc.
				% Also a comment here about how reliable the source might be, or what shape it is in, if there are large omissions or parts that are hard to understand, and so on.
			
			\begin{description}
				\item[Text]
			\end{description}
			
				\begin{tightcenter}
				67 Revelation Kirtland Sept 11\supersu{th}\supers{.} 1831
				\end{tightcenter}
				
				Directions to the Elders \&c \&c
				
				Behold thus saith the Lord your God unto you O ye Elders of my Church hearken ye \& hear \& receive my will concerning you for verily I say unto you I will that ye should overcome the world wherefore I will have compassion upon you there are those among you who have sinned but verily I say for this once for mine own glory \& for the salvation of Souls I have forgiven you your sins I will be mercyfull unto you for I have given unto you the Kingdom \& the keys of the mysteries of the kingdom shall not be taken from my Servent Joseph while he liveth in-as-much as he obeyeth mine ordinances there are those who have sought occation against him without a cause nevertheless he hath sinned but verily I say unto you I the Lord forgiveth sins unto those who confess their sins before me \& ask forgiveness who have not sinned unto death my Deciples in days of old sought occasion against one an other \& forgave not one an other in their hearts \& for this evil they were afflicted \& sorely chastened wherefore I say unto you that ye had ought to forgive one another for he that forgiveth not his brother his tresspasses standeth condemned before the Lord for there remaineth in him the greater sin I the Lord will forgive whom I will forgive but of you it is required to forgive all men \& ye had ought to say in your hearts let God Judge between me \& thee \& reward thee according to thy deeds \& he that repenteth not of his sins \& confess them not then ye \sout{ye} shall bring him before the church \& do with him as the Scriptures \sout{direct} Saith unto you either by commandment or <by> revelation \& this ye shall do that God might be glorified not because ye forgive not having not compassion but that ye may be Justified in the eyes of the law that ye may not offend him who is your lawgiver verily I say for this cause ye shall do these things Behold I the Lord was angery with him who was my Servent Ezra [Booth] \& also my Servent Isaac [Morley] for they kept not the Law neither the commandment they sought evil in their hearts \& I the Lord withheld my spirit they condemned for evil that thing in which there was no evil nevertheless I have forgiven my Servent Isaac \& also my Servent Edward [Partridge] Behold he hath sinned \& Satan Seeketh to destroy his soul but when these things are made known unto them they repenteth Of the evil \& they shall be forgiven \& now verily I say that it is expedient in me that my servent Sidney (Gilbert) after a few weeks should return upon his business \& to his agency in the Land of Zion \& that which he hath seen \& heard may be made known unto my Deciples that they perish not \& for this cause have I spoken these things \& again I say unto you that my servent Isaac may not be tempted above that which he is able to bear \& council wrongfully to your hurt I gave commandment that this farm should be sold I willeth not that my servent Frederick [G. Williams] should sell his farm for I the Lord willeth to retain a strong hold in the Land of Kirtland for the space of five years in the which I will not overthrow the wicked that thereby I may save some \& after that day I the Lord will not hold any guilty that shall go with open hearts up to the Land of Zion for I the Lord requireth the hearts of the Children of men Behold now it is called to day \& verily it is a day of Sacrifice \& a day for the tithing of my People for he that is tithed shall not be burned for after to day cometh the burning this is speaking after the manner of the Lord for verily I say tomorrow all the proud \& they that do wickedly shall be as stuble \& I will burn them up saith the Lord for I am the Lord of hosts \& I will not spare any that remaineth in Babylon wherefore if ye believe me ye will labour while it is called to day \& it is not meet that my servent Newel [K. Whitney] \& sidney [Gilbert] should sell their Store \& their Possessions here for this is not wisdom untill the residue of the Church which remaineth in this place shall go up unto the Land of Zion behold it is said in my Laws or forbidden to get in debt to thine enemies but Behold it is not said at any time that the Lord should not take when he please \& pay as seemeth him good wherefore as ye are agents \& ye are on the Lords errand \& whatever ye do according to the will of the Lord is the Lords business \& it is the Lords business to provide for his saints in these last days that they may obtain an in heritance in the Land of Zion \& Behold I the Lord declare unto you \& my words are shure \& shall not fail that they shall obtain it but all things must come to pass in its time wherefore be not weary in well doing for ye are laying the foundation of a great work \& out of small things proceedeth that which is great behold the Lord requireth the heart\sout{s} \& a willing mind \& the willing \& obedient shall eat the good of the Land of Zion in these Last days \& the rebelious shall be cut off out of the Land of Zion \& shall be sent away \& shall not inherit the Land for verily I say that the rebelious are not of the blood of Ephraim wherefore they shall be plucked out Behold I the Lord have made my Church in these last days like unto a Judge setting on a hill or in an high place to Judge the Nations for it shall come to pass that the inhabitants of Zion shall Judge all things \& all liars \& hypocrites shall be proved by them \& they which are not Apostles shall be known \& even the Judge \& his councellors if they are not faithfull in their stewardship shall be condemned \& \sout{an} others shall be planted in their stead for behold I say unto you that Zion shall flourish \& the glory of the Lord shall be upon her \& she shall be an ensighn unto the People \& these \sout{shall} shall come unto her out of every Nation under heaven \& the days shall come when the Nations of the Earth shall tremble because of her \& shall fear because of her terable ones the Lord hath spoken it amen
				
				Given at Kirtland September the 11\supersu{th}\supers{.} 1831
				
		% -----
		\subsubsection{Minutes, 11 October 1831}
		% -----
			\label{EMS-1831-JS-11-OCT-1831}
			% \label{EMS-1830-JS-DAY-3LETTERMONTH-YEAR}
		
			\begin{description}
				\item[Print Sources]
			\end{description}

				\begin{tabular}{p{0.5in}p{1in}p{0.5in}p{1in}}
				JSP	 	& D2:74--76		& JSR 		& --- \\
				DC	 	& ---		 	& HC 		& 1:219 \\
				HDDC 	& ---			& TCHC 		& 1:153--154 \\
				\end{tabular}
				% HDDC = woodford's DC diss; JSR = Marquardt's JS Revelations; TCHC = Vogel HC
			
			\begin{description}
				\item[Online Access] \href{https://www.josephsmithpapers.org/paper-summary/minutes-11-october-1831/1}{Here}
				% \item[Analysis] \hyperref[XX]{Here}
			\end{description}
			
			\begin{description}
				\item[Background]
			\end{description}
			
				% It might also be useful to give a two sentence context of the period: e.g., "This speech was given in Nauvoo to a small congregation one month prior to the destruction of the Nauvoo Expositor. These and other tensions may have been on the prophet's mind when he gave the speech."
			
			\begin{description}
				\item[Original Source]
			\end{description}
			
				% brief description of original source material, with reference to JSP or somewhere else for more info
				% Background material: it might be helpful to have a note that says whether a given document was presented as a speech, was written in a journal, etc.
				% Also a comment here about how reliable the source might be, or what shape it is in, if there are large omissions or parts that are hard to understand, and so on.
			
			\begin{description}
				\item[Text]
			\end{description}
			
				...Certain points were discussed by br Joseph Smith Jr, who said that the Elders present were to tarry untill the Morrow \& hold a meeting so that the members might understand the ancient manner of conducting meetings as they were led by the Holy Ghost Also said that this was not perfectly known by many of the Elders of this Church...
				
		% -----
		\subsubsection{Minutes, 25--26 October 1831}
		% -----
			\label{EMS-1831-JS-25-26-OCT-1831}
			% \label{EMS-1830-JS-DAY-3LETTERMONTH-YEAR}
		
			\begin{description}
				\item[Print Sources]
			\end{description}

				\begin{tabular}{p{0.5in}p{1in}p{0.5in}p{1in}}
				JSP	 	& D2:78--87		& JSR 		& --- \\
				DC	 	& ---		 	& HC 		& 1:219 \\
				HDDC 	& ---			& TCHC 		& 1:154 \\
				TPJS 	& 8--9			& 	 		&  \\
				\end{tabular}
				% HDDC = woodford's DC diss; JSR = Marquardt's JS Revelations; TCHC = Vogel HC
			
			\begin{description}
				\item[Online Access] \href{https://www.josephsmithpapers.org/paper-summary/minutes-25-26-october-1831/1}{Here}
				% \item[Analysis] \hyperref[XX]{Here}
			\end{description}
			
			\begin{description}
				\item[Background]
			\end{description}
			
				% It might also be useful to give a two sentence context of the period: e.g., "This speech was given in Nauvoo to a small congregation one month prior to the destruction of the Nauvoo Expositor. These and other tensions may have been on the prophet's mind when he gave the speech."
			
			\begin{description}
				\item[Original Source]
			\end{description}
			
				% brief description of original source material, with reference to JSP or somewhere else for more info
				% Background material: it might be helpful to have a note that says whether a given document was presented as a speech, was written in a journal, etc.
				% Also a comment here about how reliable the source might be, or what shape it is in, if there are large omissions or parts that are hard to understand, and so on.
			
			\begin{description}
				\item[Text]
			\end{description}
			
					\hypertarget{EMS-1831-JS-25-26-OCT-1831-a}%
				...Prayer
					\marginnote{a}%
				by br Simeon Carter, Glorious things, were then sung Br Sidney Rigdon arose \& said that he wished to make a few observations connected with the object of our assembling our selves together, When God works all may know it, for he always answers the prayers of the Savior for he makes his children one, for he by his Holy Spirit binds their hearts from Earth to Heaven; And in this thing God had taught his children to sing a new song even about Zion which David Spoke of, \&c, God always bears testimony by his presence in counsil to his Elders when they assemble \sout{themselves} in perfect faith and humble themselves before the Lord and their wills being swallowed up in the will of God
			
					\hypertarget{EMS-1831-JS-25-26-OCT-1831-b}%
				Br 
					\marginnote{b}%
				Joseph Smith Jr said we have assembled together to do the business of the Lord and it is through the great mercy of our God that we are spared to assemble together, many of us have went at the command of the Lord in defience of every thing evil, and obtained blessings unspeakable in consequence of which, our names are sealed in the Lambs' Book of life, for the Lord has spoken it.
					\hypertarget{EMS-1831-JS-25-26-OCT-1831-c}%
				It
					\marginnote{c}%
				is the privilege of every Elder to speak of the things of God \&c, And could we all come together with one heart and one mind in perfect faith the vail might as well be rent to day as next week or any other time and if we will but cleanse ourselves and covenant before God, to serve him, it is our privilege to have an assurence that God will protect us at all times...
				
					\hypertarget{EMS-1831-JS-25-26-OCT-1831-d}%
				...Br.
					\marginnote{d}%
				Sidney Rigdon said that he supposed that it was, saying, I bear testimony that God will have a pure people who will give up all for Christ's sake and when this is done they will be sealed up unto eternal life.
				
					\hypertarget{EMS-1831-JS-25-26-OCT-1831-e}%
				...Br.
					\marginnote{e}%
				Joseph Smith jr. said that the order of the High priesthood is that they have power given them to seal up the Saints unto eternal life. And said it was the privilege of every Elder present to be ordained to the Highpriesthood...
				
					\hypertarget{EMS-1831-JS-25-26-OCT-1831-f}%
				...Br.
					\marginnote{f}%
				Hyrum Smith said that he thought best that the information of the coming forth of the book of Mormon be related by Joseph himself to the Elders present that all might know for themselves.
				
				Br. Joseph Smith jr. said that it was not intended to tell the world all the particulars of the coming forth of the book of Mormon, \& also said that it was not expedient for him to relate these things \&c...
				
					\hypertarget{EMS-1831-JS-25-26-OCT-1831-g}%
				...Br
					\marginnote{g}%
				Joseph Smith Jr said that he intended to do his duty before the Lord and hoped that the brethren would be patient, as they had a considerable distance. also said that the promise of God was that the greatest blessings which God had to bestow should be given to those who contributed to the support of his family while translating the fulness of the Scriptures; also said until we have perfect love we are liable to fall and when we have a testimony that our names are sealed in the Lamb's Book of life we have perfect love \& then it is impossible for false Christ's to decieve us.
					\hypertarget{EMS-1831-JS-25-26-OCT-1831-h}%
				also
					\marginnote{h}%
				said that the Lord held the Church bound to provide for the families of the absent Elders while proclaiming the Gospel: further said that \sout{the} God had often sealed up the heavens because of covetousness in the Church. Said that the Lord would cut his work short in righteousness and except the church recieve the fulness of the Scriptures that they would yet fall.
				
				Adjournment moved for by the moderator until evening.

					\hypertarget{EMS-1831-JS-25-26-OCT-1831-i}%
				The
					\marginnote{i}%
				Clerk begged the privilege of laying before the Conference in the evening the necessity of sending certain Elders to visit the Churches and obtain means for the support of br. Joseph Smith jr. \& those appointed to assist him in writing \& copying the fulness of the Scriptures, as was concluded upon in a conference held in Hiram, Oct. 11. 1831...
				
					\hypertarget{EMS-1831-JS-25-26-OCT-1831-j}%
				...Br
					\marginnote{j}%
				Joseph Smith Jr, was appointed to examine these brethren, presenting themselves for ordination, after prayer said that he had a testimony that each had one tallent and if after being ordained they should hide it God would take it from them; exhorted them to pray continually in meekness...
				
		% -----
		\subsubsection{Revelation, 29 October 1831 [D\&C 66]}
		% -----
			\label{EMS-1831-JS-29-OCT-1831}
			% \label{EMS-1830-JS-DAY-3LETTERMONTH-YEAR}
		
			\begin{description}
				\item[Print Sources]
			\end{description}

				\begin{tabular}{p{0.5in}p{1in}p{0.5in}p{1in}}
				JSP	 	& D2:87--92		& JSR 		& 163--164 \\
				DC	 	& Section 66 	& HC 		& 1:220--221 \\
				HDDC 	& 838--846		& TCHC 		& 1:154--155 \\
				\end{tabular}
				% HDDC = woodford's DC diss; JSR = Marquardt's JS Revelations; TCHC = Vogel HC
			
			\begin{description}
				\item[Online Access] \href{https://www.josephsmithpapers.org/paper-summary/revelation-29-october-1831-dc-66/1}{Here}
				% \item[Analysis] \hyperref[DC66]{Here}
			\end{description}
			
			\begin{description}
				\item[Background]
			\end{description}
			
				% It might also be useful to give a two sentence context of the period: e.g., "This speech was given in Nauvoo to a small congregation one month prior to the destruction of the Nauvoo Expositor. These and other tensions may have been on the prophet's mind when he gave the speech."
			
			\begin{description}
				\item[Original Source]
			\end{description}
			
				% brief description of original source material, with reference to JSP or somewhere else for more info
				% Background material: it might be helpful to have a note that says whether a given document was presented as a speech, was written in a journal, etc.
				% Also a comment here about how reliable the source might be, or what shape it is in, if there are large omissions or parts that are hard to understand, and so on.
			
			\begin{description}
				\item[Text]
			\end{description}
			
				Behold thus saith the Lord u[n]to you my servant William [E. McLellin]. Blessed are you inasmuch as you have turned away from your inequities and have received my truths saith the Lord your Redeemer, the Saviour of the world, even of as many as believe on my name. Verily I say unto you blessed are you for receiving mine everlasting Covenant even the fulness of my Gospel sent forth unto the children of men that they might have life and be made partakers of the glories which are to be revealed in the last days as it was written by the prophets and Apostles in days of old. Verily I say unto you my servant W\supers{m.} that you are clean but not all Repent therefore of those things which are not pleasing in my sight saith the Lord for the Lord will shew them unto you. And now Verily I the \sout{the} Lord will shew unto you what I will concerning you or what is my will concerning you. Behold Verily I say unto you that it is my will that you should proclaim my Gospel from land to land and from City to City. Yea in those regions round about where it hath not been proclaimed. Tarry not many days in this place Go not up unto the Land of Zion as yet But in as much as you can send. Send, otherwis[e] think not of thy property. Go unto Eastern lands. Bear testimony in every place, unto every people and in their synagogues reasoning with the people Let my servant Samuel (Smith) go with <you> and forsake him not and give him thine instructions and he that is faithful shall be made strong in every place And I the Lord will go with you. Lay your hands upon the sick and they shall recover. Return not until I the Lord shall send you, Be patien[t] in afflictions. Ask and ye shall receive Knock and it shall be opened unto you Seek not to be cumbered. Forsake all unrighteousness. Commit not adultery, $\diamond$ temptation with which thou hast been troubled. Keep these sayings true and fait[h]ful and thou shalt magnify thine office and push many people to Zion with songs of everlasting joy upon their heads Continue in these things even unto the end and you shall have a Crown of etern[al] life on the right hand of my Father w[ho] is full of grace and truth. Verily thus saith the Lord your God, your Redeemer even Jesus Christ. Amen------
				
				A revelation given to W\supers{m.} E. M\supers{c}Lelin a true decendant from Joseph who was sold into Egypt down through the loins of Ephraim his son------
				
				Given in Portage Co. Ohio. Hiram Township on the 29\supers{th.} of October 1831
				
				\begin{tightright}
					\uline{Joseph Smith, Revelator}
				\end{tightright}
				
		% -----
		\subsubsection{Revelation, 30 October 1831 [D\&C 65]}
		% -----
			\label{EMS-1831-JS-30-OCT-1831}
			% \label{EMS-1830-JS-DAY-3LETTERMONTH-YEAR}
		
			\begin{description}
				\item[Print Sources]
			\end{description}

				\begin{tabular}{p{0.5in}p{1in}p{0.5in}p{1in}}
				JSP	 	& D2:92--94		& JSR 		& 164--165 \\
				DC	 	& Section 65 	& HC 		& 1:218 \\
				HDDC 	& 830--837		& TCHC 		& 1:153 \\
				\end{tabular}
				% HDDC = woodford's DC diss; JSR = Marquardt's JS Revelations; TCHC = Vogel HC
			
			\begin{description}
				\item[Online Access] \href{https://www.josephsmithpapers.org/paper-summary/revelation-30-october-1831-dc-65/1}{Here}
				% \item[Analysis] \hyperref[DC65]{Here}
			\end{description}
			
			\begin{description}
				\item[Background]
			\end{description}
			
				% It might also be useful to give a two sentence context of the period: e.g., "This speech was given in Nauvoo to a small congregation one month prior to the destruction of the Nauvoo Expositor. These and other tensions may have been on the prophet's mind when he gave the speech."
			
			\begin{description}
				\item[Original Source]
			\end{description}
			
				% brief description of original source material, with reference to JSP or somewhere else for more info
				% Background material: it might be helpful to have a note that says whether a given document was presented as a speech, was written in a journal, etc.
				% Also a comment here about how reliable the source might be, or what shape it is in, if there are large omissions or parts that are hard to understand, and so on.
			
			\begin{description}
				\item[Text]
			\end{description}
			
				\begin{tightcenter}
				69 Revelation Oc\supersu{t}\supers{.} 30\supersu{th}\supers{.} 1831
				\end{tightcenter}
				
				Hearken \& Lo a voice as one sent down from on high who is mighty \& powerfull whose going forth is unto the ends of the Earth yea whose voice is unto men prepare ye the way of the Lord make his paths strait The keys of the kingdom of God is committed unto man on the Earth \& from thence shall the Gospel roll forth unto the ends of the Earth as the stone which is hewn from the Mountain without hands shall roll forth untill it hath filled the whole Earth yea a voice crying prepare ye the way of the Lord prepare ye the supper of the Lamb make ready for the Bridegroom pray unto the Lord call upon his holy name make known his wonderfull work\supersu{s}\supers{.} among the people call upon the Lord that his kingdom may go forth upon the Earth that the inhabitants thereof may received it \& be prepared for the days to come in the which the Son of man Shall come down in heaven Clothed in the brightness of his glory to meet the kingdom of God which is set up on the Earth Wherefore may the kingdom of God go forth that the kingdom of heaven may come that thou O God may be glorified in heaven so on Earth that thine enemies may be subdued for thine is the \sout{kingdom} honour power \& glory forever \& ever Amen
				
		% -----
		\subsubsection{Minutes, 1--2 November 1831}
		% -----
			\label{EMS-1831-JS-1-2-NOV-1831}
			% \label{EMS-1830-JS-DAY-3LETTERMONTH-YEAR}
		
			\begin{description}
				\item[Print Sources]
			\end{description}

				\begin{tabular}{p{0.5in}p{1in}p{0.5in}p{1in}}
				JSP	 	& D2:94--98		& JSR 		& --- \\
				DC	 	& ---		 	& HC 		& 1:221--222 \\
				HDDC 	& ---			& TCHC 		& 1:155 \\
				\end{tabular}
				% HDDC = woodford's DC diss; JSR = Marquardt's JS Revelations; TCHC = Vogel HC
			
			\begin{description}
				\item[Online Access] \href{https://www.josephsmithpapers.org/paper-summary/minutes-1-2-november-1831/1}{Here}
				% \item[Analysis] \hyperref[XX]{Here}
			\end{description}
			
			\begin{description}
				\item[Background]
			\end{description}
			
				% It might also be useful to give a two sentence context of the period: e.g., "This speech was given in Nauvoo to a small congregation one month prior to the destruction of the Nauvoo Expositor. These and other tensions may have been on the prophet's mind when he gave the speech."
			
			\begin{description}
				\item[Original Source]
			\end{description}
			
				% brief description of original source material, with reference to JSP or somewhere else for more info
				% Background material: it might be helpful to have a note that says whether a given document was presented as a speech, was written in a journal, etc.
				% Also a comment here about how reliable the source might be, or what shape it is in, if there are large omissions or parts that are hard to understand, and so on.
			
			\begin{description}
				\item[Text]
			\end{description}
			
					\hypertarget{EMS-1831-JS-1-2-NOV-1831-a}%
				...Br.
					\marginnote{a}%
				Joseph Smith jr. said that inasmuch as the Lord had bestowed a great blessing upon us in giving commandments and revelations, asked the conference what testimony they were willing to attach to these commandments which should shortly be sent to the world. A number of the brethren arose and said that they were willing to testify to the world that they knew that they were of the Lord.
				
					\hypertarget{EMS-1831-JS-1-2-NOV-1831-b}%
				Revelation
					\marginnote{b}
				received relative to the same...
				
				The Revelation of last evening read by the moderator the brethren then arose in turn and bore witness to the truth of the Book of Commandments. After which br. Joseph Smith jr. arose \& expressed his feelings \& gratitude concerning the commandment \& Preface received yesterday.
				
		% -----
		\subsubsection{Revelation, 1 November 1831--A [D\&C 68]}
		% -----
			\label{EMS-1831-JS-1-NOV-1831-A}
			% \label{EMS-1830-JS-DAY-3LETTERMONTH-YEAR}
		
			\begin{description}
				\item[Print Sources]
			\end{description}

				\begin{tabular}{p{0.5in}p{1in}p{0.5in}p{1in}}
				JSP	 	& D2:98--103	& JSR 		& 169--172 \\
				DC	 	& Section 68 	& HC 		& 1:227--229 \\
				HDDC 	& 853--865		& TCHC 		& 1:158--160 \\
				\end{tabular}
				% HDDC = woodford's DC diss; JSR = Marquardt's JS Revelations; TCHC = Vogel HC
			
			\begin{description}
				\item[Online Access] \href{https://www.josephsmithpapers.org/paper-summary/revelation-1-november-1831-a-dc-68/1}{Here}
				% \item[Analysis] \hyperref[DC68]{Here}
			\end{description}
			
			\begin{description}
				\item[Background]
			\end{description}
			
				% It might also be useful to give a two sentence context of the period: e.g., "This speech was given in Nauvoo to a small congregation one month prior to the destruction of the Nauvoo Expositor. These and other tensions may have been on the prophet's mind when he gave the speech."
			
			\begin{description}
				\item[Original Source]
			\end{description}
			
				% brief description of original source material, with reference to JSP or somewhere else for more info
				% Background material: it might be helpful to have a note that says whether a given document was presented as a speech, was written in a journal, etc.
				% Also a comment here about how reliable the source might be, or what shape it is in, if there are large omissions or parts that are hard to understand, and so on.
			
			\begin{description}
				\item[Text]
			\end{description}
			
				given in Hiram Nov. 1. 1831
				
				\begin{tightcenter}
				70 A Revelation to Orson [Hyde] Luke [Johnson] \sout{\&} Lyman [Johnson] \& William [E. McLellin]
				\end{tightcenter}
				
				The mind \& will of the Lord as made known by the voice of the spirit \sout{made} <\sout{known} to> a confrence held November first 1831 concerning certain Elders who requested of the Lord to kno$\diamond$ his will concerning them \& also certain items as made known in addition to the Laws \& commandments which have been given to the church firstly my servant Orson was called by his ordinance to proclaim the everlasting Gospel by the spirit of the living God from people to people \& from land to land \sout{from} <in the> congregtions of the wicked in their Synagogues reas[o]ning with \& expounding all scriptures <un>to them \& behold \& lo this is an ensample unto all those who were ordained unto this priesthood whose mission is appointed unto them to go forth \& this is the ensample unto them that they shall speak as they are moved upon by the Holy Ghost \& whatsoever they shall speak when moved upon by the Holy Ghost shall be Scripture shall be the will of the Lord shall be the mind of the Lord shall be the voice of the Lord \& shall be the power of God unto Salvation behold this is the promise of the Lord unto you o ye my servants wherefore be of good cheer \& do not fear for I the Lord am with you \& will stand by you \& \sout{you} <ye> shall bear record of me even Jesus christ that I am the Son of God that I was that I am \& that I am to come this is the word of the Lord unto you my Servant Orson \& also to my servant Luke \& unto my servant Lyman \& unto my servant William \& unto all the faithful Elders of my church go ye unto all the world preach the gospel to every creature acting in the authority which I have given you baptising in the name of the Father \& of the Son \& \sout{of} <of the> Holy Ghost \& he that believeth \& is baptised shall be saved \& he that believeth not shall be damned \& he that believeth shall be blessed with signs following even as it is written \& unto you it shall be given to know the signs of the times \& the signs of the coming of the Son of man \& of as many as the Father shall bear record to you it shall be given to seal them up unto Eternal life Amen---
				
				And now conc[e]rning the items in addition to the Laws \& commandments they are these there rema[i]neth hereafter in the due time of the Lord other Bishops to be set apart unto the church to minister even according to the first wherefore it shall be an high priest who is worthy \& he shall be appointed by a confrenc of high priests And again no Bishop or judge which shall be set apart for this ministry shall be tried or condemned for any crime save it be before a confrence of high priests \& inasmuch as he is found guilty before a confrenc of high priests by testimony that cannot be impeached he shall be condemned or forgiven according to the Laws of the church And again inasmuch as parents have children in Zion that teach them not to understand the doctrine of repentance faith in Christ the Son of the living God \& of baptism \& the gift of the Holy Spirit by the laying <on> of the hands when eight years old the sin be upon the head of the parents for this shall be a Law unto the inhabitants of Zion \& their children shall be baptised for the remission of their sins when eight years old \& receive the laying on of the hands \& they also shall teach their children to pray \& to walk uprightly before the Lord \& the inhabitants of Zion shall also observe \sout{to} the Sabath day to keep it holy \& the inhabitants of Zion also shall remember their labors inasmuch as they are appointed to labor in all faithfulness for the idler shall be had in remembrance before the Lord now I the Lord am not well pleased with the inhabitants of Zion for there are idlers among them \& their children also are growing up in wickedness they also seek not earnestly the riches of Eternity but their eyes are full of greediness these things ought not to be \& must be done away from among them wherefore let my servant Oliver [Cowdery] cary these sayings unto the land of Zion \& a commandment I give unto them that he that obse[r]veth <not> his prayers before the Lord in the season thereof let them be had in remembrance before the judge of my people these sayings are true \& faithful wherefore transgress them not neither take therefrom behold I am Alpha \& Omega \& I come quickly Amen
				
				Given \sout{a} in Hiram November first 1831 by Joseph the Seer 
				
		% -----3j
		\subsubsection{Revelation, 1 November 1831--B [D\&C 1]}
		% -----
			\label{EMS-1831-JS-1-NOV-1831-B}
			% \label{EMS-1830-JS-DAY-3LETTERMONTH-YEAR}
		
			\begin{description}
				\item[Print Sources]
			\end{description}

				\begin{tabular}{p{0.5in}p{1in}p{0.5in}p{1in}}
				JSP	 	& D2:103--107	& JSR 		& 165--167 \\
				DC	 	& Section 1	 	& HC 		& 1:222--224 \\
				HDDC 	& 116--125		& TCHC 		& 1:155--156 \\
				\end{tabular}
				% HDDC = woodford's DC diss; JSR = Marquardt's JS Revelations; TCHC = Vogel HC
			
			\begin{description}
				\item[Online Access] \href{https://www.josephsmithpapers.org/paper-summary/revelation-1-november-1831-b-dc-1/1}{Here}
				% \item[Analysis] \hyperref[DC1]{Here}
			\end{description}
			
			\begin{description}
				\item[Background]
			\end{description}
			
				% It might also be useful to give a two sentence context of the period: e.g., "This speech was given in Nauvoo to a small congregation one month prior to the destruction of the Nauvoo Expositor. These and other tensions may have been on the prophet's mind when he gave the speech."
			
			\begin{description}
				\item[Original Source]
			\end{description}
			
				% brief description of original source material, with reference to JSP or somewhere else for more info
				% Background material: it might be helpful to have a note that says whether a given document was presented as a speech, was written in a journal, etc.
				% Also a comment here about how reliable the source might be, or what shape it is in, if there are large omissions or parts that are hard to understand, and so on.
			
			\begin{description}
				\item[Text]
			\end{description}
			
				\begin{tightcenter}
				<77 Revelation>
				\end{tightcenter}
				
				Given in Hiram Nov\supersu{m}\supers{.} 1\supersu{st}\supers{.} 1831
				
				A Preface or instructions upon the Book of Commandments which were given of the Lord unto his Church through him whom he appointed to this work by the voice of his Saints through the prayer of faith this church being organized according to the will of him who rules all things on the Sixth day of April in the year of our Lord 1830
				
				Hearken O ye People of my Church saith the voice of him who dwells on high \& whose eyes are upon all men yea verily I say hearken ye People from afar \& ye that are upon the Islands of the sea listen to gether for verily the voice of the Lord is unto all men \& there is none to escape \& there is no eye that shall not see neither ear that shall not hear neither heart that shall not be penetrated \& the rebelious shall be pierced with much sorrow for their iniquities shall be spoken upon the house tops \& their seceret acts shall be revealed \& the voice of warning shall be unto all people by the mouth of my Deciples whom I have chosen in these last days they shall go forth \& none shall stay them for I the Lord have commanded them Behold this is mine authority \& the authority of my servents \& my preface unto the Book of my Commandments which I have given them to Publish unto you O Inhabitants of the Earth wherefore fear \& tremble O ye People for what I the Lord have decreed in them shall be fulfilled \& verily I say unto you that they who go forth bearing these tidings unto the Inhabitants of the Earth to them is power given to seal both on Earth \& in Heaven the unbelieveing \& rebelious yea verily to seal them up unto the day when the wrath of God shall be poured out upon the wicked without measure unto the day when the Lord shall come to recompence unto every man according to his works \& measure to every man according to the measure which he has measured to his fellow man wherefore the voice of the Lord is unto \sout{his fellow man} unto the end of the Earth that all that will hear may hear prepare ye prepare ye for that which is to come for the Lord is nigh \& the anger of the Lord is kindled \& his sword is bathed in heaven \& it shall fall upon the inhabitants of the Earth \& the arm of the Lord shall be revealled \& the day cometh that they who will not hear the voice of the Lord neither his servants neith[er] give heed to the words of the Prophets \& Apostles shall be cut off from among the People for they have strayed from mine ordinances \& have broken mine everlasting Covenant they seek not the Lord to establish his righteousness but every man walketh in his own way \& after the Image of his own God whose Image is in the likeness of the world \& whose substance is that of an Idol which waxeth old \& shall perish in Babylon even Babylon the great, which shall fall wherefore I the Lord knowing the calamity which should come upon the inhabitants of the Earth called upon my Servent\sout{s} Joseph \& spake unto him from heaven \& gave him commandment \& also \sout{\&} gave commandments to others that they should proclaim these things unto the world \& all this that it might be fulfilled which was written by the Prophets the weak things of the world should come forth \& break down the mighty \& strong ones that \sout{men} <man> should not council his fellow man neither trust in the arm of flesh but that every man might Speak in the name of God the Lord even the Saveiour of the world that faith also might increase in the Earth that mine everlasting Covenant might be established that the fullness of my Gospel might be proclaimed by the weak \& the Simple unto the ends of the world \& before kings \& Rulers Behold I am God \& have spoken it these \sout{are} commandments are of me \& were given unto my Servents in their weakness after the manner of their Language that they might come to understanding \& in as much as they erred it might be made known \& in as much as they sought wisdom it might be \sout{made known} instructed \& in as much as they sinned they might be chastened that they might repent \& in as much as they were humble they might be made strong \& blessed from on high \& receive knowledge from time to time After they having received the record of the Nephites yea even my Servant Joseph might have power to translate through the mercy of God by the power of [God] the Book of Mormon \& also those to whom these commandments were given might have power to lay the foundation of this Church \& to bring it forth out of obscurity \& out of darkness the only true \& living Church upon the face of the whole Earth with which I the Lord am well pleased speaking unto the Church collectively \& not individually for I the Lord cannot look upon sin with the least degree of allowance nevertheless he that repenteth \& doeth the commandments of the Lord shall be forgiven \& he that repenteth not from him shall be taken even the light which he has received for my spirit shall not always strive with man saith the Lord of hosts \& again verily I say unto you O inhabitants of the Earth for I the Lord am willing to make these things known unto all flesh for I am no respector to persons \& willeth that all men shall know that the day speedily cometh the hour is not yet but is nigh at hand when peace shall be taken from the Earth \& the Devil shall have power over his own dominion \& also the Lord shall have power over his saints \& shall reign in their midst \& shall come down in Judgement upon Idumea (or the World) search these commandments for they are true \& faithfull \& the Prophecies \& promises which are in them shall all be fulfilled what I the Lord have spoken I have spoken \& I excuse not myself \& though the Heaven \& <the> Earth pass away my word shall not pass away but shall all be fulfilled whether by mine own voice or by the voice of my Servants it is the same for Behold \& Lo the Lord is God \& the Spirit beareth record \& the [record] is true \& the truth abideth for ever \& ever Amen
				
		% -----
		\subsubsection{Revelation, circa 2 November 1831 [D\&C 67]}
		% -----
			\label{EMS-1831-JS-CIR-2-NOV-1831}
			% \label{EMS-1830-JS-DAY-3LETTERMONTH-YEAR}
		
			\begin{description}
				\item[Print Sources]
			\end{description}

				\begin{tabular}{p{0.5in}p{1in}p{0.5in}p{1in}}
				JSP	 	& D2:108--110	& JSR 		& 167--168 \\
				DC	 	& Section 67 	& HC 		& 1:225 \\
				HDDC 	& 847--852		& TCHC 		& 1:156--157 \\
				\end{tabular}
				% HDDC = woodford's DC diss; JSR = Marquardt's JS Revelations; TCHC = Vogel HC
			
			\begin{description}
				\item[Online Access] \href{https://www.josephsmithpapers.org/paper-summary/revelation-circa-2-november-1831-dc-67/1}{Here}
				% \item[Analysis] \hyperref[DC67]{Here}
			\end{description}
			
			\begin{description}
				\item[Background]
			\end{description}
			
				% It might also be useful to give a two sentence context of the period: e.g., "This speech was given in Nauvoo to a small congregation one month prior to the destruction of the Nauvoo Expositor. These and other tensions may have been on the prophet's mind when he gave the speech."
			
			\begin{description}
				\item[Original Source]
			\end{description}
			
				% brief description of original source material, with reference to JSP or somewhere else for more info
				% Background material: it might be helpful to have a note that says whether a given document was presented as a speech, was written in a journal, etc.
				% Also a comment here about how reliable the source might be, or what shape it is in, if there are large omissions or parts that are hard to understand, and so on.
			
			\begin{description}
				\item[Text]
			\end{description}
			
				\begin{tightcenter}
				71 Revelation given Nov 2\supersu{nd}\supers{.} 1831
				\end{tightcenter}
				
				Behold \& hearken oh ye Elders of my Church who have assembelled yourselves together whose prayers I have heard \& whose \sout{hearts} desires have come up before me behold \& Lo mine eyes are upon you \& the heavens \& the earth are in mine hands \& the riches of eternity are mine to give[.] ye endeavour to believe that \sout{you} ye should receive the blessing which was offered unto you but behold verily I say unto you there were fears in your hearts \& verily this is the reason that ye did not receive \& now I the Lord give unto you a testimony of the truth of those commandments which \sout{I} are lying before you your eyes have been upon my Servent Joseph \& his language you have known \& his imperfections you have known \& you have sought in your hearts knowlege that you might express beyond his language this you also know now seek ye out of the Book of commandments even the least that is among them \& appoint him that is the most wise among you or if there be any among you that shall make one like unto it then ye are Justified in saying that ye do not know that is true but if you cannot make one like unto it ye are under condemnation if \sout{you} ye do not bear that it is true for ye know that there is no unrighteousness in it \& that which is righteous cometh down from above from the father of lights \& again verily I say unto you that it is your privilege \& a promise I give unto you that have been ordained unto the ministry that in as much as ye strip yourselves from Jealesies \& fears \& humble yourselves before me for ye are not sufficiently humble the veil shall \sout{not} be wrent \& you shall see me \& know that I am not with the carnal neither natural but with the spiritual for no man hath seen God at any time in the flesh but by the Spirit of God neither can any natural man abide the presence of God neither after the carnal mind ye are not \sout{to} able to abide the presence of God now neither the ministering of Angels wherefore continue in patience untill ye are perfected let not your minds turn back \& when ye are worthy in mine own due time ye shall see \& know that which was confirmed <upon you> by the hands of my Ser[v]ant Joseph Amen
				
		% -----
		\subsubsection{Testimony, circa 2 November 1831}
		% -----
			\label{EMS-1831-JS-CIR-2-NOV-1831-TEST}
			% \label{EMS-1830-JS-DAY-3LETTERMONTH-YEAR}
		
			\begin{description}
				\item[Print Sources]
			\end{description}

				\begin{tabular}{p{0.5in}p{1in}p{0.5in}p{1in}}
				JSP	 	& D2:110--114	& JSR 		& 168--169 \\
				DC	 	& --- 			& HC 		& 1:226 \\
				HDDC 	& ---			& TCHC 		& 1:157--158 \\
				\end{tabular}
				% HDDC = woodford's DC diss; JSR = Marquardt's JS Revelations; TCHC = Vogel HC
			
			\begin{description}
				\item[Online Access] \href{https://www.josephsmithpapers.org/paper-summary/testimony-circa-2-november-1831/1}{Here}
				% \item[Analysis] \hyperref[XX]{Here}
			\end{description}
			
			\begin{description}
				\item[Background]
			\end{description}
			
				% It might also be useful to give a two sentence context of the period: e.g., "This speech was given in Nauvoo to a small congregation one month prior to the destruction of the Nauvoo Expositor. These and other tensions may have been on the prophet's mind when he gave the speech."
			
			\begin{description}
				\item[Original Source]
			\end{description}
			
				% brief description of original source material, with reference to JSP or somewhere else for more info
				% Background material: it might be helpful to have a note that says whether a given document was presented as a speech, was written in a journal, etc.
				% Also a comment here about how reliable the source might be, or what shape it is in, if there are large omissions or parts that are hard to understand, and so on.
			
			\begin{description}
				\item[Text]
			\end{description}
			
				\begin{tightcenter}
					\hypertarget{EMS-1831-JS-CIR-2-NOV-1831-TEST-a}%
				<73
					\marginnote{a}
				Revelation>
				\end{tightcenter}
				
				The Testimony of the witnesses to the Book of the Lords commandments which he gave to his church through Joseph Smith Jr who was appointed by the vos [voice] of the Church for this purpose
				
					\hypertarget{EMS-1831-JS-CIR-2-NOV-1831-TEST-b}%
				We
					\marginnote{b}%
				the undersigners feel willing to bear testimony to all the world of mankind to every creature upon all the face of all the Earth upon the Islands of the Sea that god hath born record to our souls through the Holy Ghost shed forth upon us that these commandments are given by inspiration of God \& are profitable for all men \& are verily true we give this testimony unto the world the Lord being my helper \& it is through the grace of God the father \& his Son Jesus Christ that we are permitted to have this privelege of bearing this testimony unto the world in the which we rejoice exceedingly by praying the Lord always that the children of men may be profited thereby Amen
				
				\begin{multicols}{2}
						\hypertarget{EMS-1831-JS-CIR-2-NOV-1831-TEST-c}%
					Joshua
						\marginnote{c}%
					Fairchild
					
					Peter Dustin
					
					Newel Knight
					
					Levi Hancock<; never to be eraised>
					
					Thomas B Marsh
					
					\columnbreak
					
					Sidney Rigdon
					
					Orson Hyde
					
					W\supersu{m}\supers{.}. E. Mc.lel[l]in
					
					Luke Johnson
					
					Lyman Johnson
					
					Reynolds Cahoon
					
					John Corrill
					
					Parley [P.] Pratt
					
					Harv[e]y Whitlock
					
					Lyman Wight
					
					John Murdock
					
					Calvin Beebe
					
					Zebedee Coltrin
				\end{multicols}
				
		% -----
		\subsubsection{Revelation, 3 November 1831 [D\&C 133]}
		% -----
			\label{EMS-1831-JS-3-NOV-1831}
			% \label{EMS-1830-JS-DAY-3LETTERMONTH-YEAR}
		
			\begin{description}
				\item[Print Sources]
			\end{description}

				\begin{tabular}{p{0.5in}p{1in}p{0.5in}p{1in}}
				JSP	 	& D2:114--121	& JSR 		& 173--176 \\
				DC	 	& Section 133 	& HC 		& 1:229--234 \\
				HDDC 	& 1762--1783	& TCHC 		& 1:160--163 \\
				\end{tabular}
				% HDDC = woodford's DC diss; JSR = Marquardt's JS Revelations; TCHC = Vogel HC
			
			\begin{description}
				\item[Online Access] \href{https://www.josephsmithpapers.org/paper-summary/revelation-3-november-1831-dc-133/1}{Here}
				% \item[Analysis] \hyperref[DC133]{Here}
			\end{description}
			
			\begin{description}
				\item[Background]
			\end{description}
			
				% It might also be useful to give a two sentence context of the period: e.g., "This speech was given in Nauvoo to a small congregation one month prior to the destruction of the Nauvoo Expositor. These and other tensions may have been on the prophet's mind when he gave the speech."
			
			\begin{description}
				\item[Original Source]
			\end{description}
			
				% brief description of original source material, with reference to JSP or somewhere else for more info
				% Background material: it might be helpful to have a note that says whether a given document was presented as a speech, was written in a journal, etc.
				% Also a comment here about how reliable the source might be, or what shape it is in, if there are large omissions or parts that are hard to understand, and so on.
			
			\begin{description}
				\item[Text]
			\end{description}
			
				<72 A Revelation Rec\uline{d}. Nov 3, 1831>
				
				Hearken oh ye People of my Church------ saith the Lord your <God> \& hear the word of the Lord concerning you the Lord who shall suddenly come to his temple the Lord who shall come down with a curse to Judgement yea upon all the Nations that forget god \& upon all the ungodly amongst you for he shall make bear his holy arm in the eyes of all the Nations \& all the ends of the earth shall see the salvation of their god wherefore prepare ye prepare ye oh ye my People Sanctify yourselves gether ye together oh ye People of my Church upon the Land of Zion all you that have not been commanded to tarry go ye out from Babylon be ye clean that bear the vessels of the Lord call your solom assemblies \& speak often one to another \& let every man call upon the name of the Lord yea verily I say unto you again the time has come when the voice of the Lord is unto you go ye out of Babylon gether ye out from among the nations from the four winds from one end of Heaven to the other send forth the Elders of my Church unto the nations which are afar off unto the ilands of the sea send forth unto foreign lands call upon all nations firstly upon the gentiles \& then upon the Jews \& Behold \& Lo this shall be their Cry \& the voice of the Lord unto all People go ye forth unto the Land of Zion that the borders of my People may be enlarged \& that her stakes may be strengthened that Zion may go forth---unto the regions round about--- yea let the cry go forth among all people awake \& arise \& go forth to meet the Bride-groom Behold \& Lo the Bride-groom Cometh--- go ye out to meet him prepare yourselves for the great day of the Lord watch therefore for ye know neither the day nor the hour let them therefore which are among the gentiles--- flee unto Zion \& let they which be of Judah flee unto Jerusalem unto the Mountains of the Lords house go ye out from among the Nations even from Babylon From the midst of wickedness which is spiritual babylon But verily thus saith the Lord let not your flight \sout{not} be in h$\diamond$ste but let all things be prepared before you \& he that goeth let him not [look?] back--- lest sudden distruction shall come upon him hearken \& hear oh ye inhabitants of the Earth \& listen ye Elders of my Church together \& hear the voice of the Lord for he calleth upon all men \& he comandeth all men every where to repent for behold the Lord God hath sent forth the Angel \sout{with the everlasting gospel} crying through the midst of Heaven saying prepare ye the way of the Lord \& make his paths strait for the hour of his coming is nigh when the Lamb Shall stand upon Mount Zion \& with him a hundred \& forty four thousand having his fathers name written in their foreheads--- wherefore prepare ye for the coming of the Bride-groom go ye g[o] y$\diamond\diamond$ out to meet him for Behold he shall stand upon the Mount of Olivet \& upon the mighty Ocean even the great deep \& upon the Islands of the Sea \& upon the Land of Zion \& he shall utter his voice out of Zion \& he shall speak from Jerusalem \& his voice shall be heard among all people \& it shall be as the voice of many waters \& as the voice of a great thunder which shall break down the Mountains \& the valies shall not be found he shall command the great deep \& it shall be driven back into the North countries \& the Islands shall become one land \& the land of Jerusalem \& the Land of Zion shall be turned back into their own place \& the earth Shall be like as it was in the days before it was \sout{before it was} divided \& the Lord even the Saviour shall stand in the midst of his people <\& shall reign> over all \sout{the} \sout{Earth} flesh \& they \sout{which} who are in the North countries shall come in rememberanc before the Lord \& their Prophets shall hear his voice and shall no longer stay themselves \& they shall smite the rocks \& the ice shall folow down at their presenc[e] \& an high way shall be cast up in the midst of the great deep their enemies shall become a prey unto them \& <in> the barren deserts there shall come forth pools of living water \& the parched ground shall <no> longe[r] be a thirsty land \& they shall bring forth th[e]ir rich treasures unto the Children of Ephraim my servents \& the boundaries of the everlasting hills shall trembl at their presence \& these shall thy fall down \& be crowned with glory even in Zion by the hands of the Servents of the Lord even the children of Ephraim \& they shall be filled with songs of everlasting Joy behold this is the blessing of the everlasting God upon the heads of the tribes of Israel \& the richer blessing upon the head of Ephraim \& his fellows \& they also on [of] the tribe of Judah after their pain shall be Sanctified in holiness before the Lord to dwell in his presenc[e] day \& night forever \& ever \& now verily saith the Lord that these things might be known among you oh inhabitants of the Earth I have sent forth mine Angel flying throug[h] the midst of heaven having the everlasting Gospel who hath appeared unto some \& hath committed it unto man who shall appear unto many that dwell on the Earth \& this gospel shall be preached unto every Nation \& kindred \& tongue \& People \& the Servents of God shall go forth saying with a loud voice fear God \& give glory to him for the hour of his Judgement is come \& worship him that made Heaven \& earth \& the Sea \& the fountain of waters calling upon the Lord day \& night saying oh that thou wouldst rend the heavens that thou wouldest come down that the mountains would flow down at thy presence \& it shall be answered upon their heads for the presence of the Lord shall be as the melting fire that burneth \& as the fire that causeth the waters to boil oh Lord thou shalt come down to make known thy name to thine advisary \& all nations shall trembl at thy presence when thou doeth terabl things things that they look not for yea when thou comest down \& the Mountains flow down at thy presenc[e] thou shalt meet him that rejoiceth \& worketh righteousness who remember thee in thy ways for sinc[e] the begining of the world have not man heard nor perceived by the Ear neither hath the eye seen O God besides thee how great things thou prepared for him that waiteth for thee \& [it] shall be said who is this that cometh down from god in heavn with thy garments yea from the regions that is not known clothed in his gloriou[s] appearl travling in the greatness of his strength \& he shall speak I am he in righteousness mighty to save \& the Lord shall be read \sout{read} in his appearl \& his garments like him that treadeth in the wine path \& so great shall be the glory in his presence that the Sun shall hide his face in shame \& the moon shall be blown out \& the Stars shall be hurrelled from their sockets \& his voice shall be heard I have trodden the wine press alone \& have brought Judgement upon all people \& none was with me \& I have trampelled them in my fury \& I did tread upon them in mine anger \& their blood have I sprinkled upon my garments \& have stained all my raiment for this was the day of vengeance which was in my heart \& now the year of my redeemed is come \& they shall mention the loveing kindness of their Lord \& all that he hath bestowed upon them according to his goodness \& according to his loving kindness forever \& ever in all their afflictions he was afflicted \& the angel of his presenc[e] saved them \& in his love \& in his pity he redeemed them \& did bear them \& did carry them all the days of old yea \& Enoch also \& they which were with him the Prophets which were before him \& Noah also \& they which were before him \& Elijah also \& they which were before him \& from Elijah to Moses \& from Moses to John who were with Christ in his resurrection \& the Holy Apostles with Abraham Isaac \& Jacob shall be in the presenc[e] of the lamb \& the graves of the saints shall be opened \& they shall come forth \& stand on the right hand of the Lamb when he shall stand upon mount Zion \& upon the Holy City the New Jerusalem \sout{wherefo} \& they shall sing the Song of the lamb day \& night for ever \& ever--- \& for this cause that men might be partakers of the glories which were revealed the Lord sent forth the fullness of the gospel \& the everlasting covenant reas[o]ning in plainness \& simplisity to prepare the weak for those things which are coming upon the earth \& for the Lords Errand in the days when the <weak should confound the wise> \& \textit{[illegible]} little one become a strong nation \& two should put their tens of thousands to flight by the weak things of the Earth the Lord should thresh the Nations of the Earth by the power of his spirit \& for this cause these commandments were given they were commanded to be kept from the world in the day <that> they were given but now are to go forth unto all flesh \& this according to the mind \& the will of the Lord who reigneth \sout{of} over all flesh \& unto him that repenteth \& sanctifieth himself before the Lord shall be given eternal life \& they that harken not to the voice of the Lord shall be fulfilled that which was written by the Prophet Moses \sout{\& also} that they should be cut off from among the people \& also that which was written by the Prophet M$\diamond$lichi$\diamond$ for Behold the day cometh that burneth as an Oven \& all the proud \sout{\& they that do} yea \& all that do wickedly shall be stuble \& the day that cometh shall burn them up saith the Lord of hosts that it shall l[e]ave them neither root nor branch wherefore this shall be the answer of the Lord in that day when I come unto my own no man among you received me \& ye are driven out when I called again there was none of you to answer yet my arm was not shortened at all that I could not redeem neither my power to deliver Behold at my rebuke I dry up the Sea \sout{is} I make the rivers a wilderness their fish stinketh \& dieth for thirst I Cloth[e] the Heavens with blackness \& make sackloth their covering \& this shall have \sout{at} of my hand ye shall lay down with sorrow Behold \& Lo there is none to deliver you for ye obeyed not my voice when I called unto you \sout{all} out of the Heavens ye believed not my Servants \& when they were sent unto you ye received them not wherefore they sealed up the testimony \& bound up the Law \& ye were delivered up unto darkness these shall go away into outer darkness where there is weeping \& wailing \& gnashing of theeth Behold the Lord your God hath spoken it Amen
				
				Given In Hiram Portage Co Ohio
				
				Novembr 3\supersu{rd}\supers{.} 1831
				
		% -----
		\subsubsection{Minutes, 8 November 1831}
		% -----
			\label{EMS-1831-JS-8-NOV-1831}
			% \label{EMS-1830-JS-DAY-3LETTERMONTH-YEAR}
		
			\begin{description}
				\item[Print Sources]
			\end{description}

				\begin{tabular}{p{0.5in}p{1in}p{0.5in}p{1in}}
				JSP	 	& D2:121--124	& JSR 		& --- \\
				DC	 	& ---		 	& HC 		& 1:235 \\
				HDDC 	& ---			& TCHC 		& 1:163 \\
				\end{tabular}
				% HDDC = woodford's DC diss; JSR = Marquardt's JS Revelations; TCHC = Vogel HC
			
			\begin{description}
				\item[Online Access] \href{https://www.josephsmithpapers.org/paper-summary/minutes-8-november-1831/1}{Here}
				% \item[Analysis] \hyperref[DC133]{Here}
			\end{description}
			
			\begin{description}
				\item[Background]
			\end{description}
			
				% It might also be useful to give a two sentence context of the period: e.g., "This speech was given in Nauvoo to a small congregation one month prior to the destruction of the Nauvoo Expositor. These and other tensions may have been on the prophet's mind when he gave the speech."
			
			\begin{description}
				\item[Original Source]
			\end{description}
			
				% brief description of original source material, with reference to JSP or somewhere else for more info
				% Background material: it might be helpful to have a note that says whether a given document was presented as a speech, was written in a journal, etc.
				% Also a comment here about how reliable the source might be, or what shape it is in, if there are large omissions or parts that are hard to understand, and so on.
			
			\begin{description}
				\item[Text]
			\end{description}
			
					\hypertarget{EMS-1831-JS-8-NOV-1831-a}%
				...Br.
					\marginnote{a}%
				Joseph Smith jr. appointed Moderator \& John Whitmer Clerk. Opened, prayer by br. Joseph Smith jr.
				
				Remarks by br. Sidney Rigdon on the errors or mistakes which are in commandments and revelations, made either by the \sout{scribe} translation in consequence of the slow way of the scribe at the time of receiving or by the scribes themselves
				
					\hypertarget{EMS-1831-JS-8-NOV-1831-b}%
				Resolved
					\marginnote{b}%
				by this conference that Br Joseph Smith Jr correct those errors or mistakes which he may discover by the holy Spirit while \sout{receiving the} \sout{revelations} \sout{reviewing} reviewing the revelations \& commandments \& also the fulness of the scriptures...
				
		% -----
		\subsubsection{Revelation, 11 November 1831--A [D\&C 69]}
		% -----
			\label{EMS-1831-JS-11-NOV-1831-A}
			% \label{EMS-1830-JS-DAY-3LETTERMONTH-YEAR}
		
			\begin{description}
				\item[Print Sources]
			\end{description}

				\begin{tabular}{p{0.5in}p{1in}p{0.5in}p{1in}}
				JSP	 	& D1:129--132	& JSR 		& 179 \\
				DC	 	& Section 69 	& HC 		& 1:234--235 \\
				HDDC 	& 866--872		& TCHC 		& 1:163 \\
				\end{tabular}
				% HDDC = woodford's DC diss; JSR = Marquardt's JS Revelations; TCHC = Vogel HC
			
			\begin{description}
				\item[Online Access] \href{https://www.josephsmithpapers.org/paper-summary/revelation-11-november-1831-a-dc-69/1}{Here}
				% \item[Analysis] \hyperref[DC69]{Here}
			\end{description}
			
			\begin{description}
				\item[Background]
			\end{description}
			
				% It might also be useful to give a two sentence context of the period: e.g., "This speech was given in Nauvoo to a small congregation one month prior to the destruction of the Nauvoo Expositor. These and other tensions may have been on the prophet's mind when he gave the speech."
			
			\begin{description}
				\item[Original Source]
			\end{description}
			
				% brief description of original source material, with reference to JSP or somewhere else for more info
				% Background material: it might be helpful to have a note that says whether a given document was presented as a speech, was written in a journal, etc.
				% Also a comment here about how reliable the source might be, or what shape it is in, if there are large omissions or parts that are hard to understand, and so on.
			
			\begin{description}
				\item[Text]
			\end{description}
			
				\begin{tightcenter}
				<74> Received on \supersu{the}\supers{.} 11 of Oct 1831
				\end{tightcenter}
				
				Hearken unto me saith the Lord for verily I say unto you for my Servent Olivers [Oliver Cowdery's] sake it is not wisdom in me that he should be intrusted with the commandments \& moneys which he shall carry unto the Land of Zion except one go with him who will be true \& faithfull wherefore I the Lord willeth that my Servent John (Whitmer) shall go with my servent Oliver \& also that he observe to continue in writing \& makeing a history of all the important things which he shall observe \& know concerning my Church \& also that he receive council \& assistance from my Servent Oliver \& others \& also that my Saints which are abroad in the Earth should send forth their accounts to the Land of Zion for the Land of Zion shall be a seat \& a place to receive \& do all these things nevertheless let my Servnt John travel many times from place to place \& from Church to Church that he may the more easily obtain knowledge Preaching \& expounding writing cop[y]ing \& selecting \& obtain[in]g all things which shall be for the good of the Church \& for the rising generations which shall grow up on the Land of Zion to possess it from generation\sout{s} to generation\sout{s} forever \& ever Amen
				
		% -----
		\subsubsection{Revelation, 11 November 1831--B [D\&C 107 (partial)]}
		% -----
			\label{EMS-1831-JS-11-NOV-1831-B}
			% \label{EMS-1830-JS-DAY-3LETTERMONTH-YEAR}
		
			\begin{description}
				\item[Print Sources]
			\end{description}

				\begin{tabular}{p{0.5in}p{1in}p{0.5in}p{1in}}
				JSP	 	& D2:132--136	& JSR 		& 177--179 \\
				DC	 	& Section 107 	& HC 		& 2:210--217 \\
				HDDC 	& 1398--1432	& TCHC 		& 2:218--222 \\
				\end{tabular}
				% HDDC = woodford's DC diss; JSR = Marquardt's JS Revelations; TCHC = Vogel HC
			
			\begin{description}
				\item[Online Access] \href{https://www.josephsmithpapers.org/paper-summary/revelation-11-november-1831-b-dc-107-partial/1}{Here}
				% \item[Analysis] \hyperref[DC107]{Here}
			\end{description}
			
			\begin{description}
				\item[Background]
			\end{description}
			
				% It might also be useful to give a two sentence context of the period: e.g., "This speech was given in Nauvoo to a small congregation one month prior to the destruction of the Nauvoo Expositor. These and other tensions may have been on the prophet's mind when he gave the speech."
			
			\begin{description}
				\item[Original Source]
			\end{description}
			
				% brief description of original source material, with reference to JSP or somewhere else for more info
				% Background material: it might be helpful to have a note that says whether a given document was presented as a speech, was written in a journal, etc.
				% Also a comment here about how reliable the source might be, or what shape it is in, if there are large omissions or parts that are hard to understand, and so on.
			
			\begin{description}
				\item[Text]
			\end{description}
			
				\begin{tightcenter}
				<75> A Revalation given at Hiram Portage Co No\uline{v} 11\supersu{th}\supers{.} 1831
				\end{tightcenter}
				
				To the Church of Christ in the Land of Zion in addition to the Church Laws respecting Church business verily I say unto you saith the Lord of hosts there must needs be \sout{presiding} presiding Elders to preside over them who are of the office of an Elder \& also Priests over them who are of the office of a Priest \& also Teachers over them who are of the office of a Teacher in like manner And also the deacons wherefore from Deacon to Teacher \& from Teacher to Priest \& from Priest to Elder severally as they are appointed, according to the Church Articles \& Covenants then cometh the high Priest hood which is the greatest of all wherefore it must needs be that one be appointed of the high Priest hood to preside over the Priest hood \& he shall be called President of the \sout{hood} high Priest hood of the Church or in \sout{o} other \sout{high} words the Presiding high Priest \sout{hood} over the high Priesthood of the Church from the same cometh the administring of ordinances \& blessings upon the Church by the Laying on of the hands wherefore the office of a Bishop is not equal unto it for the office of a Bishop is in administering all \sout{things} temporal things nevertheless a Bishop must be chosen from the high Priesthood that he may be set apart unto the ministering of temporal things having a knowledge of them by the Spirit of truth \& also to be a Judge in Israel to do the business of the Church to sit \sout{down} in Judgement upon transgressors upon testimony it shall be laid before them according to the Laws by the assistance of his councillors whom he hath chosen or will choose among the Elders of the church thus shall he be a judge even a common judge among the inhabitants of Zion until the borders are enlarged \& it becomes necessary to have other Bishops or judges \& inasmuch as there are other Bishops appointed they shall act in the same office \& again verily I say unto you the most important business of the church \& the most difficult cases of the church inasmuch as there is not \sout{sufficient} satisfaction upon the decision of the judge it shall be handed over \& carried up unto the court of the church before the president of the high Priesthood \& the president of the Court of the high priesthood shall have power to call other high priests even twelve to assist as counsellors \& thus the president of the high priesthood \& his councellors shall have power to decide upon testimony according to the laws of the church \& after this desision it shall be had in remembrance no more before the Lord for this is the highest court of the church of God \& a final desision upon controvers[i]es all persons belonging to the church are not exempt from this court of the church \& inasmuch as the president of the high priesthood shall transgress he shall be had in remembrance before the common court of the church who shall be assisted by twelve counsellors of the high Priesthood \& their desicision upon his head shall be an end of controversy concerning him thus none shall be exempt from the justice of the Laws of God that all things may be done in order \& in solemnity before me \sout{in} according to truth \& righteousness Amen A few more words in addition to the Laws of the church And again verily I say unto you the duty of a president over the office of a Deacon is to preside over twelve Deacons to set in council with them \& to teach them their duty edifying one another as it is given according to the covenants And also the duty of the president over the office of the Teachers is to preside over twenty four of the Teachers \& to set in council with them teaching them the duties of their office as given in the covenants Also the duty of the president over the priesthood is to preside over forty eight priests \& to set in council with them \& to teach them the duties of their office as given in the covenants And again the duty of the president over the office of the Elders is to preside over ninety six Elders \& to set in council with them \& to <teach> them according to the covenants And again the duty of the president of the office of the High Priesthood is to preside over the whole church \& to be like unto Moses behold here is wisdom yea to be a Seer a revelator a translator \& a prophet having all the gifts of God which he bestoweth upon the head of the chuch Wherefore now let every man learn his \sout{duly} duty \& to act in the office in which he is appointed in all diligence he that is slothful shall not be counted worthy to stand \& he that learneth not his duty \& sheweth himself not approved shall not be counted worthy to stand even so Amen
				
		% -----
		\subsubsection{Minutes, 12 November 1831}
		% -----
			\label{EMS-1831-JS-12-NOV-1831-MIN}
			% \label{EMS-1830-JS-DAY-3LETTERMONTH-YEAR}
		
			\begin{description}
				\item[Print Sources]
			\end{description}

				\begin{tabular}{p{0.5in}p{1in}p{0.5in}p{1in}}
				JSP	 	& D2:136--139	& JSR 		& --- \\
				DC	 	& ---		 	& HC 		& 1:235--236 \\
				HDDC 	& ---			& TCHC 		& 1:163--164 \\
				TPJS	& 7--8			& 			&  \\
				\end{tabular}
				% HDDC = woodford's DC diss; JSR = Marquardt's JS Revelations; TCHC = Vogel HC
			
			\begin{description}
				\item[Online Access] \href{https://www.josephsmithpapers.org/paper-summary/minutes-12-november-1831/1}{Here}
				% \item[Analysis] \hyperref[XX]{Here}
			\end{description}
			
			\begin{description}
				\item[Background]
			\end{description}
			
				% It might also be useful to give a two sentence context of the period: e.g., "This speech was given in Nauvoo to a small congregation one month prior to the destruction of the Nauvoo Expositor. These and other tensions may have been on the prophet's mind when he gave the speech."
			
			\begin{description}
				\item[Original Source]
			\end{description}
			
				% brief description of original source material, with reference to JSP or somewhere else for more info
				% Background material: it might be helpful to have a note that says whether a given document was presented as a speech, was written in a journal, etc.
				% Also a comment here about how reliable the source might be, or what shape it is in, if there are large omissions or parts that are hard to understand, and so on.
			
			\begin{description}
				\item[Text]
			\end{description}
			
					\hypertarget{EMS-1831-JS-12-NOV-1831-MIN-a}%
				...Br.
					\marginnote{a}%
				Joseph Smith jr. said one item he wished acted upon was that our brs. Oliver Cowdery \& John Whitmer \& the sacred writings which they have entrusted to them to carry to Zion be dedicated to the Lord by the prayer of faith. Secondly, Br. Oliver has labored with me from the begining in writing \&c. Br. Martin [Harris] has labored with me from the begining, \& brs. John \& Sidney also for a considerable time, \& as these sacred writings are now going to the Church for their benefit, that we may have claim on the Church for recompence, if this conference think these things worth prizing to be had on record to show hereafter I feel that it will be according to the mind of the Spirit for by it these things were put into my heart which I know to be the Spirit of truth \&c.
				
				Voted that Joseph Smith jr. be appointed to dedicate \& consecrate these brethren \& the sacred writings \& all they have entrusted to their care, to the Lord: done accordingly.
				
					\hypertarget{EMS-1831-JS-12-NOV-1831-MIN-b}%
				After
					\marginnote{b}%
				deliberate consideration in consequence of the book of Revelation now to be printed being the foundation of the Church \& the salvation of the world \& the Keyes of the mysteries of the Kingdom, \& the riches of Eternity to the church. Voted that they be prized by this Conference to be worth to the Church the riches of the whole Earth. Speaking temporally.
				
					\hypertarget{EMS-1831-JS-12-NOV-1831-MIN-c}%
				Voted
					\marginnote{c}%
				that in consequence of the dilligence of our brethren, Joseph Smith jr. Oliver Cowdery John Whitmer \& Sidney Rigdon in bringing to light by the grace of God these sacred things, be appointed to manage them according to the Laws of the Church \& the commandments of the Lord...
				
		% -----
		\subsubsection{Revelation, 12 November 1831 [D\&C 70]}
		% -----
			\label{EMS-1831-JS-12-NOV-1831}
			% \label{EMS-1830-JS-DAY-3LETTERMONTH-YEAR}
		
			\begin{description}
				\item[Print Sources]
			\end{description}

				\begin{tabular}{p{0.5in}p{1in}p{0.5in}p{1in}}
				JSP	 	& D2:139--141	& JSR 		& 180 \\
				DC	 	& Section 70 	& HC 		& 1:236--237 \\
				HDDC 	& 873--929		& TCHC 		& 1:164 \\
				\end{tabular}
				% HDDC = woodford's DC diss; JSR = Marquardt's JS Revelations; TCHC = Vogel HC
			
			\begin{description}
				\item[Online Access] \href{https://www.josephsmithpapers.org/paper-summary/revelation-12-november-1831-dc-70/1}{Here}
				% \item[Analysis] \hyperref[DC70]{Here}
			\end{description}
			
			\begin{description}
				\item[Background]
			\end{description}
			
				% It might also be useful to give a two sentence context of the period: e.g., "This speech was given in Nauvoo to a small congregation one month prior to the destruction of the Nauvoo Expositor. These and other tensions may have been on the prophet's mind when he gave the speech."
			
			\begin{description}
				\item[Original Source]
			\end{description}
			
				% brief description of original source material, with reference to JSP or somewhere else for more info
				% Background material: it might be helpful to have a note that says whether a given document was presented as a speech, was written in a journal, etc.
				% Also a comment here about how reliable the source might be, or what shape it is in, if there are large omissions or parts that are hard to understand, and so on.
			
			\begin{description}
				\item[Text]
			\end{description}
			
				\begin{tightcenter}
				<Revelation>
				\end{tightcenter}
				
				<76> <Hiram> Nov. 12. 1831
				
				Behold \& hearken o ye inhabitants of Zion \& all ye people of my Church which are far off \& hear the word of the Lord which I give unto my servant Jos[e]ph \& also unto my servant Martin [Harris] \& also unto my servant Oliver [Cowdery] \& also my servant John [Whitmer] \& also unto my servant Sidney [Rigdon] by the way of commandment\sout{s} unto them for I give unto them a commandment W[h]erefore hearken \& hear for thus saith the Lord unto them I the Lord have appointed them \& ordained them to be stewards over the revelations \& commandments which I have given unto them \& which I shall hereafter give unto them \& an account of this stewardship will I require of them in the day of judgement wherefore I have appointed unto them \& this is their business in the church of God to manage them \& the concerns thereof yea the profits thereof wherefore a commandment I give unto them that they shall not give these things unto the church neither unto the world nevertheless inasmuch as they receive more than is for their necessities \& their wants it shall be given into my storehouse \& the benefits thereof shall be consecrated unto the inhabtants of Zion \& unto their generations inasmuch as they become heirs according to the laws of the kingdom behold this is what the Lord requires of every man in his stewardship even as I the Lord have appointed or shall hereafter apoint unto any man \& behold none is exempt from this law who belong to the church of the Living God yea neither Bishop neither the agent who keepeth the Lords storehouse neither he that is appointed in a stewardship over temporal things he that is appointed to administer spiritual things the same is worthy of his hire even as they who are appointed in a stewardship to administer in temporal things yea even more abundantly which abundance is multiplied unto them through the manifestations of the spirit nevertheless in your temporal things you shall be equal in all things \& this not grudgeingly otherwise the abundance of the manifestations of the spirit shall be withheld now this commandment I give unto my servants while they remain for a manifestations of my blessings upon their heads \& for a reward of their diligence \& for their security for food \& for raiment for an inheritance for houses \& for lands \& in whatsoever circumstances I the Lord shall place them \& whithersoever I the Lord shall send them for they have been faithful over many things \& have done well inas much as they have not sin[n]ed behold I the Lord am merciful \& will bless them \& they shall enter into the joy of these things
				
				even so Amen and again verily I say unto you that my servant William [W. Phelps] shall be included in this commandment with you in this same stewardship even so Amen
				
		% -----
		\subsubsection{Revelation, 1 December 1831 [D\&C 71]}
		% -----
			\label{EMS-1831-JS-1-DEC-1831}
			% \label{EMS-1830-JS-DAY-3LETTERMONTH-YEAR}
		
			\begin{description}
				\item[Print Sources]
			\end{description}

				\begin{tabular}{p{0.5in}p{1in}p{0.5in}p{1in}}
				JSP	 	& D2:144--146	& JSR 		& 181 \\
				DC	 	& Section 71 	& HC 		& 1:238--239 \\
				HDDC 	& 881--886		& TCHC 		& 1:165 \\
				\end{tabular}
				% HDDC = woodford's DC diss; JSR = Marquardt's JS Revelations; TCHC = Vogel HC
			
			\begin{description}
				\item[Online Access] \href{https://www.josephsmithpapers.org/paper-summary/revelation-1-december-1831-dc-71/1}{Here}
				% \item[Analysis] \hyperref[DC71]{Here}
			\end{description}
			
			\begin{description}
				\item[Background]
			\end{description}
			
				% It might also be useful to give a two sentence context of the period: e.g., "This speech was given in Nauvoo to a small congregation one month prior to the destruction of the Nauvoo Expositor. These and other tensions may have been on the prophet's mind when he gave the speech."
			
			\begin{description}
				\item[Original Source]
			\end{description}
			
				% brief description of original source material, with reference to JSP or somewhere else for more info
				% Background material: it might be helpful to have a note that says whether a given document was presented as a speech, was written in a journal, etc.
				% Also a comment here about how reliable the source might be, or what shape it is in, if there are large omissions or parts that are hard to understand, and so on.
			
			\begin{description}
				\item[Text]
			\end{description}
			
				\begin{tightcenter}
				Hiram Portage county Ohio D\supersu{ec}\supers{.} 1--- 1831
				\end{tightcenter}
				
				Behold \sout{this} thus saith the Lord unto you my Servents that the time has verily come that it is necessary and expedient in me that you should open your mouths in proclaiming my gospel the things of the kingdom expounding the misteries thereof out of the Schriptures according to that portion of spirit and power which shall be given unto you even as I will verily I say unto you proclaim unto the world in the regions round about and in the church also for the space of a season even untill it shall be made known unto you verily this is a mission for a season which I give unto you wherefore labour ye in my vinyard call upon the inhabitants of the earth and bear record and prepare the way for the \sout{revelations} commandments and <the> revelations which are to come Now behold this is wisdom whoso readeth let him understand and receive also for unto him who receiveth it shall be given more abundantly even power wherefore confound your enemies Call upon them to meet you \sout{<both in>} \sout{at publick} both in publick and in private and inasmuch as ye are faithfull their shame shall be made manifest wherefore let them bring forth their strong reasons against the Lord verily thus saith the Lord unto you there is no weapon that <is> formed against you shall prosper and if any man lift his voice against you he shall be confounded in mine own due time wherefore keep these commandments for they are true and faithfull even so amen [\textit{rest of page blank}]
				
		% -----
		\subsubsection{Revelation, 4 December 1831--A [D\&C 72:1--8]}
		% -----
			\label{EMS-1831-JS-4-DEC-1831-A}
			% \label{EMS-1830-JS-DAY-3LETTERMONTH-YEAR}
		
			\begin{description}
				\item[Print Sources]
			\end{description}

				\begin{tabular}{p{0.5in}p{1in}p{0.5in}p{1in}}
				JSP	 	& D2:146--150	& JSR 		& 181--182 \\
				DC	 	& Section 72 	& HC 		& 1:239--241 \\
				HDDC 	& 887--900		& TCHC 		& 1:166--167 \\
				\end{tabular}
				% HDDC = woodford's DC diss; JSR = Marquardt's JS Revelations; TCHC = Vogel HC
			
			\begin{description}
				\item[Online Access] \href{https://www.josephsmithpapers.org/paper-summary/revelation-4-december-1831-a-dc-721-8/1}{Here}
				% \item[Analysis] \hyperref[DC72]{Here}
			\end{description}
			
			\begin{description}
				\item[Background]
			\end{description}
			
				% It might also be useful to give a two sentence context of the period: e.g., "This speech was given in Nauvoo to a small congregation one month prior to the destruction of the Nauvoo Expositor. These and other tensions may have been on the prophet's mind when he gave the speech."
			
			\begin{description}
				\item[Original Source]
			\end{description}
			
				% brief description of original source material, with reference to JSP or somewhere else for more info
				% Background material: it might be helpful to have a note that says whether a given document was presented as a speech, was written in a journal, etc.
				% Also a comment here about how reliable the source might be, or what shape it is in, if there are large omissions or parts that are hard to understand, and so on.
			
			\begin{description}
				\item[Text]
			\end{description}
				
				\begin{tightcenter}
				Kirtland December 4\supers{th}--- 1831
				\end{tightcenter}
				
				Hearken and listen to the voice of the Lord o ye who have assembled yourselves together who are the high priests of my church to whom the kingdom and power have been given for verily thus saith the Lord it is expedient in me for a Bishop to be appointed unto you or of you \sout{to} unto the church in this part of the Lords vineyard and verily in this thing ye have done wisely for it is required of the Lord at the hand of every steward to render an account of his stewardship both in time and in eternity for he who is faithfull and wise in time is accounted worthy to inherit the mantions prepared for them of my father verily I say unto you the Elders of the church in this part of my vineyard shall render an account of their stewardship unto the Bishop which shall be appointed of me in this part of my vinyard these things shall be had on record to be handed over unto the Bishop in Zion and the duty of the Bishop shall be made known by the commandments which have been given and by the voice of the conference.
				
				And now I say unto you my servent Newel [K. Whitney] is the man who shall be appointed and ordained unto this power this is the will of the Lord your God your Redeemer even so Amen---
				
		% -----
		\subsubsection{Revelation, 4 December 1831--B [D\&C 72:9--23]}
		% -----
			\label{EMS-1831-JS-4-DEC-1831-B}
			% \label{EMS-1830-JS-DAY-3LETTERMONTH-YEAR}
		
			\begin{description}
				\item[Print Sources]
			\end{description}

				\begin{tabular}{p{0.5in}p{1in}p{0.5in}p{1in}}
				JSP	 	& D2:151--153	& JSR 		& 182--183 \\
				DC	 	& Section 72 	& HC 		& 1:239--241 \\
				HDDC 	& 887--900		& TCHC 		& 1:166--167 \\
				\end{tabular}
				% HDDC = woodford's DC diss; JSR = Marquardt's JS Revelations; TCHC = Vogel HC
			
			\begin{description}
				\item[Online Access] \href{https://www.josephsmithpapers.org/paper-summary/revelation-4-december-1831-b-dc-729-23/1}{Here}
				% \item[Analysis] \hyperref[DC72]{Here}
			\end{description}
			
			\begin{description}
				\item[Background]
			\end{description}
			
				% It might also be useful to give a two sentence context of the period: e.g., "This speech was given in Nauvoo to a small congregation one month prior to the destruction of the Nauvoo Expositor. These and other tensions may have been on the prophet's mind when he gave the speech."
			
			\begin{description}
				\item[Original Source]
			\end{description}
			
				% brief description of original source material, with reference to JSP or somewhere else for more info
				% Background material: it might be helpful to have a note that says whether a given document was presented as a speech, was written in a journal, etc.
				% Also a comment here about how reliable the source might be, or what shape it is in, if there are large omissions or parts that are hard to understand, and so on.
			
			\begin{description}
				\item[Text]
			\end{description}
			
				\begin{tightcenter}
				The duty of the Bishop as made known at the same time
				\end{tightcenter}
				
				The word of the Lord in addition to the law which has been given making known the duty of the Bishop which has been ordained unto the church in this part of the vinyard which is verily this. To keep the Lords storehouse to receive the funds of the church in this part of the vinyard to take an account of the Elders as before has been commanded and to administer to their wants who shall pay for that which they receive inasmuch as they have wherewith to pay that this also may be consecrated to the good of the church to the poor and needy and he who hath not wherewith to pay an account shall be taken and handed over to the Bishop in Zion who shall pay the debt out of that which the Lord shall put into his hands and the labours of the faithfull who labour in spiritual things in administering the gospel and the things of the kingdom unto the church and unto the world shall answer the debt unto the Bishop in Zion thus it cometh out of the church for according to the law every man who cometh up to Zion must lay all things before the Bishop in Zion. And now verily I say unto you that as every Elder in this part of the vinyard must give an account of his stewardship unto the Bishop in this part of the vinyard a certificate from the judge or Bishop in this part of the vinyard unto the Bishop in Zion rendereth every man acceptable and answereth all things for an inheritence and to be received as a wise steward and as a faithfull labourer otherwise shall not be accepted of the Bishop in Zion. And now verily I say unto you let every Elder who shall give an account unto the Bishop of the church in this part of the vinyard be recommended by the church or churches in which he labours that he may render himself and his accounts approved in all things
				
				And again let my servents who are appointed as stewards over the litterary concerns of my church have claim for assistence \sout{in all things} upon the Bishop or Bishops in all things that the revelations may be published and go forth unto the ends of the earth that they also may obtain funds which shall benefit the church in all things that they also may render themselves approved in all things and be accounted as wise stewards. And behold this shall be an ensample for all the extensive branches of my church in whatsoever land they shall be established and now I make an end of my sayings Amen
				
		% -----
		\subsubsection{Revelation, 4 December 1831--C [D\&C 72:24--26]}
		% -----
			\label{EMS-1831-JS-4-DEC-1831-C}
			% \label{EMS-1830-JS-DAY-3LETTERMONTH-YEAR}
		
			\begin{description}
				\item[Print Sources]
			\end{description}

				\begin{tabular}{p{0.5in}p{1in}p{0.5in}p{1in}}
				JSP	 	& D2:153--154	& JSR 		& 182--183 \\
				DC	 	& Section 72 	& HC 		& 1:239--241 \\
				HDDC 	& 887--900		& TCHC 		& 1:166--167 \\
				\end{tabular}
				% HDDC = woodford's DC diss; JSR = Marquardt's JS Revelations; TCHC = Vogel HC
			
			\begin{description}
				\item[Online Access] \href{https://www.josephsmithpapers.org/paper-summary/revelation-4-december-1831-c-dc-7224-26/1}{Here}
				% \item[Analysis] \hyperref[DC72]{Here}
			\end{description}
			
			\begin{description}
				\item[Background]
			\end{description}
			
				% It might also be useful to give a two sentence context of the period: e.g., "This speech was given in Nauvoo to a small congregation one month prior to the destruction of the Nauvoo Expositor. These and other tensions may have been on the prophet's mind when he gave the speech."
			
			\begin{description}
				\item[Original Source]
			\end{description}
			
				% brief description of original source material, with reference to JSP or somewhere else for more info
				% Background material: it might be helpful to have a note that says whether a given document was presented as a speech, was written in a journal, etc.
				% Also a comment here about how reliable the source might be, or what shape it is in, if there are large omissions or parts that are hard to understand, and so on.
			
			\begin{description}
				\item[Text]
			\end{description}
			
				A few words in addition to the laws of the kingdom respecting the members of the church they who are appointed by the holy spirit to go up unto Zion <\&> they \sout{who} who are priveledged to go up unto Zion let them carry up unto the bishop a certificate from three Elders of the church or a certificate from the Bishop otherwise he who shall go up unto the land of Zion shall not be accounted \sout{worthy} a wise steward this als[o] is an ensample Amen [\textit{rest of page blank}]